\newMonth{November}
\setrunningtitles{November}{November}

% ---- martyrology/mart11/mart1101.htm
\needspace{10\baselineskip}
\begin{paracol}{2}
\selectlanguage{latin}
\begin{center}{\color{gregoriocolor} Kaléndis Novémbris. 
 Luna\dots\ }\end{center}
\switchcolumn
\selectlanguage{english}
\begin{center}{\color{gregoriocolor} The 1\textsuperscript{st} Day of 
 November. The\dots\ Day of the Moon.}\end{center}
\end{paracol}

\noindent\begin{tabularx}{\linewidth}{*{19}{>{\centering\arraybackslash}X}}
 \textcolor{gregoriocolor}{a} & \textcolor{gregoriocolor}{b} & \textcolor{gregoriocolor}{c} & \textcolor{gregoriocolor}{d} & \textcolor{gregoriocolor}{e} & \textcolor{gregoriocolor}{f} & \textcolor{gregoriocolor}{g} & \textcolor{gregoriocolor}{h} & \textcolor{gregoriocolor}{i} & \textcolor{gregoriocolor}{k} & \textcolor{gregoriocolor}{l} & \textcolor{gregoriocolor}{m} & \textcolor{gregoriocolor}{n} & \textcolor{gregoriocolor}{p} & \textcolor{gregoriocolor}{q} & \textcolor{gregoriocolor}{r} & \textcolor{gregoriocolor}{s} & \textcolor{gregoriocolor}{t} & \textcolor{gregoriocolor}{u} \\
 11 & 12 & 13 & 14 & 15 & 16 & 17 & 18 & 19 & 20 & 21 & 22 & 23 & 24 & 25 & 26 & 27 & 28 & 29 \\
\end{tabularx}
\vspace{0.5\baselineskip}
\noindent\begin{tabularx}{\linewidth}{*{12}{>{\centering\arraybackslash}X}}
 \textcolor{gregoriocolor}{A} & \textcolor{gregoriocolor}{B} & \textcolor{gregoriocolor}{C} & \textcolor{gregoriocolor}{D} & \textcolor{gregoriocolor}{E} & F & \textcolor{gregoriocolor}{F} & \textcolor{gregoriocolor}{G} & \textcolor{gregoriocolor}{H} & \textcolor{gregoriocolor}{M} & \textcolor{gregoriocolor}{N} & \textcolor{gregoriocolor}{P} \\
 30 & 1 & 2 & 3 & 4 & 5 & 5 & 6 & 7 & 8 & 9 & 10 \\
\end{tabularx}

\begin{paracol}{2}
\selectlanguage{latin}
\lettrine[lines=2]{F}{estívitas} ómnium 
 Sanctórum, quam in honórem beátæ Dei Genitrícis Vírginis Maríæ et sanctórum 
 Mártyrum Bonifátius Papa Quartus, cum templum Pántheon tértio Idus Maji 
 dedicásset, célebrem et generálem instítuit agi quotánnis in urbe Roma. 
 Sed Gregórius item Quartus póstmodum decrévit, eándem festivitátem, quæ 
 váriis modis jam in divérsis Ecclésiis celebrabátur, in honórem ómnium 
 Sanctórum solémniter hac die ab univérsa Ecclésia perpétuo observári.
\switchcolumn
\selectlanguage{english}
\lettrine[lines=2]{T}{he} Festival of All Saints, which Pope Boniface IV, after the dedication 
 of the Pantheon, ordained to be kept generally and solemnly every year on 
 the 13th of May, in the city of Rome, in honour of the blessed Virgin Mary, 
 Mother of God, and of the holy martyrs. It was afterwards decreed by 
 Gregory IV that this feast, which was then celebrated in many dioceses, but 
 at different times, should be on this day kept by the whole Church in honour 
 of all the saints.
\switchcolumn*
\selectlanguage{latin}
In Pérside sanctórum 
 Mártyrum Joánnis Epíscopi, et Jacóbi Presbyteri, sub Sápore Rege.
\switchcolumn
\selectlanguage{english}
In Persia, the holy martyrs John, a 
 bishop, and James, a priest, under King Sapor.
\switchcolumn*
\selectlanguage{latin}
Tarracínæ, in Campánia, natális sancti Cæsárii Diáconi, qui, diébus multis in custódia macerátus, 
 póstea, cum sancto Juliáno Presbytero, in saccum missus et in mare 
 præcipitátus est.
\switchcolumn
\selectlanguage{english}
At Terracina in Campania, the 
 birthday of St. Caesarius, deacon, who was detained many days in prison, 
 afterwards put into a sack with the priest St. Julian, and then thrown into 
 the sea.
\switchcolumn*
\selectlanguage{latin}
In castro Divióne 
 sancti Benígni Presbyteri, qui a beáto Polycárpo missus est in Gálliam ad 
 prædicándum Evangélium; et, postquam, sub Marco Aurélio Imperatóre, a 
 Teréntio Júdice gravíssimis torméntis multiplíciter est afflíctus, tandem 
 ejus collum vecte férreo tundi et corpus láncea perforári jubétur.
\switchcolumn
\selectlanguage{english}
At Dijon, St. Benignus, a priest, 
 who was sent to France by blessed Polycarp to preach the Gospel. After 
 he had been subjected to many grievous torments by the judge Terentius, 
 under Emperor Marcus Aurelius, he was finally condemned to have his neck 
 struck with an iron bar and his body pierced with a lance.
\switchcolumn*
\selectlanguage{latin}
Damásci pássio 
 sanctórum Cæsárii, Dácii et aliórum quinque.
\switchcolumn
\selectlanguage{english}
At Damascus, the martyrdom of the 
 Saints Caesarius, Dacius, and five others.
\switchcolumn*
\selectlanguage{latin}
Eódem die sanctæ Maríæ 
 ancíllæ, quæ, Christiánæ religiónis nómine accusáta, hinc, sub Hadriáno 
 Imperatóre, diris verbéribus afflícta, equúlei extensiónem et ungulárum 
 laceratiónem passa, martyrium complévit.
\switchcolumn
\selectlanguage{english}
On the same day, St. Mary, a servant girl. Being 
 accused of professing the Christian religion in the time of Emperor Hadrian, 
 she was subjected to cruel scourging, to torture on the rack, and the 
 lacerating of her body with iron hooks, and thus completed her martyrdom
\switchcolumn*
\selectlanguage{latin}
Tarsi, in Cilícia, sanctárum Cyréniæ et Juliánæ Mártyrum, sub Maximiáno Imperatóre.
\switchcolumn
\selectlanguage{english}
At Tarsus in Cilicia, under Emperor 
 Maximian, the Saints Cyrenia and Juliana.
\switchcolumn*
\selectlanguage{latin}
Arvérnis, in Gállia, 
 sancti Austremónii, qui fuit primus ejúsdem civitátis Epíscopus.
\switchcolumn
\selectlanguage{english}
At Auvergne in France, St. 
 Austremonius, first bishop of Clermont.
\switchcolumn*
\selectlanguage{latin}
Lutétiæ Parisiórum 
 deposítio sancti Marcélli Epíscopi.
\switchcolumn
\selectlanguage{english}
At Paris, the death of St. 
 Marcellus, bishop.
\switchcolumn*
\selectlanguage{latin}
Bajócis, in Gállia, 
 sancti Vigóris Epíscopi, témpore Childebérti, Francórum Regis.
\switchcolumn
\selectlanguage{english}
At Bayeux, in the reign of the 
 Frankish king Childebert, St. Vigor, bishop.
\switchcolumn*
\selectlanguage{latin}
Andégavi, in Gállia, 
 deposítio sancti Licínii Epíscopi, venerábilis sanctitátis viri.
\switchcolumn
\selectlanguage{english}
At Angers in France, the death of 
 the aged holy man, St. Licinius, bishop.
\switchcolumn*
\selectlanguage{latin}
Tíbure sancti Severíni 
 Mónachi.
\switchcolumn
\selectlanguage{english}
At Tivoli, St. Severinus, monk.
\switchcolumn*
\selectlanguage{latin}
Lyricánti, in 
 Wastinénsi Gálliæ território, sancti Maturíni Confessóris.
\switchcolumn
\selectlanguage{english}
In Gatinais in France, St. Mathurin, 
 confessor.
\switchcolumn*
\selectlanguage{latin}
\end{paracol}

% \continueMonth{November}

% ---- martyrology/mart11/mart1102.htm
\needspace{10\baselineskip}
\begin{paracol}{2}
\selectlanguage{latin}
\begin{center}{\color{gregoriocolor} Quarto Nonas Novémbris. 
 Luna\dots\ }\end{center}
\switchcolumn
\selectlanguage{english}
\begin{center}{\color{gregoriocolor} The 2\textsuperscript{nd} Day of 
 November. The\dots\ Day of the Moon.}\end{center}
\end{paracol}

\noindent\begin{tabularx}{\linewidth}{*{19}{>{\centering\arraybackslash}X}}
 \textcolor{gregoriocolor}{a} & \textcolor{gregoriocolor}{b} & \textcolor{gregoriocolor}{c} & \textcolor{gregoriocolor}{d} & \textcolor{gregoriocolor}{e} & \textcolor{gregoriocolor}{f} & \textcolor{gregoriocolor}{g} & \textcolor{gregoriocolor}{h} & \textcolor{gregoriocolor}{i} & \textcolor{gregoriocolor}{k} & \textcolor{gregoriocolor}{l} & \textcolor{gregoriocolor}{m} & \textcolor{gregoriocolor}{n} & \textcolor{gregoriocolor}{p} & \textcolor{gregoriocolor}{q} & \textcolor{gregoriocolor}{r} & \textcolor{gregoriocolor}{s} & \textcolor{gregoriocolor}{t} & \textcolor{gregoriocolor}{u} \\
 12 & 13 & 14 & 15 & 16 & 17 & 18 & 19 & 20 & 21 & 22 & 23 & 24 & 25 & 26 & 27 & 28 & 29 & 30 \\
\end{tabularx}
\vspace{0.5\baselineskip}
\noindent\begin{tabularx}{\linewidth}{*{12}{>{\centering\arraybackslash}X}}
 \textcolor{gregoriocolor}{A} & \textcolor{gregoriocolor}{B} & \textcolor{gregoriocolor}{C} & \textcolor{gregoriocolor}{D} & \textcolor{gregoriocolor}{E} & F & \textcolor{gregoriocolor}{F} & \textcolor{gregoriocolor}{G} & \textcolor{gregoriocolor}{H} & \textcolor{gregoriocolor}{M} & \textcolor{gregoriocolor}{N} & \textcolor{gregoriocolor}{P} \\
 1 & 2 & 3 & 4 & 5 & 6 & 6 & 7 & 8 & 9 & 10 & 11 \\
\end{tabularx}

\begin{paracol}{2}

% TODO: add the rubrics for November 2nd
\selectlanguage{latin}
	\rubric{The latin rubric...}
\switchcolumn
\selectlanguage{english}
	\rubric{If the second day of November should occur on a Sunday so that the Feast of All
Souls is transferred to the Monday, the reading of the Martyrology today should
	commence with the announcement of the transferred Feast as follows:}
\switchcolumn*

\selectlanguage{latin}
\lettrine[lines=1]{C}{ommemorátio} ómnium Fidélium Defunctórum.
\switchcolumn
\selectlanguage{english}
\lettrine[lines=1]{T}{he} Commemoration of all the Faithful Departed.
\switchcolumn*

\selectlanguage{latin}
	\rubric{The latin rubric...}
\switchcolumn
\selectlanguage{english}
	\rubric{Otherwise the reading of the Martyrology commences as follows:}
\switchcolumn*

\selectlanguage{latin}
Si dies secunda Novembris incidat in Dominicam, prædicta verba leguntur 
 primo loco, die sequenti.
\switchcolumn
\selectlanguage{english}
If the second 
 day of November should occur on a Sunday, the above words are omitted and read 
 instead at the beginning on the following day.
\switchcolumn*
\selectlanguage{latin}
Pœtovióne, in Pannónia 
 superióre, natális sancti Victoríni, ejúsdem civitátis Epíscopi, qui, post 
 multa édita scripta (ut sanctus Hierónymus testátur), in persecutióne 
 Diocletiáni, martyrio coronátus est.
\switchcolumn
\selectlanguage{english}
At Pettau in Styria, the birthday of 
 St. Victorinus, bishop of that city, who, after publishing many writings, as 
 is attested to by St. Jerome, was crowned with martyrdom in the persecution 
 of Diocletian.
\switchcolumn*
\selectlanguage{latin}
Tergéste pássio beáti 
 Justi, qui in eádem persecutióne, sub Manátio Præside, martyrium consummávit.
\switchcolumn
\selectlanguage{english}
At Trieste, blessed Justus, who 
 fulfilled his martyrdom in the same persecution under the governor Manatius.
\switchcolumn*
\selectlanguage{latin}
Sebáste, in Arménia, sanctórum Cartérii, Styríaci, Tobíæ, Eudóxii, Agápii et Sociórum Mártyrum, 
 sub Licínio Imperatóre.
\switchcolumn
\selectlanguage{english}
At Sebaste in Armenia, the Saints 
 Carterius, Styriacus, Tobias, Eudoxius, Agapius, and their companions, 
 martyrs under Emperor Licinius.
\switchcolumn*
\selectlanguage{latin}
In Pérside sanctórum 
 Mártyrum Acíndyni, Pegásii, Aphthónii, Elpidíphori et Anempodísti, cum 
 plúrimis Sóciis.
\switchcolumn
\selectlanguage{english}
In Persia, the holy martyrs 
 Acindynus, Pegasius, Aphthonius, Elpiderphorus, and Anempodistus, with many 
 companions.
\switchcolumn*
\selectlanguage{latin}
In Africa natális 
 sanctórum Mártyrum Públii, Victóris, Hermétis, et Pápiæ.
\switchcolumn
\selectlanguage{english}
In Africa, the birthday of the holy 
 martyrs Publius, Victor, Hermes, and Papias.
\switchcolumn*
\selectlanguage{latin}
Tarsi, in Cilícia, 
 sanctæ Eustóchii, Vírginis et Mártyris; quæ, sub Juliáno Apóstata, post dira 
 torménta, in oratióne réddidit spíritum.
\switchcolumn
\selectlanguage{english}
At Tarsus in Cilicia, in the reign 
 of Julian the Apostate, St. Eustochium, virgin and martyr, who breathed her 
 last in prayer in the midst of severe torments.
\switchcolumn*
\selectlanguage{latin}
Laodicéæ, in Syria, 
 sancti Theódoti Epíscopi, qui non solum verbis, sed rebus quoque et 
 virtútibus fuit ornátus.
\switchcolumn
\selectlanguage{english}
At Laodicea in Syria, St. Theodotus, 
 a bishop powerful in words and adorned with good works and virtues.
\switchcolumn*
\selectlanguage{latin}
Viénnæ, in Gállia, 
 sancti Geórgii Epíscopi.
\switchcolumn
\selectlanguage{english}
At Vienne in France, the bishop St. 
 George.
\switchcolumn*
\selectlanguage{latin}
In monastério Agaunénsi, 
 in Gállia, sancti Ambrósii Abbátis.
\switchcolumn
\selectlanguage{english}
In the monastery of St. Moritz in 
 Switzerland, St. Ambrose, abbot.
\switchcolumn*
\selectlanguage{latin}
Cyri, in Syria, sancti 
 Marciáni Confessóris.
\switchcolumn
\selectlanguage{english}
At Cyrus in Syria, St. Marcian, 
 confessor.
\switchcolumn*
\selectlanguage{latin}
\end{paracol}

% TODO: check the martyrology to see how this is to be formatted...

% ---- martyrology/mart11/mart1103.htm
\needspace{10\baselineskip}
\begin{paracol}{2}
\selectlanguage{latin}
\begin{center}{\color{gregoriocolor}Tértio Nonas Novémbris, Luna\dots\ }\end{center}
\switchcolumn
\selectlanguage{english}
\begin{center}{\color{gregoriocolor}The 3\textsuperscript{rd} Day of 
 November, the... Day of the Moon.}\end{center}
\end{paracol}

\noindent\begin{tabularx}{\linewidth}{*{19}{>{\centering\arraybackslash}X}}
 \textcolor{gregoriocolor}{a} & \textcolor{gregoriocolor}{b} & \textcolor{gregoriocolor}{c} & \textcolor{gregoriocolor}{d} & \textcolor{gregoriocolor}{e} & \textcolor{gregoriocolor}{f} & \textcolor{gregoriocolor}{g} & \textcolor{gregoriocolor}{h} & \textcolor{gregoriocolor}{i} & \textcolor{gregoriocolor}{k} & \textcolor{gregoriocolor}{l} & \textcolor{gregoriocolor}{m} & \textcolor{gregoriocolor}{n} & \textcolor{gregoriocolor}{p} & \textcolor{gregoriocolor}{q} & \textcolor{gregoriocolor}{r} & \textcolor{gregoriocolor}{s} & \textcolor{gregoriocolor}{t} & \textcolor{gregoriocolor}{u} \\
 13 & 14 & 15 & 16 & 17 & 18 & 19 & 20 & 21 & 22 & 23 & 24 & 25 & 26 & 27 & 28 & 29 & 30 & 1 \\
\end{tabularx}
\vspace{0.5\baselineskip}
\noindent\begin{tabularx}{\linewidth}{*{12}{>{\centering\arraybackslash}X}}
 \textcolor{gregoriocolor}{A} & \textcolor{gregoriocolor}{B} & \textcolor{gregoriocolor}{C} & \textcolor{gregoriocolor}{D} & \textcolor{gregoriocolor}{E} & F & \textcolor{gregoriocolor}{F} & \textcolor{gregoriocolor}{G} & \textcolor{gregoriocolor}{H} & \textcolor{gregoriocolor}{M} & \textcolor{gregoriocolor}{N} & \textcolor{gregoriocolor}{P} \\
	2 & 3 & 4 & 5 & 6 & 7 & 7 & 8 & 9 & 10 & 11 & 12 \\
\end{tabularx}

\begin{paracol}{2}
\selectlanguage{latin}
\lettrine[lines=2]{T}{értio} Nonas Novémbris. 
 Luna\dots\
\switchcolumn
\selectlanguage{english}
\lettrine[lines=2]{T}{he} 3\textsuperscript{rd} Day of 
 November. The\dots\ Day of the Moon.
\switchcolumn*
\selectlanguage{latin}
Commemorátio ómnium 
 Fidélium Defunctórum.
\switchcolumn
\selectlanguage{english}
The Commemoration of all the 
 Faithful Departed.
\switchcolumn*
\selectlanguage{latin}
Medioláni natália 
 sancti Cároli Borromæi Cardinális, Epíscopi Mediolanénsis et Confessóris, 
 quem, sanctitáte conspícuum et miráculis clarum, Paulus Papa Quintus in 
 Sanctórum númerum rétulit. Ipsíus tamen festívitas sequénti die 
 celebrátur.
\switchcolumn
\selectlanguage{english}
At Milan, St. Charles Borromeo, 
 cardinal, bishop of that city, and confessor, who was ranked among the 
 saints by Paul V for the holiness of his life and for his renown for 
 miracles. His feast is observed on the following day.
\switchcolumn*
\selectlanguage{latin}
Eódem die natális 
 quoque sancti Quarti, Apostolórum discípuli.
\switchcolumn
\selectlanguage{english}
On the same day, the birthday of St. 
 Quartus, a disciple of the apostles.
\switchcolumn*
\selectlanguage{latin}
Vitérbii sanctórum 
 Mártyrum Valentíni Presbyteri, et Hilárii Diáconi, qui, in persecutióne 
 Maximiáni, ob Christi fidem, cum saxi póndere in Tíberim præcipitáti et inde 
 ab Angelo divínitus erépti, demum, abscíssis cervícibus, corónam martyrii 
 percepérunt.
\switchcolumn
\selectlanguage{english}
At Viterbo, during the persecution 
 of Maximian, the holy martyrs Valentine, a priest, and Hilary, a deacon. 
 For their attachment to the faith of Christ, they were tied to a stone and 
 cast into the Tiber, but being miraculously delivered by an angel, they were 
 beheaded, and thus crowned with the glory of martyrdom.
\switchcolumn*
\selectlanguage{latin}
Cæsaréæ, in Cappadócia, sanctórum Mártyrum Germáni, Theóphili, Cæsárii et Vitális; qui, in 
 persecutióne Deciána, óptime duxérunt martyrium.
\switchcolumn
\selectlanguage{english}
At Caesarea in Cappadocia, the holy 
 martyrs Germanus, Theophilus, Caesarius, and Vitalis, who nobly endured 
 martyrdom in the Decian persecution.
\switchcolumn*
\selectlanguage{latin}
Cæsaraugústæ, in 
 Hispánia, sanctórum innumerabílium Mártyrum, qui, sub Hispaniárum Præside 
 Daciáno, mirabíliter occubuérunt pro Christo.
\switchcolumn
\selectlanguage{english}
At Saragossa in Spain, the countless 
 holy martyrs who lay down their lives with admirable fervour for the faith 
 of Christ under Dacian, governor of Spain.
\switchcolumn*
\selectlanguage{latin}
In Anglia sanctæ 
 Wenefrídæ, Vírginis et Mártyris.
\switchcolumn
\selectlanguage{english}
In England, St. Winifred, virgin and 
 martyr.
\switchcolumn*
\selectlanguage{latin}
In monastério 
 Clarævallénsi, in Gállia, deposítio sancti Malachíæ, Connerthénsis in 
 Hibérnia Epíscopi, qui multis virtútibus suo témpore cláruit; cujus vitam 
 sanctum Bernárdus Abbas conscrípsit.
\switchcolumn
\selectlanguage{english}
In the monastery of Clairvaux in 
 France, the death of St. Malachy, bishop of Armagh in Ireland, who won 
 renown in his own days for his many virtues, and whose life was written by 
 St. Bernard the Abbot.
\switchcolumn*
\selectlanguage{latin}
Eódem die sancti 
 Hubérti, Tungrénsis Epíscopi.
\switchcolumn
\selectlanguage{english}
On the same day, St. Hubert, bishop 
 of Tongres.
\switchcolumn*
\selectlanguage{latin}
Viénnæ, in Gállia, 
 sancti Domni, Epíscopi et Confessóris.
\switchcolumn
\selectlanguage{english}
At Vienne in France, St. Domnus, 
 bishop and confessor.
\switchcolumn*
\selectlanguage{latin}
Item deposítio sancti 
 Pirmíni, Meldénsis Epíscopi.
\switchcolumn
\selectlanguage{english}
Also, the death of St. Pirmin, 
 bishop of Meaux.
\switchcolumn*
\selectlanguage{latin}
Urgéllæ, in Hispánia 
 Tarraconénsi, sancti Hermengáudii Epíscopi.
\switchcolumn
\selectlanguage{english}
At Urgel in Spain, Bishop St. 
 Hermengaud.
\switchcolumn*
\selectlanguage{latin}
Romæ sanctæ Sílviæ, 
 matris sancti Gregórii Papæ.
\switchcolumn
\selectlanguage{english}
At Rome, St. Sylvia, mother of Pope 
 St. Gregory.
\switchcolumn*
\selectlanguage{latin}
\end{paracol}


% ---- martyrology/mart11/mart1104.htm
\needspace{10\baselineskip}
\begin{paracol}{2}
\selectlanguage{latin}
\begin{center}{\color{gregoriocolor}Prídie Nonas 
 Novémbris, Luna\dots\ }\end{center}
\switchcolumn
\selectlanguage{english}
\begin{center}{\color{gregoriocolor}The 4\textsuperscript{th} Day of 
 November, the... Day of the Moon.}\end{center}
\end{paracol}

\noindent\begin{tabularx}{\linewidth}{*{19}{>{\centering\arraybackslash}X}}
 \textcolor{gregoriocolor}{a} & \textcolor{gregoriocolor}{b} & \textcolor{gregoriocolor}{c} & \textcolor{gregoriocolor}{d} & \textcolor{gregoriocolor}{e} & \textcolor{gregoriocolor}{f} & \textcolor{gregoriocolor}{g} & \textcolor{gregoriocolor}{h} & \textcolor{gregoriocolor}{i} & \textcolor{gregoriocolor}{k} & \textcolor{gregoriocolor}{l} & \textcolor{gregoriocolor}{m} & \textcolor{gregoriocolor}{n} & \textcolor{gregoriocolor}{p} & \textcolor{gregoriocolor}{q} & \textcolor{gregoriocolor}{r} & \textcolor{gregoriocolor}{s} & \textcolor{gregoriocolor}{t} & \textcolor{gregoriocolor}{u} \\
	14 & 15 & 16 & 17 & 18 & 19 & 20 & 21 & 22 & 23 & 24 & 25 & 26 & 27 & 28 & 29 & 30 & 1 & 2 \\
\end{tabularx}
\vspace{0.5\baselineskip}
\noindent\begin{tabularx}{\linewidth}{*{12}{>{\centering\arraybackslash}X}}
 \textcolor{gregoriocolor}{A} & \textcolor{gregoriocolor}{B} & \textcolor{gregoriocolor}{C} & \textcolor{gregoriocolor}{D} & \textcolor{gregoriocolor}{E} & F & \textcolor{gregoriocolor}{F} & \textcolor{gregoriocolor}{G} & \textcolor{gregoriocolor}{H} & \textcolor{gregoriocolor}{M} & \textcolor{gregoriocolor}{N} & \textcolor{gregoriocolor}{P} \\
	3 & 4 & 5 & 6 & 7 & 8 & 8 & 9 & 10 & 11 & 12 & 13\\
\end{tabularx}


\begin{paracol}{2}
\selectlanguage{latin}
\lettrine[lines=2]{P}{rídie} Nonas Novémbris. 
 Luna\dots\
\switchcolumn
\selectlanguage{english}
\lettrine[lines=2]{T}{he} 4\textsuperscript{th} Day of 
 November. The\dots\ Day of the Moon.
\switchcolumn*
\selectlanguage{latin}
Sancti Cároli Borromæi Cardinális, Epíscopi Mediolanénsis et Confessóris, 
 qui migrávit in cælum prídie hujus diéi.
\switchcolumn
\selectlanguage{english}
St. Charles Borromeo, 
 cardinal, bishop of Milan, and confessor, whose birthday is on the day 
 previous.
\switchcolumn*
\selectlanguage{latin}
Bonóniæ sanctórum 
 Mártyrum Vitális et Agrícolæ; quorum prior postérioris servus ántea fuit, 
 póstea consors et colléga martyrii. In ipsum porro Vitálem 
 persecutóres ómnia tormentórum génera ita exercuérunt, ut non esset in 
 córpore ejus sine vúlnere locus, quæ ipse constánter pérferens, in oratióne 
 spíritum Deo réddidit; Agrícolam vero, plúrimis clavis cruci affigéntes, 
 interemérunt. Eórum translatióni sanctus Ambrósius cum interésset, 
 Mártyris clavos, sánguinem triumphálem et crucis lignum se collegísse refert, 
 ac sub sacris altáribus condidísse.
\switchcolumn
\selectlanguage{english}
At Bologna, the holy martyrs Vitalis 
 and Agricola. The former was first the servant of the latter, and 
 afterwards his partner and companion in martyrdom. He was subjected by 
 the persecutors to all kinds of torments, so that there was no part of his 
 body without wounds. After having suffered with constancy, he yielded 
 up his soul unto God in prayer. Agricola was put to death by being 
 fastened to a cross with many nails. St. Ambrose relates that being 
 present at the translation, he took the martyr's nails, his glorious blood, 
 and the wood of his cross, and deposited them under consecrated altars.
\switchcolumn*
\selectlanguage{latin}
In cœnóbio Cervi 
 Frígidi, territórii Meldénsis, natális sancti Felícis Valésii, Presbyteri et 
 Confessóris; qui Fundátor fuit Ordinis sanctíssimæ Trinitátis redemptiónis 
 captivórum. Ipsíus autem festum, ex dispositióne Innocéntii Papæ 
 Undécimi, celebrátur duodécimo Kaléndas Decémbris.
\switchcolumn
\selectlanguage{english}
In the monastery of Cerfroid, in the 
 territory of Meaux, St. Felix of Valois, priest and confessor, and founder of the Order of the Most 
 Holy Trinity for the Redemption of Captives, whose feast is celebrated on 
 the 20th of November by order of Pope Innocent XI.
\switchcolumn*
\selectlanguage{latin}
Eódem die natális 
 sanctórum Philólogi et Pátrobi, sancti Pauli Apóstoli discipulórum.
\switchcolumn
\selectlanguage{english}
On the same day, the birthday of the 
 Saints Philologus and Patrobas, disciples of the apostle St. Paul.
\switchcolumn*
\selectlanguage{latin}
Augustodúni sancti 
 Próculi, Epíscopi et Mártyris.
\switchcolumn
\selectlanguage{english}
At Autun, St. Proculus, bishop and 
 martyr.
\switchcolumn*
\selectlanguage{latin}
Myræ, in Lycia, 
 sanctórum Mártyrum Nicándri Epíscopi, et Hermæ Presbyteri, sub Libánio 
 Præside.
\switchcolumn
\selectlanguage{english}
At Myra in Lycia, under the governor 
 Libanius, the holy martyrs Nicander, a bishop, and Hermes, a priest.
\switchcolumn*
\selectlanguage{latin}
In pago Vilcassíno, in 
 Gállia, sancti Clari, Presbyteri et Mártyris.
\switchcolumn
\selectlanguage{english}
In the district of Vexin in France, 
 St. Clarus, priest and martyr.
\switchcolumn*
\selectlanguage{latin}
Ephesi sancti Porphyrii 
 Mártyris, sub Aureliáno Imperatóre.
\switchcolumn
\selectlanguage{english}
At Ephesus, St. Porphyrias, a martyr 
 under Emperor Aurelian.
\switchcolumn*
\selectlanguage{latin}
Apud Ruthénos, in 
 Gállia, beáti Amántii Epíscopi, cujus vita éxstitit sanctitáte et miráculis 
 gloriósa.
\switchcolumn
\selectlanguage{english}
At Rodez in France, blessed Bishop 
 Amantius, whose life stood out glorious by his sanctity and miracles.
\switchcolumn*
\selectlanguage{latin}
Romæ natális sancti 
 Piérii, Presbyteri Alexandríni, qui, in divínis Scriptúris nobíliter 
 erudítus, vita puríssimus, et ad Christiánam philosophíam nudus pénitus et 
 expedítus, sub Caro et Diocletiáno Princípibus, regénte Alexandrínam 
 Ecclésiam Theóna, florentíssime dócuit pópulum, et divérsos tractátus édidit; 
 post persecutiónem vero, omne vitæ suæ tempus Romæ versátus, in pace quiévit.
\switchcolumn
\selectlanguage{english}
At Rome, the birthday of St. Pierius, 
 priest of Alexandria, who was well versed in the Holy Scriptures, most pure 
 in his life, and highly skilled in Christian philosophy. He taught the 
 people and became famous under Emperors Carus and Diocletian, when Theonas 
 governed the Church of Alexandria. After the persecution, he spent the 
 remainder of his life at Rome, where he died in peace.
\switchcolumn*
\selectlanguage{latin}
In Bithynia sancti 
 Joannícii Abbátis.
\switchcolumn
\selectlanguage{english}
In Bithynia, St. Joannicius, abbot.
\switchcolumn*
\selectlanguage{latin}
Apud Albam Regálem, in 
 Pannónia, deposítio beáti Emeríci Confessóris, qui fuit fílius sancti 
 Stéphani, Hungarórum Regis.
\switchcolumn
\selectlanguage{english}
In Hungary at Alba Regalis, the 
 death of blessed Emeric, confessor, the son of St. Stephen, king of Hungary.
\switchcolumn*
\selectlanguage{latin}
Tréviris sanctæ Modéstæ 
 Vírginis.
\switchcolumn
\selectlanguage{english}
At Treves, St. Modesta, virgin.
\switchcolumn*
\selectlanguage{latin}
\end{paracol}


% ---- martyrology/mart11/mart1105.htm
\needspace{10\baselineskip}
\begin{paracol}{2}
\selectlanguage{latin}
\begin{center}{\color{gregoriocolor} Nonis Novémbris. 
 Luna\dots\ }\end{center}
\switchcolumn
\selectlanguage{english}
\begin{center}{\color{gregoriocolor} The 5\textsuperscript{th} Day of 
 November. The\dots\ Day of the Moon.}\end{center}
\end{paracol}

\noindent\begin{tabularx}{\linewidth}{*{19}{>{\centering\arraybackslash}X}}
 \textcolor{gregoriocolor}{a} & \textcolor{gregoriocolor}{b} & \textcolor{gregoriocolor}{c} & \textcolor{gregoriocolor}{d} & \textcolor{gregoriocolor}{e} & \textcolor{gregoriocolor}{f} & \textcolor{gregoriocolor}{g} & \textcolor{gregoriocolor}{h} & \textcolor{gregoriocolor}{i} & \textcolor{gregoriocolor}{k} & \textcolor{gregoriocolor}{l} & \textcolor{gregoriocolor}{m} & \textcolor{gregoriocolor}{n} & \textcolor{gregoriocolor}{p} & \textcolor{gregoriocolor}{q} & \textcolor{gregoriocolor}{r} & \textcolor{gregoriocolor}{s} & \textcolor{gregoriocolor}{t} & \textcolor{gregoriocolor}{u} \\
 15 & 16 & 17 & 18 & 19 & 20 & 21 & 22 & 23 & 24 & 25 & 26 & 27 & 28 & 29 & 30 & 1 & 2 & 3 \\
\end{tabularx}
\vspace{0.5\baselineskip}
\noindent\begin{tabularx}{\linewidth}{*{12}{>{\centering\arraybackslash}X}}
 \textcolor{gregoriocolor}{A} & \textcolor{gregoriocolor}{B} & \textcolor{gregoriocolor}{C} & \textcolor{gregoriocolor}{D} & \textcolor{gregoriocolor}{E} & F & \textcolor{gregoriocolor}{F} & \textcolor{gregoriocolor}{G} & \textcolor{gregoriocolor}{H} & \textcolor{gregoriocolor}{M} & \textcolor{gregoriocolor}{N} & \textcolor{gregoriocolor}{P} \\
 4 & 5 & 6 & 7 & 8 & 9 & 9 & 10 & 11 & 12 & 13 & 14 \\
\end{tabularx}

\begin{paracol}{2}
\selectlanguage{latin}
\lettrine[lines=2]{S}{ancti} Zacharíæ, 
 Sacerdótis et Prophétæ, qui pater éxstitit beáti Joánnis Baptístæ, 
 Præcursóris Dómini.
\switchcolumn
\selectlanguage{english}
\lettrine[lines=2]{S}{t.} Zachary, priest and prophet, the 
 father of blessed John Baptist, Forerunner of our Lord.
\switchcolumn*
\selectlanguage{latin}
Item sanctæ Elísabeth, 
 ejúsdem sanctíssimi Præcursóris matris.
\switchcolumn
\selectlanguage{english}
Also, St. Elizabeth, mother of the 
 same most holy Forerunner.
\switchcolumn*
\selectlanguage{latin}
Tarracínæ, in Campánia, natális sanctórum Mártyrum Felícis Presbyteri, et Eusébii Mónachi. Ex 
 his Eusébius, cum sepelísset sanctos Mártyres Juliánum et Cæsárium, et 
 multos convérteret ad fidem Christi, quos sanctus Felix Presbyter baptizábat, 
 una cum ipso Felíce tentus est; et, ad Júdicis forum ducti nec superáti, 
 inde in cárcerem inclúsi, ambo, nocte eádem, cum sacrificáre noluíssent, 
 decolláti sunt.
\switchcolumn
\selectlanguage{english}
At Terracina in Campania, the 
 birthday of the holy martyrs Felix, a priest, and Eusebius, a monk. 
 The latter buried the holy martyrs Julian and Caesarius, and converted to 
 the faith of Christ many whom the priest St. Felix baptized. They were 
 arrested together, and both were led to the tribunal of the judge, who could 
 not succeed in intimidating them; they were shut up in prison, and as they 
 refused to offer sacrifice, were beheaded that same night.
\switchcolumn*
\selectlanguage{latin}
Eméssæ, in Phœnícia, 
 sanctórum Mártyrum Galatiónis et Epistémis cónjugis, qui in Décii 
 persecutióne, flagris cæsi, mánibus pedibúsque et lingua ínsuper mutiláti, 
 dénique, truncáto cápite, martyrium consummárunt.
\switchcolumn
\selectlanguage{english}
At Emesa in Phoenicia, during the 
 persecution of Decius, the holy martyrs Galation and his wife Epistemis, who 
 were scourged, had their hands, feet, and tongue mutilated, and finally 
 fulfilled their martyrdom by beheading.
\switchcolumn*
\selectlanguage{latin}
Item sanctórum Mártyrum 
 Domníni, Theótimi, Philóthei, Silváni et Sociórum, sub Maximíno Imperatóre.
\switchcolumn
\selectlanguage{english}
Also, the holy martyrs Dominus, 
 Theotimus, Philotheus, Silvanus, and their companions, under Emperor 
 Maximinus.
\switchcolumn*
\selectlanguage{latin}
Medioláni sancti Magni, 
 Epíscopi et Confessóris.
\switchcolumn
\selectlanguage{english}
At Milan, St. Magnus, bishop and 
 confessor.
\switchcolumn*
\selectlanguage{latin}
Bríxiæ sancti 
 Dominatóris Epíscopi.
\switchcolumn
\selectlanguage{english}
At Brescia, St. Dominator, bishop.
\switchcolumn*
\selectlanguage{latin}
Tréviris sancti Fibítii, 
 qui ex Abbáte factus est ejúsdem civitátis Epíscopus.
\switchcolumn
\selectlanguage{english}
At Treves, St. Fibitius, first an 
 abbot and then bishop of that city.
\switchcolumn*
\selectlanguage{latin}
Aureliánis, in Gállia, 
 sancti Læti, Presbyteri et Confessóris.
\switchcolumn
\selectlanguage{english}
At Orleans in France, St. Laetus, 
 priest and confessor.
\switchcolumn*
\selectlanguage{latin}
\end{paracol}


% ---- martyrology/mart11/mart1106.htm
\needspace{10\baselineskip}
\begin{paracol}{2}
\selectlanguage{latin}
\begin{center}{\color{gregoriocolor} Octávo Idus Novémbris. 
 Luna\dots\ }\end{center}
\switchcolumn
\selectlanguage{english}
\begin{center}{\color{gregoriocolor} The 6\textsuperscript{th} Day of 
 November. The\dots\ Day of the Moon.}\end{center}
\end{paracol}

\noindent\begin{tabularx}{\linewidth}{*{19}{>{\centering\arraybackslash}X}}
 \textcolor{gregoriocolor}{a} & \textcolor{gregoriocolor}{b} & \textcolor{gregoriocolor}{c} & \textcolor{gregoriocolor}{d} & \textcolor{gregoriocolor}{e} & \textcolor{gregoriocolor}{f} & \textcolor{gregoriocolor}{g} & \textcolor{gregoriocolor}{h} & \textcolor{gregoriocolor}{i} & \textcolor{gregoriocolor}{k} & \textcolor{gregoriocolor}{l} & \textcolor{gregoriocolor}{m} & \textcolor{gregoriocolor}{n} & \textcolor{gregoriocolor}{p} & \textcolor{gregoriocolor}{q} & \textcolor{gregoriocolor}{r} & \textcolor{gregoriocolor}{s} & \textcolor{gregoriocolor}{t} & \textcolor{gregoriocolor}{u} \\
 16 & 17 & 18 & 19 & 20 & 21 & 22 & 23 & 24 & 25 & 26 & 27 & 28 & 29 & 30 & 1 & 2 & 3 & 4 \\
\end{tabularx}
\vspace{0.5\baselineskip}
\noindent\begin{tabularx}{\linewidth}{*{12}{>{\centering\arraybackslash}X}}
 \textcolor{gregoriocolor}{A} & \textcolor{gregoriocolor}{B} & \textcolor{gregoriocolor}{C} & \textcolor{gregoriocolor}{D} & \textcolor{gregoriocolor}{E} & F & \textcolor{gregoriocolor}{F} & \textcolor{gregoriocolor}{G} & \textcolor{gregoriocolor}{H} & \textcolor{gregoriocolor}{M} & \textcolor{gregoriocolor}{N} & \textcolor{gregoriocolor}{P} \\
 5 & 6 & 7 & 8 & 9 & 10 & 10 & 11 & 12 & 13 & 14 & 15 \\
\end{tabularx}

\begin{paracol}{2}
\selectlanguage{latin}
\lettrine[lines=2]{B}{arcinóne,} in Hispánia, 
 sancti Sevéri, Epíscopi et Mártyris; qui ob fidem cathólicam, confósso per 
 clavum cápite, martyrii corónam accépit.
\switchcolumn
\selectlanguage{english}
\lettrine[lines=2]{A}{t} Barcelona in Spain, St. Severus, 
 bishop and martyr, who had his head pierced with a spike, and thus received 
 the crown of martyrdom for the sake of the Catholic faith.
\switchcolumn*
\selectlanguage{latin}
Thiníssæ, in Africa, 
 natális sancti Felícis Mártyris, qui, conféssus et ad torménta dilátus, álio 
 die (ut refert sanctus Augustínus, Psalmum in ejus festivitáte ad pópulum 
 expónens) invéntus est in cárcere exánimis.
\switchcolumn
\selectlanguage{english}
At Tunis in Africa, the birthday of 
 St. Felix, martyr, who, having confessed Christ, was sent to prison. 
 His sentence had been deferred, but the next day he was found dead, as is 
 related by St. Augustine when he was expounding on a psalm to the people on 
 the feast of the saint.
\switchcolumn*
\selectlanguage{latin}
Theópoli, quæ est 
 Antiochía, sanctórum decem Mártyrum, qui a Saracénis passi legúntur.
\switchcolumn
\selectlanguage{english}
At Theopolis, which is Antioch, ten 
 holy martyrs who are said to have suffered at the hands of the Saracens.
\switchcolumn*
\selectlanguage{latin}
In Phrygia sancti 
 Attici Mártyris.
\switchcolumn
\selectlanguage{english}
In Phrygia, St. Atticus, martyr.
\switchcolumn*
\selectlanguage{latin}
Apud Bergas, in 
 Flándria, deposítio sancti Winóci Abbátis, qui, virtútibus et miráculis 
 clarus, étiam frátribus sibi súbditis multo témpore ministrávit.
\switchcolumn
\selectlanguage{english}
At Berg in Flanders, the death of 
 St. Winoc, abbot, who was renowned for virtues and miracles, and for a long 
 time was servant to the brethren subject to him.
\switchcolumn*
\selectlanguage{latin}
Fundis, in Látio, 
 sancti Felícis Mónachi.
\switchcolumn
\selectlanguage{english}
At Fondi in Lazio, St. Felix, monk.
\switchcolumn*
\selectlanguage{latin}
Lemóvicis, in Aquitánia, sancti Leonárdi Confessóris, qui fuit beáti Remígii Epíscopi discípulus. 
 Hic, nóbili génere ortus, solitáriam vitam delégit, et sanctitáte ac 
 miráculis cláruit; ejúsque virtus præcípue in liberándis captívis enítuit.
\switchcolumn
\selectlanguage{english}
At Limoges in Aquitaine, St. 
 Leonard, confessor, disciple of the blessed bishop Remigius, who was born of 
 a noble family but chose to lead a solitary life. He was celebrated 
 for holiness and miracles, but his virtue shone particularly in the 
 deliverance of captives.
\switchcolumn*
\selectlanguage{latin}
\end{paracol}


% ---- martyrology/mart11/mart1107.htm
\needspace{10\baselineskip}
\begin{paracol}{2}
\selectlanguage{latin}
\begin{center}{\color{gregoriocolor} Séptimo Idus Novémbris. 
 Luna\dots\ }\end{center}
\switchcolumn
\selectlanguage{english}
\begin{center}{\color{gregoriocolor} The 7\textsuperscript{th} Day of 
 November. The\dots\ Day of the Moon.}\end{center}
\end{paracol}

\noindent\begin{tabularx}{\linewidth}{*{19}{>{\centering\arraybackslash}X}}
 \textcolor{gregoriocolor}{a} & \textcolor{gregoriocolor}{b} & \textcolor{gregoriocolor}{c} & \textcolor{gregoriocolor}{d} & \textcolor{gregoriocolor}{e} & \textcolor{gregoriocolor}{f} & \textcolor{gregoriocolor}{g} & \textcolor{gregoriocolor}{h} & \textcolor{gregoriocolor}{i} & \textcolor{gregoriocolor}{k} & \textcolor{gregoriocolor}{l} & \textcolor{gregoriocolor}{m} & \textcolor{gregoriocolor}{n} & \textcolor{gregoriocolor}{p} & \textcolor{gregoriocolor}{q} & \textcolor{gregoriocolor}{r} & \textcolor{gregoriocolor}{s} & \textcolor{gregoriocolor}{t} & \textcolor{gregoriocolor}{u} \\
 17 & 18 & 19 & 20 & 21 & 22 & 23 & 24 & 25 & 26 & 27 & 28 & 29 & 30 & 1 & 2 & 3 & 4 & 5 \\
\end{tabularx}
\vspace{0.5\baselineskip}
\noindent\begin{tabularx}{\linewidth}{*{12}{>{\centering\arraybackslash}X}}
 \textcolor{gregoriocolor}{A} & \textcolor{gregoriocolor}{B} & \textcolor{gregoriocolor}{C} & \textcolor{gregoriocolor}{D} & \textcolor{gregoriocolor}{E} & F & \textcolor{gregoriocolor}{F} & \textcolor{gregoriocolor}{G} & \textcolor{gregoriocolor}{H} & \textcolor{gregoriocolor}{M} & \textcolor{gregoriocolor}{N} & \textcolor{gregoriocolor}{P} \\
 6 & 7 & 8 & 9 & 10 & 11 & 11 & 12 & 13 & 14 & 15 & 16 \\
\end{tabularx}

\begin{paracol}{2}
\selectlanguage{latin}
\lettrine[lines=2]{P}{atávii} deposítio 
 sancti Prosdócimi, qui fuit primus ejúsdem civitátis Epíscopus. Hic, a 
 beáto Petro Apóstolo Epíscopus ordinátus, ad prædicándum Dei verbum ad 
 prædíctam civitátem missus est; ibíque, multis virtútibus et prodígiis 
 corúscans, beáto fine quiévit.
\switchcolumn
\selectlanguage{english}
\lettrine[lines=2]{A}{t} Padua, the death of St. 
 Prosdocimus, consecrated as first bishop of that city by the blessed apostle 
 Peter. He was sent to that city to preach the word of God and there he 
 died a holy death, celebrated for many virtues and miracles.
\switchcolumn*
\selectlanguage{latin}
Perúsiæ sancti 
 Herculáni, Epíscopi et Mártyris.
\switchcolumn
\selectlanguage{english}
At Perugia, St. Herculanus, bishop 
 and martyr.
\switchcolumn*
\selectlanguage{latin}
Apud Swelménsem 
 civitátem, in Germánia, pássio sancti Engelbérti, Epíscopi Coloniénsis, qui, 
 cum illuc ex óppido Sosátio ad templum dedicándum pérgeret, a sicáriis 
 intercéptus in via multísque vulnéribus cæsus, gloriósum pro defensióne 
 ecclesiásticæ libertátis et Románæ Ecclésiæ obediéntia martyrium súbiit.
\switchcolumn
\selectlanguage{english}
At Schwelm in Germany, the martyrdom 
 of St. Engelbert, bishop of Cologne. He was on his way from that city 
 to the town of Essen in order to consecrate a church, when he was set upon 
 by ruffians on the road and slain by their many blows. Thus he 
 suffered martyrdom in defence of Church liberty and for obedience to the 
 Roman Church.
\switchcolumn*
\selectlanguage{latin}
Eódem die sancti 
 Amaránthi Mártyris, qui apud Albigénsem urbem, in Gállia, exácto agónis 
 fidélis cursu, sepúltus, vivit in glória.
\switchcolumn
\selectlanguage{english}
The same day, St. Amaranthus, 
 martyr. After successfully fulfilling the course of his trials he was 
 buried in the city of Albi, but lives in eternal glory.
\switchcolumn*
\selectlanguage{latin}
Melitínæ, in Arménia, pássio sanctórum Hierónis, Nicándri, Hesychii et aliórum trigínta; qui in 
 persecutióne Diocletiáni, sub Lysia Præside, coronáti sunt.
\switchcolumn
\selectlanguage{english}
At Melitina in Armenia, the 
 martyrdom of the Saints Hiero, Nicander, Hesychius, and thirty others, who 
 were crowned in the persecution of Diocletian under the governor Lysias.
\switchcolumn*
\selectlanguage{latin}
Amphípoli, in 
 Macedónia, sanctórum Mártyrum Aucti, Tauriónis et Thessalonícæ.
\switchcolumn
\selectlanguage{english}
At Amphipolis in Macedonia, the holy 
 martyrs Auctus, Taurio, and Thessalonica.
\switchcolumn*
\selectlanguage{latin}
Ancyræ, in Galátia, pássio sanctórum Melasíppi, Antónii et Carínæ, sub Juliáno Apóstata.
\switchcolumn
\selectlanguage{english}
At Ancyra in Galatia, the martyrdom 
 of Saints Melasippus, Anthony and Carina, under Julian the Apostate.
\switchcolumn*
\selectlanguage{latin}
Alexandríæ beáti 
 Achíllæ Epíscopi, qui eruditióne, fide, conversatióne ac móribus fuit 
 insígnis.
\switchcolumn
\selectlanguage{english}
At Alexandria, the blessed Achilles, 
 bishop, renowned for his learning, faith, and purity of life.
\switchcolumn*
\selectlanguage{latin}
In Frísia deposítio 
 sancti Willibrórdi, Epíscopi Trajecténsis; qui, a beáto Sérgio Papa ordinátus Epíscopus, in Frísia et Dánia Evangélium prædicávit.
\switchcolumn
\selectlanguage{english}
In Friesland, the death of St. 
 Willibrord, bishop of Utrecht, who was consecrated bishop by blessed Pope 
 Sergius, and preached the Gospel in Friesland and Denmark.
\switchcolumn*
\selectlanguage{latin}
Metis, in Gállia, 
 sancti Rufi, Epíscopi et Confessóris.
\switchcolumn
\selectlanguage{english}
At Metz, St. Rufus, bishop and 
 confessor.
\switchcolumn*
\selectlanguage{latin}
Argentoráti sancti 
 Floréntii Epíscopi.
\switchcolumn
\selectlanguage{english}
At Strasbourg, St. Florentius, 
 bishop.
\switchcolumn*
\selectlanguage{latin}
\end{paracol}


% ---- martyrology/mart11/mart1108.htm
\needspace{10\baselineskip}
\begin{paracol}{2}
\selectlanguage{latin}
\begin{center}{\color{gregoriocolor} Sexto Idus Novémbris. 
 Luna\dots\ }\end{center}
\switchcolumn
\selectlanguage{english}
\begin{center}{\color{gregoriocolor} The 8\textsuperscript{th} Day of 
 November. The\dots\ Day of the Moon.}\end{center}
\end{paracol}

\noindent\begin{tabularx}{\linewidth}{*{19}{>{\centering\arraybackslash}X}}
 \textcolor{gregoriocolor}{a} & \textcolor{gregoriocolor}{b} & \textcolor{gregoriocolor}{c} & \textcolor{gregoriocolor}{d} & \textcolor{gregoriocolor}{e} & \textcolor{gregoriocolor}{f} & \textcolor{gregoriocolor}{g} & \textcolor{gregoriocolor}{h} & \textcolor{gregoriocolor}{i} & \textcolor{gregoriocolor}{k} & \textcolor{gregoriocolor}{l} & \textcolor{gregoriocolor}{m} & \textcolor{gregoriocolor}{n} & \textcolor{gregoriocolor}{p} & \textcolor{gregoriocolor}{q} & \textcolor{gregoriocolor}{r} & \textcolor{gregoriocolor}{s} & \textcolor{gregoriocolor}{t} & \textcolor{gregoriocolor}{u} \\
 18 & 19 & 20 & 21 & 22 & 23 & 24 & 25 & 26 & 27 & 28 & 29 & 30 & 1 & 2 & 3 & 4 & 5 & 6 \\
\end{tabularx}
\vspace{0.5\baselineskip}
\noindent\begin{tabularx}{\linewidth}{*{12}{>{\centering\arraybackslash}X}}
 \textcolor{gregoriocolor}{A} & \textcolor{gregoriocolor}{B} & \textcolor{gregoriocolor}{C} & \textcolor{gregoriocolor}{D} & \textcolor{gregoriocolor}{E} & F & \textcolor{gregoriocolor}{F} & \textcolor{gregoriocolor}{G} & \textcolor{gregoriocolor}{H} & \textcolor{gregoriocolor}{M} & \textcolor{gregoriocolor}{N} & \textcolor{gregoriocolor}{P} \\
 7 & 8 & 9 & 10 & 11 & 12 & 12 & 13 & 14 & 15 & 16 & 17 \\
\end{tabularx}

\begin{paracol}{2}
\selectlanguage{latin}
\lettrine[lines=1]{O}{ctáva} ómnium 
 Sanctórum.
\switchcolumn
\selectlanguage{english}
\lettrine[lines=1]{T}{he} Octave of All Saints.
\switchcolumn*
\selectlanguage{latin}
Romæ, via Lavicána, 
 tértio ab Urbe milliário, pássio sanctórum Mártyrum Cláudii, Nicóstrati, 
 Symphoriáni, Castórii et Simplícii, qui, primo in cárcerem missi, deínde 
 scorpiónibus gravíssime cæsi, tandem, cum ex fide Christi dimovéri non 
 possent, a Diocletiáno jussi sunt in flúvium præcípites dari.
\switchcolumn
\selectlanguage{english}
At Rome, on the Lavican Way, three 
 miles from the city, the martyrdom of the Saints Claudius, Nicostratus, 
 Symphorian, Castorius, and Simplicius. They were first sent to prison, 
 then scourged with whips set with metal, but since they could not be made to 
 forsake the faith of Christ, Diocletian ordered them to be thrown into the 
 river.
\switchcolumn*
\selectlanguage{latin}
Ibídem, via Lavicána, 
 natális sanctórum Quátuor Coronatórum fratrum, id est Sevéri, Severiáni, 
 Carpóphori et Victoríni; qui, sub eódem Imperatóre, íctibus plumbatárum 
 usque ad mortem cæsi sunt. Horum autem nómina, quæ póstea, interjéctis 
 annis, Dómino revelánte, osténsa sunt, cum mínime reperíri tunc potuíssent, 
 statútum fuit ut anniversária dies ipsórum, una cum illis quinque, sub 
 nómine sanctórum Quátuor Coronatórum recolerétur; qui mos, étiam postquam 
 reveláta sunt, in Ecclésia perseverávit.
\switchcolumn
\selectlanguage{english}
Also, on the Lavican Way, the 
 birthday of the saintly brothers, Severus, Severian, Carpophorus, and 
 Victorinus, called the Four Crowned, who were scourged to death with leaded 
 whips, during the reign of the same emperor. Because their names, 
 known some years afterwards by revelation, could not then be ascertained, it 
 was ordered that their anniversary should be commemorated with the preceding 
 five, under the name of the Four Saints Crowned. This custom was 
 retained by the Church, even after their names had been revealed.
\switchcolumn*
\selectlanguage{latin}
Item Romæ sancti 
 Deúsdedit Papæ Primi, qui tanti mériti fuit, ut leprósum ósculo a lepra 
 sanáverit.
\switchcolumn
\selectlanguage{english}
Also at Rome, St. Deusdedit, pope, 
 whose merit was so great that he cured a leper by kissing him.
\switchcolumn*
\selectlanguage{latin}
In vico Blexen, ad 
 Visúrgim flúvium, in Germánia, sancti Willehádi, qui primus éxstitit 
 Breménsis civitátis Epíscopus; atque, una cum sancto Bonifátio, cujus 
 discípulus fuit, in Frísia et Saxónia Evangélium propagávit.
\switchcolumn
\selectlanguage{english}
In the village of Plexem, on the 
 Weser River in Germany, St. Willehad, first bishop of Bremen, who, together 
 with St. Boniface, whose disciple he was, spread the Gospel in Friesland and 
 Saxony.
\switchcolumn*
\selectlanguage{latin}
Suessíone, in Gálliis, 
 sancti Godefrídi, Ambianénsis Epíscopi, magnæ sanctitátis viri.
\switchcolumn
\selectlanguage{english}
At Soissons in France, St. Godfrey, 
 bishop of Amiens, a man of great sanctity.
\switchcolumn*
\selectlanguage{latin}
Apud Virodúnum, in 
 Gállia, sancti Mauri, Epíscopi et Confessóris.
\switchcolumn
\selectlanguage{english}
At Verdun in France, St. Maur, 
 bishop and confessor.
\switchcolumn*
\selectlanguage{latin}
Turónis, in Gállia, 
 sancti Clári Presbyteri, cujus sanctus Paulínus epitáphium scripsit.
\switchcolumn
\selectlanguage{english}
At Tours in France, St. Clarus, a 
 priest whose epitaph was written by St. Paulinus.
\switchcolumn*
\selectlanguage{latin}
\end{paracol}


% ---- martyrology/mart11/mart1109.htm
\needspace{10\baselineskip}
\begin{paracol}{2}
\selectlanguage{latin}
\begin{center}{\color{gregoriocolor} Quinto Idus Novémbris. 
 Luna\dots\ }\end{center}
\switchcolumn
\selectlanguage{english}
\begin{center}{\color{gregoriocolor} The 9\textsuperscript{th} Day of 
 November. The\dots\ Day of the Moon.}\end{center}
\end{paracol}

\noindent\begin{tabularx}{\linewidth}{*{19}{>{\centering\arraybackslash}X}}
 \textcolor{gregoriocolor}{a} & \textcolor{gregoriocolor}{b} & \textcolor{gregoriocolor}{c} & \textcolor{gregoriocolor}{d} & \textcolor{gregoriocolor}{e} & \textcolor{gregoriocolor}{f} & \textcolor{gregoriocolor}{g} & \textcolor{gregoriocolor}{h} & \textcolor{gregoriocolor}{i} & \textcolor{gregoriocolor}{k} & \textcolor{gregoriocolor}{l} & \textcolor{gregoriocolor}{m} & \textcolor{gregoriocolor}{n} & \textcolor{gregoriocolor}{p} & \textcolor{gregoriocolor}{q} & \textcolor{gregoriocolor}{r} & \textcolor{gregoriocolor}{s} & \textcolor{gregoriocolor}{t} & \textcolor{gregoriocolor}{u} \\
 19 & 20 & 21 & 22 & 23 & 24 & 25 & 26 & 27 & 28 & 29 & 30 & 1 & 2 & 3 & 4 & 5 & 6 & 7 \\
\end{tabularx}
\vspace{0.5\baselineskip}
\noindent\begin{tabularx}{\linewidth}{*{12}{>{\centering\arraybackslash}X}}
 \textcolor{gregoriocolor}{A} & \textcolor{gregoriocolor}{B} & \textcolor{gregoriocolor}{C} & \textcolor{gregoriocolor}{D} & \textcolor{gregoriocolor}{E} & F & \textcolor{gregoriocolor}{F} & \textcolor{gregoriocolor}{G} & \textcolor{gregoriocolor}{H} & \textcolor{gregoriocolor}{M} & \textcolor{gregoriocolor}{N} & \textcolor{gregoriocolor}{P} \\
 8 & 9 & 10 & 11 & 12 & 13 & 13 & 14 & 15 & 16 & 17 & 18 \\
\end{tabularx}

\begin{paracol}{2}
\selectlanguage{latin}
\lettrine[lines=2]{R}{omæ,} in Lateráno, 
 Dedicátio Basílicæ sanctíssimi Salvatóris, quæ ómnium Urbis et Orbis 
 Ecclesiárum est mater et caput.
\switchcolumn
\selectlanguage{english}
\lettrine[lines=2]{A}{t} Rome in the Lateran, the 
 Dedication of the Basilica of the Saviour, which is the Mother and Head of 
 all churches in the city and the world.
\switchcolumn*
\selectlanguage{latin}
Amaséæ, in Ponto, 
 natális sancti Theodóri mílitis, qui, témpore Maximiáni Imperatóris, pro 
 Christiánæ fídei confessióne, fórtiter cæsus et in cárcerem missus; deínde, 
 Dómino sibi apparénte ac monénte ut constánter et viríliter ágeret, 
 relevátus est; novíssime, postquam in equúleo suspénsus et úngulis ita 
 excarnificátus est, ut ejus interióra apparérent nuda, ardéntibus ígnibus 
 comburéndus tráditur. Ipsíus vero laudes sanctus Gregórius Nyssénus 
 præcláro encómio celebrávit.
\switchcolumn
\selectlanguage{english}
At Amasea in Pontus, the birthday of 
 St. Theodore, a soldier, in the time of Emperor Maximian. For the 
 confession of Christ he was severely scourged and sent to prison, where he 
 was comforted by an apparition of our Lord, who exhorted him to act with 
 courage and constancy. He was finally stretched on the rack, lacerated 
 with iron hooks until his bowels were laid bare, then cast into the flames 
 to be burned alive. His glorious deeds have been celebrated in a 
 eulogy by Gregory of Nyssa.
\switchcolumn*
\selectlanguage{latin}
Tyánæ, in Cappadócia, pássio sancti Oréstis, sub Diocletiáno Imperatóre.
\switchcolumn
\selectlanguage{english}
At Tyana in Cappadocia, the 
 martyrdom of St. Orestes under Emperor Diocletian.
\switchcolumn*
\selectlanguage{latin}
Thessalonícæ sancti 
 Alexándri Mártyris, sub Maximiáno Príncipe.
\switchcolumn
\selectlanguage{english}
At Thessalonica, under Emperor 
 Maximian, St. Alexander, martyr.
\switchcolumn*
\selectlanguage{latin}
Apud Bitúricas, in 
 Aquitánia, sancti Ursíni Confessóris, qui, Romæ ordinátus a successóribus 
 Apostolórum, primus eídem Bituricénsi urbi destinátur Epíscopus.
\switchcolumn
\selectlanguage{english}
At Bourges in Aquitaine, St. Ursinus, 
 confessor, who was ordained at Rome by the successors of the apostles and 
 appointed first bishop of that city.
\switchcolumn*
\selectlanguage{latin}
Neápoli, in Campánia, sancti Agrippíni Epíscopi, miráculis clari.
\switchcolumn
\selectlanguage{english}
At Naples in Campania, St. 
 Agrippinus, bishop, renowned for miracles.
\switchcolumn*
\selectlanguage{latin}
Constantinópoli 
 sanctárum Vírginum Eustóliæ Románæ, et Sópatræ, fíliæ Maurítii Imperatóris.
\switchcolumn
\selectlanguage{english}
At Constantinople, the holy virgins 
 Eustolia, a Roman maiden, and Sopatra, the daughter of Emperor Maurice
\switchcolumn*
\selectlanguage{latin}
Beryti, in Syria, 
 commemorátio Imáginis Salvatóris, quæ, a Judæis crucifíxa, tam copiósum 
 emísit sánguinem, ut Orientáles et Occidentáles Ecclésiæ ex eo ubértim 
 accéperint.
\switchcolumn
\selectlanguage{english}
At Berytus in Syria, the 
 Commemoration of the Image of our Saviour, which, being fastened to a cross 
 by the Jews, poured out blood so plentifully that the Eastern and Western 
 Churches received abundantly of it.
\switchcolumn*
\selectlanguage{latin}
\end{paracol}


% ---- martyrology/mart11/mart1110.htm
\needspace{10\baselineskip}
\begin{paracol}{2}
\selectlanguage{latin}
\begin{center}{\color{gregoriocolor} Quarto Idus Novémbris. 
 Luna\dots\ }\end{center}
\switchcolumn
\selectlanguage{english}
\begin{center}{\color{gregoriocolor} The 10\textsuperscript{th} Day of 
 November. The\dots\ Day of the Moon.}\end{center}
\end{paracol}

\noindent\begin{tabularx}{\linewidth}{*{19}{>{\centering\arraybackslash}X}}
 \textcolor{gregoriocolor}{a} & \textcolor{gregoriocolor}{b} & \textcolor{gregoriocolor}{c} & \textcolor{gregoriocolor}{d} & \textcolor{gregoriocolor}{e} & \textcolor{gregoriocolor}{f} & \textcolor{gregoriocolor}{g} & \textcolor{gregoriocolor}{h} & \textcolor{gregoriocolor}{i} & \textcolor{gregoriocolor}{k} & \textcolor{gregoriocolor}{l} & \textcolor{gregoriocolor}{m} & \textcolor{gregoriocolor}{n} & \textcolor{gregoriocolor}{p} & \textcolor{gregoriocolor}{q} & \textcolor{gregoriocolor}{r} & \textcolor{gregoriocolor}{s} & \textcolor{gregoriocolor}{t} & \textcolor{gregoriocolor}{u} \\
 20 & 21 & 22 & 23 & 24 & 25 & 26 & 27 & 28 & 29 & 30 & 1 & 2 & 3 & 4 & 5 & 6 & 7 & 8 \\
\end{tabularx}
\vspace{0.5\baselineskip}
\noindent\begin{tabularx}{\linewidth}{*{12}{>{\centering\arraybackslash}X}}
 \textcolor{gregoriocolor}{A} & \textcolor{gregoriocolor}{B} & \textcolor{gregoriocolor}{C} & \textcolor{gregoriocolor}{D} & \textcolor{gregoriocolor}{E} & F & \textcolor{gregoriocolor}{F} & \textcolor{gregoriocolor}{G} & \textcolor{gregoriocolor}{H} & \textcolor{gregoriocolor}{M} & \textcolor{gregoriocolor}{N} & \textcolor{gregoriocolor}{P} \\
 9 & 10 & 11 & 12 & 13 & 14 & 14 & 15 & 16 & 17 & 18 & 19 \\
\end{tabularx}

\begin{paracol}{2}
\selectlanguage{latin}
\lettrine[lines=2]{N}{eápoli,} in Campánia, natális sancti Andréæ Avellini, Clérici Reguláris et Confessóris, sanctitáte 
 et salútis proximórum procurándæ stúdio præcélebris, quem, miráculis clarum, 
 Clemens Undécimus, Póntifex Máximus, Sanctórum catálogo adscrípsit.
\switchcolumn
\selectlanguage{english}
\lettrine[lines=2]{A}{t} Naples in Campania, the birthday 
 of St. Andrew Avellini, Cleric Regular and confessor, celebrated for his 
 sanctity, his zeal in procuring the salvation of souls, and renowned for his 
 miracles. He was inscribed on the catalogue of the Saints by Pope 
 Clement XI.
\switchcolumn*
\selectlanguage{latin}
Eódem die natális 
 quoque sanctórum Mártyrum Tryphónis, et Respícii, ac Nymphæ Vírginis.
\switchcolumn
\selectlanguage{english}
On the same day, the birthday of the 
 holy martyrs Trypho and Respicius, and the virgin Nympha.
\switchcolumn*
\selectlanguage{latin}
Romæ item natális 
 sancti Leónis Papæ Primi, Confessóris et Ecclésiæ Doctóris; qui, virtútum 
 excéllens méritis, dictus est Magnus. Ejus tempóribus celebráta fuit 
 sancta Synodus Chalcedonénsis, in qua ipse per legátos damnávit Eutychen; 
 cujus étiam Synodi decréta póstmodum auctoritáte sua confirmávit. 
 Tandem, cum sanxísset multa luculentérque scripsísset, Pastor bonus, de 
 sancta Dei Ecclésia et univérso grege Domínico óptime méritus, quiévit in 
 pace. Ipsíus tamen festívitas tértio Idus Aprílis celebrátur.
\switchcolumn
\selectlanguage{english}
At Rome, Pope St. Leo I, confessor 
 and doctor of the Church, surnamed the Great because of his extraordinary 
 merits. During his pontificate the holy Council of Chalcedon was held 
 which condemned Eutyches through his legates, and whose decrees were 
 afterwards given the seal of his authority. After meriting the 
 gratitude of the Church of God and the whole flock of Christ by the many 
 decrees which he issued, and by the many excellent works which he wrote, 
 this good and zealous shepherd rested in peace. His feast is 
 celebrated on the 11th of April.
\switchcolumn*
\selectlanguage{latin}
Icónii, in Lycaónia, 
 sanctárum mulíerum Tryphénnæ et Tryphósæ, quæ, beáti Pauli prædicatióne et 
 exémplo Theclæ, in Christiána disciplína plúrimum profecérunt.
\switchcolumn
\selectlanguage{english}
At Iconium in Lycaonia, the holy 
 women Tryphenna and Tryphosa, who profited by the preaching of blessed Paul 
 and the example of Thecla to make great progress in Christian perfection.
\switchcolumn*
\selectlanguage{latin}
Antiochíæ sanctórum 
 Demétrii Epíscopi, Aniáni Diáconi, Eustósii et aliórum vigínti Mártyrum.
\switchcolumn
\selectlanguage{english}
At Antioch, Saints Demetrius, 
 bishop, Anian, deacon, Eustosius, and twenty other martyrs.
\switchcolumn*
\selectlanguage{latin}
In território Agathénsi, 
 in Gállia, sanctórum Mártyrum Tibérii, Modésti et Floréntiæ; qui, témpore 
 Diocletiáni, váriis torméntis cruciáti, martyrium complevérunt.
\switchcolumn
\selectlanguage{english}
In the diocese of Agde in France, 
 the holy martyrs Tiberius, Modestus, and Florence, who were subjected to 
 diverse torments and fulfilled their martyrdom in the time of Diocletian.
\switchcolumn*
\selectlanguage{latin}
Ravénnæ sancti Probi 
 Epíscopi, miráculis clari.
\switchcolumn
\selectlanguage{english}
At Ravenna, St. Probus, a bishop 
 renowned for miracles.
\switchcolumn*
\selectlanguage{latin}
Aureliánis, in Gállia, 
 sancti Monitóris, Epíscopi et Confessóris.
\switchcolumn
\selectlanguage{english}
At Orleans in France, St. Monitor, 
 bishop and confessor.
\switchcolumn*
\selectlanguage{latin}
In Anglia sancti Justi 
 Epíscopi, qui, una cum Augustíno, Mellíto et áliis a beáto Gregório Papa in 
 eam ínsulam missus ad prædicándum Evangélium, ibídem, sanctitáte célebris, 
 obdormívit in Dómino.
\switchcolumn
\selectlanguage{english}
In England, St. Justus, bishop, who 
 was sent by Pope Gregory with Augustine, Mellitus, and others to preach the 
 Gospel in that country. There he went to repose in the Lord, 
 celebrated for his sanctity.
\switchcolumn*
\selectlanguage{latin}
In óppido Milledúno, in 
 Gállia, sancti Leónis Confessóris.
\switchcolumn
\selectlanguage{english}
In the town of Melun in France, St. 
 Leo, confessor.
\switchcolumn*
\selectlanguage{latin}
In Paro ínsula sanctæ 
 Theoctístis Vírginis.
\switchcolumn
\selectlanguage{english}
In the island of Paros, St. 
 Theoctistis, virgin.
\switchcolumn*
\selectlanguage{latin}
\end{paracol}


% ---- martyrology/mart11/mart1111.htm
\needspace{10\baselineskip}
\begin{paracol}{2}
\selectlanguage{latin}
\begin{center}{\color{gregoriocolor} Tértio Idus Novémbris. 
 Luna\dots\ }\end{center}
\switchcolumn
\selectlanguage{english}
\begin{center}{\color{gregoriocolor} The 
 11\textsuperscript{th} Day of 
 November. The\dots\ Day of the Moon.}\end{center}
\end{paracol}

\noindent\begin{tabularx}{\linewidth}{*{19}{>{\centering\arraybackslash}X}}
 \textcolor{gregoriocolor}{a} & \textcolor{gregoriocolor}{b} & \textcolor{gregoriocolor}{c} & \textcolor{gregoriocolor}{d} & \textcolor{gregoriocolor}{e} & \textcolor{gregoriocolor}{f} & \textcolor{gregoriocolor}{g} & \textcolor{gregoriocolor}{h} & \textcolor{gregoriocolor}{i} & \textcolor{gregoriocolor}{k} & \textcolor{gregoriocolor}{l} & \textcolor{gregoriocolor}{m} & \textcolor{gregoriocolor}{n} & \textcolor{gregoriocolor}{p} & \textcolor{gregoriocolor}{q} & \textcolor{gregoriocolor}{r} & \textcolor{gregoriocolor}{s} & \textcolor{gregoriocolor}{t} & \textcolor{gregoriocolor}{u} \\
 21 & 22 & 23 & 24 & 25 & 26 & 27 & 28 & 29 & 30 & 1 & 2 & 3 & 4 & 5 & 6 & 7 & 8 & 9 \\
\end{tabularx}
\vspace{0.5\baselineskip}
\noindent\begin{tabularx}{\linewidth}{*{12}{>{\centering\arraybackslash}X}}
 \textcolor{gregoriocolor}{A} & \textcolor{gregoriocolor}{B} & \textcolor{gregoriocolor}{C} & \textcolor{gregoriocolor}{D} & \textcolor{gregoriocolor}{E} & F & \textcolor{gregoriocolor}{F} & \textcolor{gregoriocolor}{G} & \textcolor{gregoriocolor}{H} & \textcolor{gregoriocolor}{M} & \textcolor{gregoriocolor}{N} & \textcolor{gregoriocolor}{P} \\
 10 & 11 & 12 & 13 & 14 & 15 & 15 & 16 & 17 & 18 & 19 & 20 \\
\end{tabularx}

\begin{paracol}{2}
\selectlanguage{latin}
\lettrine[lines=2]{T}{urónis,} in Gállia, 
 natális beáti Martíni Epíscopi et Confessóris; cujus vita tantis éxstitit 
 miráculis gloriósa, ut trium mortuórum suscitátor esse merúerit.
\switchcolumn
\selectlanguage{english}
\lettrine[lines=2]{A}{t} Tours in France, the birthday of 
 blessed Martin, bishop and confessor, whose life was so renowned for 
 miracles that he received the power to raise three persons from the dead.
\switchcolumn*
\selectlanguage{latin}
Cotyǽi, in Phrygia, 
 insígnis pássio sancti Mennæ, Ægyptii mílitis, qui in persecutióne 
 Diocletiáni, postquam, abjécto milítiæ cíngulo, méruit cælésti Regi secréta 
 conversatióne in erémo militáre, procéssit in públicum, et, se Christiánum 
 líbera voce declárans, primo diris cruciátibus examinátur; novíssime, fixis 
 in oratióne génibus, Dómino Jesu Christo grátias agens, gládio cæsus est, ac 
 multis post mortem miráculis cláruit.
\switchcolumn
\selectlanguage{english}
At Cotyaeum in Phrygia, during the 
 persecution of Diocletian, the celebrated martyrdom of St. Mennas, an 
 Egyptian soldier, who cast off the military belt and obtained the grace of 
 serving the King of heaven secretly in the desert. Afterwards, coming 
 out publicly and freely declaring himself a Christian, he was first 
 subjected to severe torments; and finally kneeling in prayer, giving thanks 
 to our Lord Jesus Christ, he was slain with the sword. After his death 
 he became renowned for many miracles.
\switchcolumn*
\selectlanguage{latin}
Ravénnæ sanctórum 
 Mártyrum Valentíni, Feliciáni et Victoríni; qui in Diocletiáni persecutióne 
 coronáti sunt.
\switchcolumn
\selectlanguage{english}
At Ravenna, the holy martyrs 
 Valentine, Felician, and Victorinus, who were crowned during the persecution 
 of Diocletian.
\switchcolumn*
\selectlanguage{latin}
In Mesopotámia sancti 
 Athenodóri Mártyris, qui, sub eódem Diocletiáno et Eléusio Præside, ígnibus 
 cruciátus et áliis supplíciis tortus, demum cápitis damnátus est, et, cum 
 cárnifex corruísset neque ullus álius gládio illum feríre ausus esset, orans 
 obdormívit in Dómino.
\switchcolumn
\selectlanguage{english}
In Mesopotamia, St. Athenodorus, 
 martyr, who was subjected to fire and other torments under the same 
 Diocletian and the governor Eleusius. He was at length sentenced to be 
 beheaded, but when the executioner fell to the ground and no other person 
 would dare to strike him with the sword, he passed to his repose in the Lord 
 while praying.
\switchcolumn*
\selectlanguage{latin}
Lugdúni, in Gállia, 
 sancti Veráni Epíscopi, cujus vita fuit fide et virtútum méritis illústris.
\switchcolumn
\selectlanguage{english}
At Lyons in France, St. Veranus, 
 bishop, whose life was illustrious for his faith and his other virtues.
\switchcolumn*
\selectlanguage{latin}
Constantinópoli sancti 
 Theodóri, Abbátis Studítæ, qui, pro fide cathólica advérsus Iconoclástas 
 strénue pugnans, factus est apud univérsam Ecclésiam cathólicam célebris.
\switchcolumn
\selectlanguage{english}
At Constantinople, St. Theodore, 
 abbot of Studium, who fought valiantly for the Catholic faith against the 
 Iconoclasts, and became famed throughout the universal Church.
\switchcolumn*
\selectlanguage{latin}
In monastério Cryptæ 
 Ferrátæ, in agro Tusculáno, sancti Bartholomæi Abbátis, qui fuit sócius 
 beáti Nili, ejúsque vitam conscrípsit.
\switchcolumn
\selectlanguage{english}
In the monastery of Grottaferrata, 
 in the Tuscan plain, the holy abbot Bartholomew, a companion of blessed 
 Nilus, whose life he wrote.
\switchcolumn*
\selectlanguage{latin}
In província Sámnii 
 beáti Mennæ solitárii, cujus virtútes et mirácula sanctus Gregórius Papa 
 commémorat.
\switchcolumn
\selectlanguage{english}
In the province of Abruzzi, blessed 
 Mennas, a solitary whose virtues and miracles are mentioned by Pope St. 
 Gregory.
\switchcolumn*
\selectlanguage{latin}
\end{paracol}


% ---- martyrology/mart11/mart1112.htm
\needspace{10\baselineskip}
\begin{paracol}{2}
\selectlanguage{latin}
\begin{center}{\color{gregoriocolor} Prídie Idus Novémbris. 
 Luna\dots\ }\end{center}
\switchcolumn
\selectlanguage{english}
\begin{center}{\color{gregoriocolor} The 
 12\textsuperscript{th} Day of 
 November. The\dots\ Day of the Moon.}\end{center}
\end{paracol}

\noindent\begin{tabularx}{\linewidth}{*{19}{>{\centering\arraybackslash}X}}
 \textcolor{gregoriocolor}{a} & \textcolor{gregoriocolor}{b} & \textcolor{gregoriocolor}{c} & \textcolor{gregoriocolor}{d} & \textcolor{gregoriocolor}{e} & \textcolor{gregoriocolor}{f} & \textcolor{gregoriocolor}{g} & \textcolor{gregoriocolor}{h} & \textcolor{gregoriocolor}{i} & \textcolor{gregoriocolor}{k} & \textcolor{gregoriocolor}{l} & \textcolor{gregoriocolor}{m} & \textcolor{gregoriocolor}{n} & \textcolor{gregoriocolor}{p} & \textcolor{gregoriocolor}{q} & \textcolor{gregoriocolor}{r} & \textcolor{gregoriocolor}{s} & \textcolor{gregoriocolor}{t} & \textcolor{gregoriocolor}{u} \\
 22 & 23 & 24 & 25 & 26 & 27 & 28 & 29 & 30 & 1 & 2 & 3 & 4 & 5 & 6 & 7 & 8 & 9 & 10 \\
\end{tabularx}
\vspace{0.5\baselineskip}
\noindent\begin{tabularx}{\linewidth}{*{12}{>{\centering\arraybackslash}X}}
 \textcolor{gregoriocolor}{A} & \textcolor{gregoriocolor}{B} & \textcolor{gregoriocolor}{C} & \textcolor{gregoriocolor}{D} & \textcolor{gregoriocolor}{E} & F & \textcolor{gregoriocolor}{F} & \textcolor{gregoriocolor}{G} & \textcolor{gregoriocolor}{H} & \textcolor{gregoriocolor}{M} & \textcolor{gregoriocolor}{N} & \textcolor{gregoriocolor}{P} \\
 11 & 12 & 13 & 14 & 15 & 16 & 16 & 17 & 18 & 19 & 20 & 21 \\
\end{tabularx}

\begin{paracol}{2}
\selectlanguage{latin}
\lettrine[lines=2]{S}{ancti} Martíni Primi, 
 Papæ et Mártyris, cujus dies natális sextodécimo Kaléndas Octóbris 
 recensétur.
\switchcolumn
\selectlanguage{english}
\lettrine[lines=2]{T}{he} Feast of St. Martin I, pope and 
 martyr, whose birthday is mentioned on the 16th day of September.
\switchcolumn*
\selectlanguage{latin}
Vitépsci, in Polónia, 
 pássio sancti Jósaphat, e sancti Basilíi Ordine, Epíscopi Polocénsis et 
 Mártyris; qui a schismáticis, in ódium cathólicæ unitátis et veritátis, 
 crudéliter interféctus est, et a Pio Papa Nono inter sanctos Mártyres 
 adscríptus. Ejus tamen festívitas recólitur décimo octávo Kaléndas 
 Decémbris.
\switchcolumn
\selectlanguage{english}
At Witebsk in Poland, the martyrdom 
 of St. Josaphat, of the Order of St. Basil, a Polish archbishop and martyr, 
 who was cruelly slain by schismatics through hatred of Catholic unity and 
 truth. He was canonized by Pope Pius IX, and his feast is observed on 
 the 16th of November.
\switchcolumn*
\selectlanguage{latin}
Complúti, in Hispánia, 
 natális sancti Dídaci Confessóris, ex Ordine Minórum, humilitáte célebris; 
 quem Xystus Quintus, Póntifex Máximus, Sanctórum catálogo adscrípsit. 
 Ipsíus autem festum sequénti die celebrátur.
\switchcolumn
\selectlanguage{english}
At Alcala in Spain, the birthday of 
 St. Didacus, confessor, a member of the Order of Friars Minor well known for 
 his humility. Pope Sixtus V included him in the catalogue of the 
 saints and his feast is celebrated on the day following.
\switchcolumn*
\selectlanguage{latin}
In Asia pássio 
 sanctórum Aurélii et Públii Episcopórum.
\switchcolumn
\selectlanguage{english}
In Asia, the martyrdom of the holy 
 bishops Aurelius and Publius.
\switchcolumn*
\selectlanguage{latin}
Eschæ, in Bélgio, 
 sancti Livíni, Epíscopi et Mártyris; qui, cum plúrimos ad Christi fidem 
 convertísset, a Pagánis necátus est. Ipsíus vero corpus ad Portum 
 Gandæ póstea translátum fuit.
\switchcolumn
\selectlanguage{english}
At Eschen in Belgium, St. Livinus, 
 bishop and martyr. After converting many to the faith he was slain by 
 heathens. His body, however, was afterwards translated to Ghent.
\switchcolumn*
\selectlanguage{latin}
Apud Casimíriam, in 
 Polónia, sanctórum Mártyrum Eremitárum Benedícti, Joánnis, Matthæi, Isaac et 
 Christiáni; qui a prædónibus, divíno inténti servítio, dire vexáti sunt et 
 gládiis occísi.
\switchcolumn
\selectlanguage{english}
At Gnesen in Poland, the holy 
 hermits and martyrs Benedict, John, Matthew, Isaac, and Christian. 
 They were savagely attacked by robbers and slain by the sword while there 
 were at prayer.
\switchcolumn*
\selectlanguage{latin}
Apud óppidum Sergíniam, 
 in território Senonénsi, sancti Patérni, Mónachi et Mártyris; qui, dum eídem 
 occurréntes in ipsíus óppidi silva latrónes ad emendándam vitam incitáret, 
 ab illis trucidátus est.
\switchcolumn
\selectlanguage{english}
In the neighbourhood of Sens, St. 
 Paternus, monk and martyr. He had met some robbers in a nearby forest, 
 and for attempting to persuade them to amend their lives, they slew him.
\switchcolumn*
\selectlanguage{latin}
Avenióne sancti Rufi, 
 qui éxstitit primus ejúsdem civitátis Epíscopus.
\switchcolumn
\selectlanguage{english}
At Avignon, St. Rufus, the first 
 bishop of that city.
\switchcolumn*
\selectlanguage{latin}
Colóniæ Agrippínæ 
 deposítio sancti Cunibérti Epíscopi.
\switchcolumn
\selectlanguage{english}
At Cologne, the death of St. 
 Cunibert, bishop.
\switchcolumn*
\selectlanguage{latin}
Turiasóne, in Hispánia 
 Tarraconénsi, beáti Æmiliáni Presbyteri, qui innúmeris miráculis cláruit; 
 cujus admirábilem vitam sanctus Bráulio, Cæsaraugustánus Epíscopus, 
 descrípsit.
\switchcolumn
\selectlanguage{english}
At Tarazona in Aragon, blessed 
 Emilian, a priest favoured with many miracles. His admirable life was 
 recorded by St. Braulio, bishop of Saragossa.
\switchcolumn*
\selectlanguage{latin}
Constantinópoli sancti 
 Nili Abbátis, qui, sub Theodósio junióre, ex Præfécto ejúsdem civitátis 
 factus Mónachus, doctrína et sanctitáte cláruit.
\switchcolumn
\selectlanguage{english}
At Constantinople, St. Nilus, abbot, 
 who resigned as prefect of the city to become a monk, and during the reign 
 of Theodosius the Younger became distinguished for his learning and 
 sanctity.
\switchcolumn*
\selectlanguage{latin}
\end{paracol}


% ---- martyrology/mart11/mart1113.htm
\needspace{10\baselineskip}
\begin{paracol}{2}
\selectlanguage{latin}
\begin{center}{\color{gregoriocolor} Idibus Novémbris. 
 Luna\dots\ }\end{center}
\switchcolumn
\selectlanguage{english}
\begin{center}{\color{gregoriocolor} The 
 13\textsuperscript{th} Day of 
 November. The\dots\ Day of the Moon.}\end{center}
\end{paracol}

\noindent\begin{tabularx}{\linewidth}{*{19}{>{\centering\arraybackslash}X}}
 \textcolor{gregoriocolor}{a} & \textcolor{gregoriocolor}{b} & \textcolor{gregoriocolor}{c} & \textcolor{gregoriocolor}{d} & \textcolor{gregoriocolor}{e} & \textcolor{gregoriocolor}{f} & \textcolor{gregoriocolor}{g} & \textcolor{gregoriocolor}{h} & \textcolor{gregoriocolor}{i} & \textcolor{gregoriocolor}{k} & \textcolor{gregoriocolor}{l} & \textcolor{gregoriocolor}{m} & \textcolor{gregoriocolor}{n} & \textcolor{gregoriocolor}{p} & \textcolor{gregoriocolor}{q} & \textcolor{gregoriocolor}{r} & \textcolor{gregoriocolor}{s} & \textcolor{gregoriocolor}{t} & \textcolor{gregoriocolor}{u} \\
 23 & 24 & 25 & 26 & 27 & 28 & 29 & 30 & 1 & 2 & 3 & 4 & 5 & 6 & 7 & 8 & 9 & 10 & 11 \\
\end{tabularx}
\vspace{0.5\baselineskip}
\noindent\begin{tabularx}{\linewidth}{*{12}{>{\centering\arraybackslash}X}}
 \textcolor{gregoriocolor}{A} & \textcolor{gregoriocolor}{B} & \textcolor{gregoriocolor}{C} & \textcolor{gregoriocolor}{D} & \textcolor{gregoriocolor}{E} & F & \textcolor{gregoriocolor}{F} & \textcolor{gregoriocolor}{G} & \textcolor{gregoriocolor}{H} & \textcolor{gregoriocolor}{M} & \textcolor{gregoriocolor}{N} & \textcolor{gregoriocolor}{P} \\
 12 & 13 & 14 & 15 & 16 & 17 & 17 & 18 & 19 & 20 & 21 & 22 \\
\end{tabularx}

\begin{paracol}{2}
\selectlanguage{latin}
\lettrine[lines=2]{S}{ancti} Dídaci, ex 
 Ordine Minórum, Confessóris; cujus dies natális recólitur prídie hujus diéi.
\switchcolumn
\selectlanguage{english}
\lettrine[lines=2]{S}{t.} Didacus, confessor of the Order 
 of Friars Minor, whose birthday occurred on the preceding day.
\switchcolumn*
\selectlanguage{latin}
Ravénnæ natális 
 sanctórum Mártyrum Valentíni, Solutóris et Victóris; qui sub Diocletiáno 
 Imperatóre passi sunt.
\switchcolumn
\selectlanguage{english}
At Ravenna, the birthday of the holy 
 martyrs Valentine, Salutor, and Victor, who suffered under Emperor 
 Diocletian.
\switchcolumn*
\selectlanguage{latin}
Aquis, in província 
 Narbonénsi, beáti Mítrii, claríssimi Mártyris.
\switchcolumn
\selectlanguage{english}
At Aix, in the province of Narbonne, 
 the renowned martyr, blessed Mitrius.
\switchcolumn*
\selectlanguage{latin}
Cæsaréæ, in Palæstína, 
 pássio sanctórum Antoníni, Zébinæ, Germáni et Ennathæ Vírginis. Hæc, 
 sub Galério Maximiáno Imperatóre, verbéribus cæsa, igne cremáta est; illi 
 vero, cum intrépidi ac líbera voce Firmiliánum Præsidem, diis immolántem, 
 impietátis argúerent, cápite cæsi sunt.
\switchcolumn
\selectlanguage{english}
At Caesarea in Palestine, the 
 martyrdom of the Saints Antoninus, Zebina, Germanus, and the virgin Ennatha. 
 Under Galerius Maximian, Ennatha was scourged and burned alive, while the 
 others, for boldly reproaching the governor Firmilian for his idolatry in 
 sacrificing to the gods, were beheaded.
\switchcolumn*
\selectlanguage{latin}
In Africa sanctórum 
 Mártyrum Hispanórum Arcádii, Paschásii, Probi et Eutychiáni; qui, in 
 persecutióne Wandálica, cum in Ariánam perfídiam nullátenus declináre 
 pateréntur, hinc a Genseríco, Rege Ariáno, primum proscrípti, deínde acti in 
 exsílium atque atrocíssimis supplíciis cruciáti, postrémum divérso mortis 
 génere interémpti sunt. Tunc et Paulílli puéruli, germáni sanctórum 
 Paschásii et Eutychiáni, constántia enítuit; qui, cum de fide cathólica 
 nullátenus posset avélli, fústibus diu cæsus est, atque ad ínfimam 
 servitútem damnátus.
\switchcolumn
\selectlanguage{english}
In Africa, the holy martyrs Arcadius, 
 Paschasius, Probus, and Eutychian, Spaniards who absolutely refused to yield 
 to the Arian perfidy, during the persecution of the Vandals. 
 Accordingly, they were condemned by the Arian king Genseric, driven into exile, and finally, after being subjected to fearful tortures, were put to 
 death in divers manners. At that time there was also seen the 
 constancy of the small boy Paulillus, brother of the Saints Paschasius and 
 Eutychian. Because he could not be turned from the Catholic faith, he 
 was long beaten with rods and sentenced to the lowest servitude.
\switchcolumn*
\selectlanguage{latin}
Romæ sancti Nicolái Papæ Primi, vigóre apostólico præstántis.
\switchcolumn
\selectlanguage{english}
At Rome, Pope St. Nicholas, 
 distinguished for the apostolic spirit.
\switchcolumn*
\selectlanguage{latin}
Turónis, in Gállia, 
 sancti Brítii Epíscopi, qui fuit discípulus beáti Martíni Epíscopi.
\switchcolumn
\selectlanguage{english}
At Tours in France, St. Brice, 
 bishop, a disciple of the blessed Bishop Martin.
\switchcolumn*
\selectlanguage{latin}
Toléti, in Hispánia, 
 sancti Eugénii Epíscopi.
\switchcolumn
\selectlanguage{english}
At Toledo in Spain, St. Eugene, 
 bishop.
\switchcolumn*
\selectlanguage{latin}
Arvérnis, in Gállia, 
 sancti Quinctiáni Epíscopi.
\switchcolumn
\selectlanguage{english}
In Auvergne in France, St. Quinctian, 
 bishop.
\switchcolumn*
\selectlanguage{latin}
Cremónæ, in Insúbria, 
 sancti Homobóni Confessóris; quem, miráculis clarum, Innocéntius Papa 
 Tértius in Sanctórum númerum rétulit.
\switchcolumn
\selectlanguage{english}
At Cremona, in the duchy of Milan, 
 St. Homobonus, confessor, renowned for miracles, whom Innocent III placed 
 among the saints.
\switchcolumn*
\selectlanguage{latin}
\end{paracol}


% ---- martyrology/mart11/mart1114.htm
\needspace{10\baselineskip}
\begin{paracol}{2}
\selectlanguage{latin}
\begin{center}{\color{gregoriocolor} Décimo octávo Kaléndas Decémbris. 
 Luna\dots\ }\end{center}
\switchcolumn
\selectlanguage{english}
\begin{center}{\color{gregoriocolor} The 
 14\textsuperscript{th} Day of 
 November. The\dots\ Day of the Moon.}\end{center}
\end{paracol}

\noindent\begin{tabularx}{\linewidth}{*{19}{>{\centering\arraybackslash}X}}
 \textcolor{gregoriocolor}{a} & \textcolor{gregoriocolor}{b} & \textcolor{gregoriocolor}{c} & \textcolor{gregoriocolor}{d} & \textcolor{gregoriocolor}{e} & \textcolor{gregoriocolor}{f} & \textcolor{gregoriocolor}{g} & \textcolor{gregoriocolor}{h} & \textcolor{gregoriocolor}{i} & \textcolor{gregoriocolor}{k} & \textcolor{gregoriocolor}{l} & \textcolor{gregoriocolor}{m} & \textcolor{gregoriocolor}{n} & \textcolor{gregoriocolor}{p} & \textcolor{gregoriocolor}{q} & \textcolor{gregoriocolor}{r} & \textcolor{gregoriocolor}{s} & \textcolor{gregoriocolor}{t} & \textcolor{gregoriocolor}{u} \\
 24 & 25 & 26 & 27 & 28 & 29 & 30 & 1 & 2 & 3 & 4 & 5 & 6 & 7 & 8 & 9 & 10 & 11 & 12 \\
\end{tabularx}
\vspace{0.5\baselineskip}
\noindent\begin{tabularx}{\linewidth}{*{12}{>{\centering\arraybackslash}X}}
 \textcolor{gregoriocolor}{A} & \textcolor{gregoriocolor}{B} & \textcolor{gregoriocolor}{C} & \textcolor{gregoriocolor}{D} & \textcolor{gregoriocolor}{E} & F & \textcolor{gregoriocolor}{F} & \textcolor{gregoriocolor}{G} & \textcolor{gregoriocolor}{H} & \textcolor{gregoriocolor}{M} & \textcolor{gregoriocolor}{N} & \textcolor{gregoriocolor}{P} \\
 13 & 14 & 15 & 16 & 17 & 18 & 18 & 19 & 20 & 21 & 22 & 23 \\
\end{tabularx}

\begin{paracol}{2}
\selectlanguage{latin}
\lettrine[lines=2]{S}{ancti} Jósaphat, e 
 sancti Basilíi Ordine, Epíscopi Polocénsis et Mártyris, cujus dies natális 
 recensétur prídie Idus Novémbris.
\switchcolumn
\selectlanguage{english}
\lettrine[lines=2]{S}{t.} Josaphat, of the Order of St. 
 Basil, archbishop and martyr of Poland, whose birthday was observed on the 
 12th of November.
\switchcolumn*
\selectlanguage{latin}
Gangris, in Paphlagónia, 
 sancti Hypátii Epíscopi, qui, a magna Nicæna Synodo rédiens, a Novatiánis 
 hæréticis in via lapídibus impetítus, Martyr occúbuit.
\switchcolumn
\selectlanguage{english}
At Gangra in Paphlagonia, St. 
 Hypatius, bishop, who on his way home from the great Council of Nicaea, was 
 attacked with stones by the Novatian heretics, and died a martyr.
\switchcolumn*
\selectlanguage{latin}
Heracléæ, in Thrácia, 
 natális sanctórum Mártyrum Clementíni, Theodóti et Philómeni.
\switchcolumn
\selectlanguage{english}
At Heraclea in Thrace, the birthday 
 of the holy martyrs Clementinus, Theodotus and Philomenus.
\switchcolumn*
\selectlanguage{latin}
Alexandríæ sancti 
 Serapiónis Mártyris, quem persecutóres, sub Décio Príncipe, ita 
 crudelíssimis affecérunt supplíciis, ut cunctas et junctúras membrórum prius 
 sólverent, eúmque póstea de superióribus domus suæ præcipitárent; atque ille 
 sic gloriósus Christi Martyr efficerétur.
\switchcolumn
\selectlanguage{english}
At Alexandria, St. Serapion, martyr, 
 whom the persecutors under Emperor Decius subjected to torments so cruel 
 that all his limbs were disjointed. He became a martyr of Christ by 
 being hurled from the upper part of the house.
\switchcolumn*
\selectlanguage{latin}
Trecis, in Gállia, 
 sancti Venerándi Mártyris, sub Aureliáno Imperatóre.
\switchcolumn
\selectlanguage{english}
At Troyes in France, under Emperor 
 Aurelian, St. Venerandus, martyr.
\switchcolumn*
\selectlanguage{latin}
In Gállia sanctæ 
 Venerándæ Vírginis, quæ, sub Antoníno Imperatóre et Asclepíade Præside, 
 martyrii corónam accépit.
\switchcolumn
\selectlanguage{english}
Also in France, the holy virgin 
 Veneranda, who received the crown of martyrdom under Emperor Antoninus and 
 the governor Asclepiades.
\switchcolumn*
\selectlanguage{latin}
Eméssæ, in Phœnícia, 
 pássio plurimárum sanctárum mulíerum, quæ, sub sævíssimo Arabum duce Mady, 
 ob Christi fidem, crudelíssime tortæ atque necátæ sunt.
\switchcolumn
\selectlanguage{english}
At Emesa in Phoenicia, the martyrdom 
 of many holy women, who were barbarously tortured and massacred for the 
 faith of Christ under Mady, a savage Arabian chief.
\switchcolumn*
\selectlanguage{latin}
Bonóniæ sancti Jucúndi, 
 Epíscopi et Confessóris.
\switchcolumn
\selectlanguage{english}
At Bologna, St. Jucundus, bishop and 
 confessor.
\switchcolumn*
\selectlanguage{latin}
Augæ, in Gállia, 
 tránsitus sancti Lauréntii, Epíscopi Dublinénsis.
\switchcolumn
\selectlanguage{english}
At Eu in France, St. Laurence, 
 bishop of Dublin.
\switchcolumn*
\selectlanguage{latin}
Algáriæ, in Africa, 
 beáti Serapiónis, qui, primus ex Ordine beátæ Maríæ de Mercéde redemptiónis 
 captivórum, pro captívis fidélibus rediméndis et Christiánæ fídei 
 prædicatióne actus in crucem et membrátim disséctus, martyrii palmam méruit 
 obtinére.
\switchcolumn
\selectlanguage{english}
At Algiers in Africa, blessed 
 Serapion, of the Order of Our Blessed Lady of Ransom. For the 
 redemption of the faithful in captivity and the preaching of the Christian 
 faith, he was the first of his Order to merit the palm of martyrdom, being 
 crucified and torn limb from limb.
\switchcolumn*
\selectlanguage{latin}
\end{paracol}


% ---- martyrology/mart11/mart1115.htm
\needspace{10\baselineskip}
\begin{paracol}{2}
\selectlanguage{latin}
\begin{center}{\color{gregoriocolor} Décimo séptimo Kaléndas Decémbris. 
 Luna\dots\ }\end{center}
\switchcolumn
\selectlanguage{english}
\begin{center}{\color{gregoriocolor} The 
 15\textsuperscript{th} Day of 
 November. The\dots\ Day of the Moon.}\end{center}
\end{paracol}

\noindent\begin{tabularx}{\linewidth}{*{19}{>{\centering\arraybackslash}X}}
 \textcolor{gregoriocolor}{a} & \textcolor{gregoriocolor}{b} & \textcolor{gregoriocolor}{c} & \textcolor{gregoriocolor}{d} & \textcolor{gregoriocolor}{e} & \textcolor{gregoriocolor}{f} & \textcolor{gregoriocolor}{g} & \textcolor{gregoriocolor}{h} & \textcolor{gregoriocolor}{i} & \textcolor{gregoriocolor}{k} & \textcolor{gregoriocolor}{l} & \textcolor{gregoriocolor}{m} & \textcolor{gregoriocolor}{n} & \textcolor{gregoriocolor}{p} & \textcolor{gregoriocolor}{q} & \textcolor{gregoriocolor}{r} & \textcolor{gregoriocolor}{s} & \textcolor{gregoriocolor}{t} & \textcolor{gregoriocolor}{u} \\
 25 & 26 & 27 & 28 & 29 & 30 & 1 & 2 & 3 & 4 & 5 & 6 & 7 & 8 & 9 & 10 & 11 & 12 & 13 \\
\end{tabularx}
\vspace{0.5\baselineskip}
\noindent\begin{tabularx}{\linewidth}{*{12}{>{\centering\arraybackslash}X}}
 \textcolor{gregoriocolor}{A} & \textcolor{gregoriocolor}{B} & \textcolor{gregoriocolor}{C} & \textcolor{gregoriocolor}{D} & \textcolor{gregoriocolor}{E} & F & \textcolor{gregoriocolor}{F} & \textcolor{gregoriocolor}{G} & \textcolor{gregoriocolor}{H} & \textcolor{gregoriocolor}{M} & \textcolor{gregoriocolor}{N} & \textcolor{gregoriocolor}{P} \\
 14 & 15 & 16 & 17 & 18 & 19 & 19 & 20 & 21 & 22 & 23 & 24 \\
\end{tabularx}

\begin{paracol}{2}
\selectlanguage{latin}
\lettrine[lines=2]{C}{olóniæ} Agrippínæ 
 sancti Albérti Epíscopi et Confessóris, ex Ordine Prædicatórum, cognoménto 
 Magni, sanctitáte et doctrína célebris, quem Pius Papa Undécimus Doctórem 
 universális Ecclésiæ declarávit, et Pius Duodécimus cultórum scientiárum 
 naturálium cæléstem apud Deum Patrónum constítuit.
\switchcolumn
\selectlanguage{english}
\lettrine[lines=2]{A}{t} Cologne, St. Albert, surnamed the 
 Great, bishop and confessor of the Order of Preachers, renowned for his 
 holiness and learning. Pope Pius XI appointed him as Doctor of the 
 universal Church, and Pius XII appointed him as heavenly patron of those 
 studying the natural sciences.
\switchcolumn*
\selectlanguage{latin}
Eódem die natális 
 sancti Eugénii, Epíscopi Toletáni et Mártyris; qui fuit beáti Dionysii 
 Areopagítæ discípulus, et in território Parisiénsi, consummáto martyrii 
 cursu, beátæ passiónis corónam percépit a Dómino. Ipsíus autem corpus 
 Tolétum, in Hispánia, póstea fuit translátum.
\switchcolumn
\selectlanguage{english}
Also, the birthday of St. Eugene, 
 bishop of Toledo and martyr, disciple of blessed Denis the Areopagite. 
 His martyrdom was completed near Paris, and he received from our Lord a 
 crown for his blessed sufferings. His body was afterwards translated 
 to Toledo in Spain.
\switchcolumn*
\selectlanguage{latin}
Nolæ, in Campánia, beáti Felícis, Epíscopi et Mártyris; qui, a quintodécimo ætátis suæ anno, 
 miráculis cláruit, et, sub Marciáno Præside, cum áliis trigínta Sóciis, 
 agónem martyrii complévit.
\switchcolumn
\selectlanguage{english}
At Nola in Campania, blessed Felix, 
 bishop and martyr, who was renowned for miracles from his fifteenth year. 
 He completed the combats of his martyrdom with thirty others, under the 
 governor Marcian.
\switchcolumn*
\selectlanguage{latin}
Edéssæ, in Mesopotámia, pássio sancti Abibi Diáconi, qui sub Licínio Imperatóre et Lysánia Præside, 
 únguibus lacerátus, in ignem conjéctus est.
\switchcolumn
\selectlanguage{english}
At Edessa in Mesopotamia, the 
 martyrdom of St. Abibus, deacon, who was torn with iron hooks and cast into 
 the fire in the time of Emperor Licinius and the governor Lysanias.
\switchcolumn*
\selectlanguage{latin}
Ibídem sanctórum 
 Mártyrum Guríæ et Samónæ, sub Diocletiáno Imperatóre et Antoníno Præside.
\switchcolumn
\selectlanguage{english}
In the same place, the holy martyrs 
 Gurias and Samonas, under Emperor Diocletian and the governor Antoninus.
\switchcolumn*
\selectlanguage{latin}
In Africa sanctórum 
 Mártyrum Secúndi, Fidentiáni et Várici.
\switchcolumn
\selectlanguage{english}
In Africa, the holy martyrs Secundus, 
 Fidentian, and Varicus.
\switchcolumn*
\selectlanguage{latin}
Apud Arcum, in 
 território Santonénsi, natális sancti Machúti, Aleténsis in Gállia Epíscopi; 
 qui, in Anglia natus, a primævo ætátis suæ tirocínio miráculis emícuit.
\switchcolumn
\selectlanguage{english}
At Archingeay, in the neighbourhood 
 of Saintes, the birthday of St. Malo, bishop of Aleth, in France. He 
 was born in England and from his earliest years was famed for his miracles.
\switchcolumn*
\selectlanguage{latin}
Verónæ sancti Lupérii, 
 Epíscopi et Confessóris.
\switchcolumn
\selectlanguage{english}
At Verona, St. Luperius, bishop and 
 confessor.
\switchcolumn*
\selectlanguage{latin}
Kahlembérgæ, prope 
 Vindobónam, in Austria, sancti Leopóldi, ejúsdem provínciæ Austriæ 
 Marchiónis, quem Innocéntius Papa Octávus in Sanctórum númerum adscrípsit.
\switchcolumn
\selectlanguage{english}
At Klosterneuburg, near Vienna in 
 Austria, St. Leopold, margrave of that province of Austria. He was 
 placed on the canon of the saints by Pope Innocent VIII.
\switchcolumn*
\selectlanguage{latin}
\end{paracol}


% ---- martyrology/mart11/mart1116.htm
\needspace{10\baselineskip}
\begin{paracol}{2}
\selectlanguage{latin}
\begin{center}{\color{gregoriocolor} Sextodécimo Kaléndas Decémbris. 
 Luna\dots\ }\end{center}
\switchcolumn
\selectlanguage{english}
\begin{center}{\color{gregoriocolor} The 
 16\textsuperscript{th} Day of 
 November. The\dots\ Day of the Moon.}\end{center}
\end{paracol}

\noindent\begin{tabularx}{\linewidth}{*{19}{>{\centering\arraybackslash}X}}
 \textcolor{gregoriocolor}{a} & \textcolor{gregoriocolor}{b} & \textcolor{gregoriocolor}{c} & \textcolor{gregoriocolor}{d} & \textcolor{gregoriocolor}{e} & \textcolor{gregoriocolor}{f} & \textcolor{gregoriocolor}{g} & \textcolor{gregoriocolor}{h} & \textcolor{gregoriocolor}{i} & \textcolor{gregoriocolor}{k} & \textcolor{gregoriocolor}{l} & \textcolor{gregoriocolor}{m} & \textcolor{gregoriocolor}{n} & \textcolor{gregoriocolor}{p} & \textcolor{gregoriocolor}{q} & \textcolor{gregoriocolor}{r} & \textcolor{gregoriocolor}{s} & \textcolor{gregoriocolor}{t} & \textcolor{gregoriocolor}{u} \\
 26 & 27 & 28 & 29 & 30 & 1 & 2 & 3 & 4 & 5 & 6 & 7 & 8 & 9 & 10 & 11 & 12 & 13 & 14 \\
\end{tabularx}
\vspace{0.5\baselineskip}
\noindent\begin{tabularx}{\linewidth}{*{12}{>{\centering\arraybackslash}X}}
 \textcolor{gregoriocolor}{A} & \textcolor{gregoriocolor}{B} & \textcolor{gregoriocolor}{C} & \textcolor{gregoriocolor}{D} & \textcolor{gregoriocolor}{E} & F & \textcolor{gregoriocolor}{F} & \textcolor{gregoriocolor}{G} & \textcolor{gregoriocolor}{H} & \textcolor{gregoriocolor}{M} & \textcolor{gregoriocolor}{N} & \textcolor{gregoriocolor}{P} \\
 15 & 16 & 17 & 18 & 19 & 20 & 20 & 21 & 22 & 23 & 24 & 25 \\
\end{tabularx}

\begin{paracol}{2}
\selectlanguage{latin}
\lettrine[lines=2]{S}{anctæ} Gertrúdis 
 Vírginis, cujus natális sequénti die recensétur.
\switchcolumn
\selectlanguage{english}
\lettrine[lines=2]{S}{t.} Gertrude, virgin, whose birthday 
 is on the 17th of November.
\switchcolumn*
\selectlanguage{latin}
Edimbúrgi, in Scótia, sanctæ Margarítæ Víduæ, Scotórum Regínæ, amóre in páuperes et voluntária 
 paupertáte célebris. Ipsíus tamen festívitas quarto Idus Júnii 
 celebrátur.
\switchcolumn
\selectlanguage{english}
At Edinburgh in Scotland, the 
 birthday of St. Margaret, queen of the Scots and widow, renowned for her 
 love of the poor and her voluntary poverty. Her feast is celebrated on 
 the 10th of June.
\switchcolumn*
\selectlanguage{latin}
In Africa sanctórum 
 Mártyrum Rufíni, Marci, Valérii et Sociórum.
\switchcolumn
\selectlanguage{english}
In Africa, the holy martyrs Rufinus, 
 Mark, Valerius, and their fellows.
\switchcolumn*
\selectlanguage{latin}
Eódem die sanctórum 
 Mártyrum Elpídii, Marcélli, Eustóchii et Sociórum; ex quibus Elpídius, cum 
 esset órdinis Senatórii et coram Juliáno Apóstata Christiánam fidem 
 constantíssime profiterétur, ídeo, primum equis indómitis, una cum Sóciis, 
 alligátus atque pertráctus, deínde, in ignem conjéctus, gloriósum martyrium 
 consummávit.
\switchcolumn
\selectlanguage{english}
The same day, the holy martyrs 
 Elpidius, Marcellus, Eustochius, and their companions. Elpidius, who 
 was a senator, perseveringly confessed the Christian faith before Julian the 
 Apostate, and, with his companions, was tied to wild horses and dragged by 
 them, thus fulfilling a glorious martyrdom.
\switchcolumn*
\selectlanguage{latin}
Lugdúni, in Gállia, 
 natális sancti Euchérii, Epíscopi et Confessóris, viri admirándæ fídei et 
 doctrínæ. Hic, ex nobilíssimo Senatórum órdine ad religiósam vitam 
 habitúmque convérsus, diu, intra septa spelúncæ sponte conclúsus, in 
 oratiónibus et jejúniis Christo servívit; deínde apud præfátam urbem in 
 pontificáli Cáthedra, revelánte Angelo, solémniter collocátus est.
\switchcolumn
\selectlanguage{english}
At Lyons in France, the birthday of 
 St. Eucherius, bishop and confessor, a man of extraordinary faith and 
 learning. He renounced the senatorial dignity to embrace the religious 
 life, and for a long time voluntarily shut himself up in a cave, where he 
 served Christ in prayer and fasting. Afterwards, through the 
 revelation of an angel, he was solemnly installed in the episcopal chair of 
 the city of Lyons.
\switchcolumn*
\selectlanguage{latin}
Patávii sancti Fidéntii 
 Epíscopi.
\switchcolumn
\selectlanguage{english}
At Padua, St. Fidentius, bishop.
\switchcolumn*
\selectlanguage{latin}
Cantuáriæ, in Anglia, 
 sancti Edmúndi, Epíscopi et Confessóris; qui, pro Ecclésiæ suæ júribus 
 tuéndis in exsílium actus, apud Provínum, Sénonum óppidum, sanctíssime óbiit; 
 et Sanctórum cánoni ab Innocéntio Papa Quarto adscríptus est.
\switchcolumn
\selectlanguage{english}
At Canterbury in England, St. 
 Edmund, archbishop and confessor, who was sent into exile for having 
 maintained the rights of his church. He died a most holy death at 
 Provins, a town near Sens, and was canonized by Innocent IV.
\switchcolumn*
\selectlanguage{latin}
Eódem die deposítio 
 sancti Othmári Abbátis.
\switchcolumn
\selectlanguage{english}
The same day, the death of St. 
 Othmar, abbot.
\switchcolumn*
\selectlanguage{latin}
\end{paracol}


% ---- martyrology/mart11/mart1117.htm
\needspace{10\baselineskip}
\begin{paracol}{2}
\selectlanguage{latin}
\begin{center}{\color{gregoriocolor} Quintodécimo Kaléndas Decémbris. 
 Luna\dots\ }\end{center}
\switchcolumn
\selectlanguage{english}
\begin{center}{\color{gregoriocolor} The 
 17\textsuperscript{th} Day of 
 November. The\dots\ Day of the Moon.}\end{center}
\end{paracol}

\noindent\begin{tabularx}{\linewidth}{*{19}{>{\centering\arraybackslash}X}}
 \textcolor{gregoriocolor}{a} & \textcolor{gregoriocolor}{b} & \textcolor{gregoriocolor}{c} & \textcolor{gregoriocolor}{d} & \textcolor{gregoriocolor}{e} & \textcolor{gregoriocolor}{f} & \textcolor{gregoriocolor}{g} & \textcolor{gregoriocolor}{h} & \textcolor{gregoriocolor}{i} & \textcolor{gregoriocolor}{k} & \textcolor{gregoriocolor}{l} & \textcolor{gregoriocolor}{m} & \textcolor{gregoriocolor}{n} & \textcolor{gregoriocolor}{p} & \textcolor{gregoriocolor}{q} & \textcolor{gregoriocolor}{r} & \textcolor{gregoriocolor}{s} & \textcolor{gregoriocolor}{t} & \textcolor{gregoriocolor}{u} \\
 27 & 28 & 29 & 30 & 1 & 2 & 3 & 4 & 5 & 6 & 7 & 8 & 9 & 10 & 11 & 12 & 13 & 14 & 15 \\
\end{tabularx}
\vspace{0.5\baselineskip}
\noindent\begin{tabularx}{\linewidth}{*{12}{>{\centering\arraybackslash}X}}
 \textcolor{gregoriocolor}{A} & \textcolor{gregoriocolor}{B} & \textcolor{gregoriocolor}{C} & \textcolor{gregoriocolor}{D} & \textcolor{gregoriocolor}{E} & F & \textcolor{gregoriocolor}{F} & \textcolor{gregoriocolor}{G} & \textcolor{gregoriocolor}{H} & \textcolor{gregoriocolor}{M} & \textcolor{gregoriocolor}{N} & \textcolor{gregoriocolor}{P} \\
 16 & 17 & 18 & 19 & 20 & 21 & 21 & 22 & 23 & 24 & 25 & 26 \\
\end{tabularx}

\begin{paracol}{2}
\selectlanguage{latin}
\lettrine[lines=2]{N}{eocæsaréæ,} in Ponto, 
 natális sancti Gregórii, Epíscopi et Confessóris, doctrína et sanctitáte 
 illústris, qui propter signa atque mirácula, quæ cum multa Ecclesiárum 
 glória perpetrávit, Thaumatúrgus est appellátus.
\switchcolumn
\selectlanguage{english}
\lettrine[lines=2]{A}{t} Neocaesarea in Pontus, the 
 birthday of St. Gregory, bishop and confessor, illustrious for his learning 
 and sanctity. The signs and miracles which he wrought to the great 
 glory of the Church gained for him the surname Wonderworker.
\switchcolumn*
\selectlanguage{latin}
Helpíthi, in Saxónia, 
 item natális sanctæ Gertrúdis Vírginis, ex Ordine sancti Benedícti, quæ dono 
 revelatiónum clara éxstitit. Ipsíus tamen festívitas prídie hujus diéi 
 celebrátur.
\switchcolumn
\selectlanguage{english}
At Hedelfs in Saxony, the birthday 
 of St. Gertrude, virgin of the Order of St. Benedict, who was famous for her 
 revelations. Her feast is observed on the preceding day.
\switchcolumn*
\selectlanguage{latin}
In Palæstína sanctórum 
 Mártyrum Alphæi et Zachæi, qui primo anno persecutiónis Diocletiáni, post 
 multa torménta, capitálem senténtiam subiére.
\switchcolumn
\selectlanguage{english}
In Palestine, in the first year of 
 Diocletian's persecution, the holy martyrs Alpheus and Zachaeus, who 
 underwent beheading after many tortures.
\switchcolumn*
\selectlanguage{latin}
Córdubæ, in Hispánia, 
 sanctórum Mártyrum Acíscli et Victóriæ germanórum, qui, in eádem 
 persecutióne, Diónis Præsidis jussu sævíssime cruciáti sunt, et illústri 
 passióne corónas a Dómino meruérunt.
\switchcolumn
\selectlanguage{english}
At Cordova in Spain, during the same 
 persecution, the holy martyrs Acisclus and his sister Victoria, who were 
 most cruelly tortured by order of the governor Dion, and thus merited to be 
 crowned by our Lord for their glorious sufferings.
\switchcolumn*
\selectlanguage{latin}
Alexandríæ sancti 
 Dionysii Epíscopi, summæ eruditiónis viri, qui, multis confessiónibus clarus 
 et pro passiónum tormentorúmque diversitáte magníficus, plenus diérum 
 Conféssor quiévit, Valeriáni et Galliéni Imperatórum tempóribus.
\switchcolumn
\selectlanguage{english}
At Alexandria, St. Denis, bishop, a 
 man of very great learning. In the time of Emperors Valerian and 
 Gallienus, renowned for often having confessed the faith, and illustrious 
 for the various sufferings and torments he had endured, full of days he 
 rested in peace a confessor.
\switchcolumn*
\selectlanguage{latin}
Aureliánis, in Gállia, 
 sancti Aniáni Epíscopi, cujus mortem in conspéctu Dómini pretiósam mirácula 
 crebra testántur.
\switchcolumn
\selectlanguage{english}
At Orleans in France, St. Anian, 
 bishop, the value of whose death in the sight of the Lord is attested by 
 frequent miracles.
\switchcolumn*
\selectlanguage{latin}
In Británnia sancti 
 Hugónis Epíscopi, qui, ex Mónacho Carthusiáno ad Ecclésiam Lincolniénsem 
 regéndam vocátus, multis cláruit miráculis, et sancto fine quiévit.
\switchcolumn
\selectlanguage{english}
In England, St. Hugh, bishop, who 
 was called to rule the church of Lincoln. He ended his holy life in 
 peace, renowned for many miracles.
\switchcolumn*
\selectlanguage{latin}
Turónis, in Gállia, 
 sancti Gregórii Epíscopi.
\switchcolumn
\selectlanguage{english}
At Tours in France, St. Gregory, 
 bishop.
\switchcolumn*
\selectlanguage{latin}
Floréntiæ sancti 
 Eugénii Confessóris, qui fuit Diáconus beáti Zenóbii, ejúsdem civitátis 
 Epíscopi.
\switchcolumn
\selectlanguage{english}
At Florence, St. Eugene, confessor, 
 the deacon of blessed Zenobius, bishop of that city.
\switchcolumn*
\selectlanguage{latin}
\end{paracol}


% ---- martyrology/mart11/mart1118.htm
\needspace{10\baselineskip}
\begin{paracol}{2}
\selectlanguage{latin}
\begin{center}{\color{gregoriocolor} Quartodécimo Kaléndas Decémbris. 
 Luna\dots\ }\end{center}
\switchcolumn
\selectlanguage{english}
\begin{center}{\color{gregoriocolor} The 
 18\textsuperscript{th} Day of 
 November. The\dots\ Day of the Moon.}\end{center}
\end{paracol}

\noindent\begin{tabularx}{\linewidth}{*{19}{>{\centering\arraybackslash}X}}
 \textcolor{gregoriocolor}{a} & \textcolor{gregoriocolor}{b} & \textcolor{gregoriocolor}{c} & \textcolor{gregoriocolor}{d} & \textcolor{gregoriocolor}{e} & \textcolor{gregoriocolor}{f} & \textcolor{gregoriocolor}{g} & \textcolor{gregoriocolor}{h} & \textcolor{gregoriocolor}{i} & \textcolor{gregoriocolor}{k} & \textcolor{gregoriocolor}{l} & \textcolor{gregoriocolor}{m} & \textcolor{gregoriocolor}{n} & \textcolor{gregoriocolor}{p} & \textcolor{gregoriocolor}{q} & \textcolor{gregoriocolor}{r} & \textcolor{gregoriocolor}{s} & \textcolor{gregoriocolor}{t} & \textcolor{gregoriocolor}{u} \\
 28 & 29 & 30 & 1 & 2 & 3 & 4 & 5 & 6 & 7 & 8 & 9 & 10 & 11 & 12 & 13 & 14 & 15 & 16 \\
\end{tabularx}
\vspace{0.5\baselineskip}
\noindent\begin{tabularx}{\linewidth}{*{12}{>{\centering\arraybackslash}X}}
 \textcolor{gregoriocolor}{A} & \textcolor{gregoriocolor}{B} & \textcolor{gregoriocolor}{C} & \textcolor{gregoriocolor}{D} & \textcolor{gregoriocolor}{E} & F & \textcolor{gregoriocolor}{F} & \textcolor{gregoriocolor}{G} & \textcolor{gregoriocolor}{H} & \textcolor{gregoriocolor}{M} & \textcolor{gregoriocolor}{N} & \textcolor{gregoriocolor}{P} \\
 17 & 18 & 19 & 20 & 21 & 22 & 22 & 23 & 24 & 25 & 26 & 27 \\
\end{tabularx}

\begin{paracol}{2}
\selectlanguage{latin}
\lettrine[lines=2]{R}{omæ} Dedicátio 
 Basilicárum sanctórum Petri et Pauli Apostolórum. Eárum primam, 
 restitútam in ampliórem formam, Summus Póntifex Urbánus Octávus consecrávit 
 hac ipsa recurrénte die; álteram vero, miserándo incéndio pénitus consúmptam, 
 ac magnificéntius reædificátam, Pius Nonus die décima Decémbris
 solémni ritu 
 consecrávit, ejúsque ánnuam commemoratiónem hodiérna die agéndam indíxit.
\switchcolumn
\selectlanguage{english}
\lettrine[lines=2]{A}{t} Rome, the dedication of the 
 basilicas of the holy apostles Peter and Paul. The former, having been 
 enlarged, was on this day solemnly consecrated by Urban VIII; while the 
 latter, more beautifully rebuilt after its total destruction by fire, was 
 solemnly dedicated on the 10th of December by Pius IX, though the feast in 
 commemoration of that event was transferred to this day.
\switchcolumn*
\selectlanguage{latin}
Antiochíæ natális 
 sancti Románi Mártyris, qui, témpore Galérii Imperatóris, cum Asclepíades 
 Præféctus in Ecclésiam irrúmperet eámque fúnditus conarétur evértere, 
 céteros Christiános hortátus est ut ei contradícerent, ideóque, post dira 
 torménta et abscissiónem linguæ (sine qua tamen Dei præcónia loquebátur), in cárcere strangulátus láqueo, célebri martyrio coronátur. Passus est 
 étiam ante ipsum puérulus, nómine Bárula, qui, cum fuísset ab eódem 
 interrogátus Præfécto utrum mélius esset unum Deum cólere an plures deos, 
 atque in unum Deum, quem Christiáni colunt, credéndum esse respondísset, 
 proptérea, verbéribus cæsus, jussus est decollári.
\switchcolumn
\selectlanguage{english}
At Antioch, the birthday of St. 
 Romanus, martyr, in the time of Emperor Galerius. When the prefect 
 Asclepiades attacked the Church and attempted to destroy it, Romanus 
 exhorted the Christians to resist him. After being subjected to severe 
 torments and the cutting out of his tongue (without which, however, he spake 
 the praises of God), he was strangled in prison and crowned with glorious 
 martyrdom. Before him suffered a young boy named Barulas, who being 
 asked by him whether it was better to worship one God or several gods, and 
 having answered that we must believe in the one God whom the Christians 
 adore, was scourged and beheaded.
\switchcolumn*
\selectlanguage{latin}
Item Antiochíæ sancti 
 Hesychii Mártyris, qui, cum esset miles, et præcéptum audísset ut quisquis 
 non sacrificáret idólis, cíngulum milítiæ depóneret, repénte cíngulum solvit; 
 ob quam causam, ingénti saxo in déxtera ejus ligáto, in flúvium præcipitári 
 jussus est.
\switchcolumn
\selectlanguage{english}
Also at Antioch, the holy martyr 
 Hesychius, a soldier. Hearing the order that anyone refusing to 
 sacrifice to idols should lay aside his military belt, he immediately took 
 off his. For this reason he was cast into the river with a large stone 
 tied to his right hand.
\switchcolumn*
\selectlanguage{latin}
Eódem die sanctórum 
 Orículi et Sociórum, qui, in persecutióne Wandálica, pro fide cathólica 
 passi sunt.
\switchcolumn
\selectlanguage{english}
On the same day, St. Oriculus and 
 his companions, who suffered for the Catholic faith in the Vandal 
 persecution.
\switchcolumn*
\selectlanguage{latin}
Mogúntiæ sancti Máximi 
 Epíscopi, qui, témpore Constántii multa passus ab Ariánis, Conféssor occúbuit.
\switchcolumn
\selectlanguage{english}
At Mainz, St. Maximus, bishop, who 
 suffered greatly at the hands of the Arians, and died a confessor in the 
 time of Constantius.
\switchcolumn*
\selectlanguage{latin}
Turónis, in Gállia, 
 tránsitus beáti Odónis, Abbátis Cluniacénsis.
\switchcolumn
\selectlanguage{english}
At Tours in France, the passing of 
 blessed Odo, abbot of Cluny.
\switchcolumn*
\selectlanguage{latin}
Antiochíæ sancti Thomæ 
 Mónachi, quem Antiochéni, ob sedátam ejus précibus pestem, solemnitáte ánnua 
 coluérunt.
\switchcolumn
\selectlanguage{english}
At Antioch, St. Thomas, a monk 
 honoured with an annual solemnity by the people of Antioch, for bringing the 
 end of a plague by his prayers.
\switchcolumn*
\selectlanguage{latin}
Lucæ, in Túscia, 
 Translátio sancti Frigdiáni, Epíscopi et Confessóris.
\switchcolumn
\selectlanguage{english}
At Lucca in Tuscany, the translation 
 of St. Frigidian, bishop and confessor.
\switchcolumn*
\selectlanguage{latin}
\end{paracol}


% ---- martyrology/mart11/mart1119.htm
\needspace{10\baselineskip}
\begin{paracol}{2}
\selectlanguage{latin}
\begin{center}{\color{gregoriocolor} Tertiodécimo Kaléndas Decémbris. 
 Luna\dots\ }\end{center}
\switchcolumn
\selectlanguage{english}
\begin{center}{\color{gregoriocolor} The 
 19\textsuperscript{th} Day of 
 November. The\dots\ Day of the Moon.}\end{center}
\end{paracol}

\noindent\begin{tabularx}{\linewidth}{*{19}{>{\centering\arraybackslash}X}}
 \textcolor{gregoriocolor}{a} & \textcolor{gregoriocolor}{b} & \textcolor{gregoriocolor}{c} & \textcolor{gregoriocolor}{d} & \textcolor{gregoriocolor}{e} & \textcolor{gregoriocolor}{f} & \textcolor{gregoriocolor}{g} & \textcolor{gregoriocolor}{h} & \textcolor{gregoriocolor}{i} & \textcolor{gregoriocolor}{k} & \textcolor{gregoriocolor}{l} & \textcolor{gregoriocolor}{m} & \textcolor{gregoriocolor}{n} & \textcolor{gregoriocolor}{p} & \textcolor{gregoriocolor}{q} & \textcolor{gregoriocolor}{r} & \textcolor{gregoriocolor}{s} & \textcolor{gregoriocolor}{t} & \textcolor{gregoriocolor}{u} \\
 29 & 30 & 1 & 2 & 3 & 4 & 5 & 6 & 7 & 8 & 9 & 10 & 11 & 12 & 13 & 14 & 15 & 16 & 17 \\
\end{tabularx}
\vspace{0.5\baselineskip}
\noindent\begin{tabularx}{\linewidth}{*{12}{>{\centering\arraybackslash}X}}
 \textcolor{gregoriocolor}{A} & \textcolor{gregoriocolor}{B} & \textcolor{gregoriocolor}{C} & \textcolor{gregoriocolor}{D} & \textcolor{gregoriocolor}{E} & F & \textcolor{gregoriocolor}{F} & \textcolor{gregoriocolor}{G} & \textcolor{gregoriocolor}{H} & \textcolor{gregoriocolor}{M} & \textcolor{gregoriocolor}{N} & \textcolor{gregoriocolor}{P} \\
 18 & 19 & 20 & 21 & 22 & 23 & 23 & 24 & 25 & 26 & 27 & 28 \\
\end{tabularx}

\begin{paracol}{2}
\selectlanguage{latin}
\lettrine[lines=2]{I}{n} óppido Marpúrgi, in 
 Germánia, deposítio sanctæ Elísabeth Víduæ, Regis Hungarórum Andréæ fíliæ, 
 ex tértio Ordine sancti Francísci, quæ, pietátis opéribus assídue inténta, 
 miráculis clara migrávit ad Dóminum.
\switchcolumn
\selectlanguage{english}
\lettrine[lines=2]{A}{t} Marburg in Germany, the death of 
 St. Elizabeth, widow, daughter of King Andrew of Hungary, and member of the 
 Third Order of St. Francis. After a life spent in the performance of 
 works of piety, she went to heaven, having a reputation for miracles.
\switchcolumn*
\selectlanguage{latin}
Sancti Pontiáni, Papæ 
 et Mártyris; cujus dies natális tértio Kaléndas Novémbris recensétur.
\switchcolumn
\selectlanguage{english}
St. Pontianus, pope and martyr, whose 
 birthday occurs on the 30th of October.
\switchcolumn*
\selectlanguage{latin}
Samaríæ, in Palæstína, 
 sancti Abdíæ Prophétæ.
\switchcolumn
\selectlanguage{english}
At Samaria in Palestine, the holy 
 prophet Abdias.
\switchcolumn*
\selectlanguage{latin}
Romæ, via Appia, 
 natális sancti Máximi, Presbyteri et Mártyris, qui, in persecutióne 
 Valeriáni passus, pósitus est ad sanctum Xystum.
\switchcolumn
\selectlanguage{english}
At Rome, on the Appian Way, the 
 birthday of St. Maximus, priest and martyr, who suffered in the persecution 
 of Valerian and was buried near St. Sixtus.
\switchcolumn*
\selectlanguage{latin}
In civitáte Astiagénsi, 
 in Hispánia, beáti Crispíni Epíscopi, qui, cápite amputáto, martyrii glóriam 
 adéptus est.
\switchcolumn
\selectlanguage{english}
At Ecijo in Spain, blessed Bishop 
 Crispin, who obtained the glory of martyrdom by beheading.
\switchcolumn*
\selectlanguage{latin}
Eódem die sancti Fausti, 
 Diáconi Alexandríni, qui primum, in persecutióne Valeriáni, cum sancto 
 Dionysio, in exsílium missus est; deínde, ætáte longævus, in persecutióne 
 Diocletiáni, animadvérsus gládio martyrium consummávit.
\switchcolumn
\selectlanguage{english}
St. Faustus, deacon of Alexandria, 
 who had been banished with St. Denis in the persecution of Valerian; later, 
 in the persecution of Diocletian, being advanced in age, his martyrdom was 
 accomplished by the sword.
\switchcolumn*
\selectlanguage{latin}
Cæsaréæ, in Cappadócia, sancti Bárlaam Mártyris, qui, agréstis licet et rudis, Christi sapiéntia 
 munítus tyránnum vicit, et ignem ipsum per invíctam fídei constántiam 
 superávit; in cujus die natáli sanctus Basilíus Magnus célebrem hábuit 
 oratiónem.
\switchcolumn
\selectlanguage{english}
At Caesarea in Cappadocia, St. 
 Barlaam, martyr, who, though unpolished and ignorant, was armed with the 
 wisdom of Christ to overcome the tyrant, and by the constancy of his faith, 
 subdue fire itself. On his birthday, St. Basil the Great delivered a 
 celebrated sermon.
\switchcolumn*
\selectlanguage{latin}
Viénnæ, in Gállia, 
 sanctórum Mártyrum Severíni, Exsupérii et Feliciáni; quorum córpora, post multa annórum currícula, ipsis revelántibus, invénta, et a Pontífice, clero 
 et pópulo illíus urbis honorífice subláta, condígno honóre cóndita sunt.
\switchcolumn
\selectlanguage{english}
At Vienne in France, the holy 
 martyrs Severinus, Exuperius and Felician. Their bodies, after the 
 lapse of many years, were found through their own revelation, and being 
 taken up with due honours by the bishop, clergy, and people of that city, 
 were buried with becoming solemnity.
\switchcolumn*
\selectlanguage{latin}
In Isáuria pássio 
 sanctórum Azæ et Sociórum centum quinquagínta mílitum, sub Diocletiáno 
 Imperatóre et Aquilíno Tribúno.
\switchcolumn
\selectlanguage{english}
In Isauria the martyrdom of St. Azas 
 and his soldier companions, to the number of one hundred and fifty, under 
 Emperor Diocletian and the tribune Aquilinus.
\switchcolumn*
\selectlanguage{latin}
\end{paracol}


% ---- martyrology/mart11/mart1120.htm
\needspace{10\baselineskip}
\begin{paracol}{2}
\selectlanguage{latin}
\begin{center}{\color{gregoriocolor} Duodécimo Kaléndas Decémbris. 
 Luna\dots\ }\end{center}
\switchcolumn
\selectlanguage{english}
\begin{center}{\color{gregoriocolor} The 
 20\textsuperscript{th} Day of 
 November. The\dots\ Day of the Moon.}\end{center}
\end{paracol}

\noindent\begin{tabularx}{\linewidth}{*{19}{>{\centering\arraybackslash}X}}
 \textcolor{gregoriocolor}{a} & \textcolor{gregoriocolor}{b} & \textcolor{gregoriocolor}{c} & \textcolor{gregoriocolor}{d} & \textcolor{gregoriocolor}{e} & \textcolor{gregoriocolor}{f} & \textcolor{gregoriocolor}{g} & \textcolor{gregoriocolor}{h} & \textcolor{gregoriocolor}{i} & \textcolor{gregoriocolor}{k} & \textcolor{gregoriocolor}{l} & \textcolor{gregoriocolor}{m} & \textcolor{gregoriocolor}{n} & \textcolor{gregoriocolor}{p} & \textcolor{gregoriocolor}{q} & \textcolor{gregoriocolor}{r} & \textcolor{gregoriocolor}{s} & \textcolor{gregoriocolor}{t} & \textcolor{gregoriocolor}{u} \\
 30 & 1 & 2 & 3 & 4 & 5 & 6 & 7 & 8 & 9 & 10 & 11 & 12 & 13 & 14 & 15 & 16 & 17 & 18 \\
\end{tabularx}
\vspace{0.5\baselineskip}
\noindent\begin{tabularx}{\linewidth}{*{12}{>{\centering\arraybackslash}X}}
 \textcolor{gregoriocolor}{A} & \textcolor{gregoriocolor}{B} & \textcolor{gregoriocolor}{C} & \textcolor{gregoriocolor}{D} & \textcolor{gregoriocolor}{E} & F & \textcolor{gregoriocolor}{F} & \textcolor{gregoriocolor}{G} & \textcolor{gregoriocolor}{H} & \textcolor{gregoriocolor}{M} & \textcolor{gregoriocolor}{N} & \textcolor{gregoriocolor}{P} \\
 19 & 20 & 21 & 22 & 23 & 24 & 24 & 25 & 26 & 27 & 28 & 29 \\
\end{tabularx}

\begin{paracol}{2}
\selectlanguage{latin}
\lettrine[lines=2]{S}{ancti} Felícis Valésii, 
 Presbyteri et Confessóris, qui Ordinis sanctíssimæ Trinitátis redemptiónis 
 captivórum éxstitit Fundátor, ac prídie Nonas Novémbris obdormívit in 
 Dómino.
\switchcolumn
\selectlanguage{english}
\lettrine[lines=2]{S}{t.} Felix of Valois, priest and 
 confessor, who founded the Order of the Most Holy Trinity for the Redemption 
 of Captives, and who fell asleep in the Lord on the 4th of November.
\switchcolumn*
\selectlanguage{latin}
In Pérside pássio 
 sanctórum Nersæ Epíscopi, et Sociórum.
\switchcolumn
\selectlanguage{english}
In Persia, the martyrdom of St. 
 Nersas, bishop, and his companions.
\switchcolumn*
\selectlanguage{latin}
Messánæ, in Sicília, 
 sanctórum Mártyrum Ampeli et Caji.
\switchcolumn
\selectlanguage{english}
At Messina in Sicily, the holy 
 martyrs Ampelus and Caius.
\switchcolumn*
\selectlanguage{latin}
Tauríni sanctórum 
 Mártyrum Octávii, Solutóris et Adventóris, Thebánæ legiónis mílitum; qui, 
 sub Maximiáno Imperatóre, egrégie decertántes, martyrio coronáti sunt.
\switchcolumn
\selectlanguage{english}
At Turin, the holy martyrs Octavius, 
 Solutor, and Adventor, soldiers of the Theban Legion, who fought valiantly 
 for the faith under Emperor Maximian and who were crowned with martyrdom.
\switchcolumn*
\selectlanguage{latin}
Cæsaréæ, in Palæstína, 
 sancti Agápii Mártyris, qui, sub Galério Maximiáno Imperatóre, damnátus ad 
 béstias et ab iis nil læsus, tandem, lapídibus ac pedes appénsis, in mare 
 demérgitur.
\switchcolumn
\selectlanguage{english}
At Caesarea in Palestine, in the 
 time of Emperor Galerius Maximian, the holy martyr Agapius, who was 
 condemned to be devoured by the beasts; but being unhurt by them, he was 
 cast into the sea with stones tied to his feet.
\switchcolumn*
\selectlanguage{latin}
Doróstori, in Mysia 
 inferióre, sancti Dásii Mártyris, qui, cum in festo Satúrni nollet 
 impudicítiis ejus consentíre, sub Basso Præfécto cæsus est.
\switchcolumn
\selectlanguage{english}
At Silistria in Rumania, St. Dasius, 
 bishop, who, for refusing to consent to the unholy rites of the Saturnalia, 
 was put to death under the governor Bassus.
\switchcolumn*
\selectlanguage{latin}
Nicéæ, in Bithynia, 
 sanctórum Mártyrum Eustáchii, Thespésii et Anatólii, in persecutióne 
 Maximíni.
\switchcolumn
\selectlanguage{english}
At Nicaea in Bithynia, the holy 
 martyrs Eustace, Thespesius, and Anatolius, in the persecution of Maximinus.
\switchcolumn*
\selectlanguage{latin}
Heracléæ, in Thrácia, 
 sanctórum Mártyrum Bassi, Dionysii, Agapíti et aliórum quadragínta.
\switchcolumn
\selectlanguage{english}
At Heraclea in Thrace, the holy 
 martyrs Bassus, Denis, Agapitus, and forty others.
\switchcolumn*
\selectlanguage{latin}
In Anglia sancti 
 Eadmúndi, Regis et Mártyris.
\switchcolumn
\selectlanguage{english}
In England, St. Edmund, king and 
 martyr.
\switchcolumn*
\selectlanguage{latin}
Constantinópoli sancti 
 Gregórii Decapolítæ, qui ob cultum sanctárum Imáginum multa passus est.
\switchcolumn
\selectlanguage{english}
At Constantinople, St. Gregory of 
 Decapolis, who suffered many things for the veneration of sacred images.
\switchcolumn*
\selectlanguage{latin}
Medioláni sancti 
 Benígni Epíscopi, qui, in magna barbarórum perturbatióne, commíssam sibi 
 Ecclésiam summa constántia et religióne administrávit.
\switchcolumn
\selectlanguage{english}
At Milan, St. Benignus, bishop, who, 
 amid great troubles caused by the barbarians, governed the Church entrusted 
 to him with greatest constancy and piety.
\switchcolumn*
\selectlanguage{latin}
Cabillóne, in Gálliis, 
 sancti Silvéstri Epíscopi, qui quadragésimo secúndo sui sacerdótii anno, 
 plenus diérum atque virtútum, migrávit ad Dóminum.
\switchcolumn
\selectlanguage{english}
At Chalons in France, St. Sylvester, 
 bishop, who went to God in the forty-second year of his priesthood, full of 
 days and virtues.
\switchcolumn*
\selectlanguage{latin}
Verónæ sancti Simplícii, 
 Epíscopi et Confessóris.
\switchcolumn
\selectlanguage{english}
At Verona, St. Simplicius, bishop 
 and confessor.
\switchcolumn*
\selectlanguage{latin}
Hildeshémii, in Saxónia, 
 sancti Bernwárdi, Epíscopi et Confessóris, qui a Cælestíno Papa Tértio in 
 Sanctórum númerum adscríptus est.
\switchcolumn
\selectlanguage{english}
At Hildesheim in Saxony, St. 
 Bernward, bishop and confessor, who was numbered among the saints by Pope 
 Celestine III.
\switchcolumn*
\selectlanguage{latin}
\end{paracol}


% ---- martyrology/mart11/mart1121.htm
\needspace{10\baselineskip}
\begin{paracol}{2}
\selectlanguage{latin}
\begin{center}{\color{gregoriocolor} Undécimo Kaléndas Decémbris. 
 Luna\dots\ }\end{center}
\switchcolumn
\selectlanguage{english}
\begin{center}{\color{gregoriocolor} The 
 21\textsuperscript{st} Day of 
 November. The\dots\ Day of the Moon.}\end{center}
\end{paracol}

\noindent\begin{tabularx}{\linewidth}{*{19}{>{\centering\arraybackslash}X}}
 \textcolor{gregoriocolor}{a} & \textcolor{gregoriocolor}{b} & \textcolor{gregoriocolor}{c} & \textcolor{gregoriocolor}{d} & \textcolor{gregoriocolor}{e} & \textcolor{gregoriocolor}{f} & \textcolor{gregoriocolor}{g} & \textcolor{gregoriocolor}{h} & \textcolor{gregoriocolor}{i} & \textcolor{gregoriocolor}{k} & \textcolor{gregoriocolor}{l} & \textcolor{gregoriocolor}{m} & \textcolor{gregoriocolor}{n} & \textcolor{gregoriocolor}{p} & \textcolor{gregoriocolor}{q} & \textcolor{gregoriocolor}{r} & \textcolor{gregoriocolor}{s} & \textcolor{gregoriocolor}{t} & \textcolor{gregoriocolor}{u} \\
 1 & 2 & 3 & 4 & 5 & 6 & 7 & 8 & 9 & 10 & 11 & 12 & 13 & 14 & 15 & 16 & 17 & 18 & 19 \\
\end{tabularx}
\vspace{0.5\baselineskip}
\noindent\begin{tabularx}{\linewidth}{*{12}{>{\centering\arraybackslash}X}}
 \textcolor{gregoriocolor}{A} & \textcolor{gregoriocolor}{B} & \textcolor{gregoriocolor}{C} & \textcolor{gregoriocolor}{D} & \textcolor{gregoriocolor}{E} & F & \textcolor{gregoriocolor}{F} & \textcolor{gregoriocolor}{G} & \textcolor{gregoriocolor}{H} & \textcolor{gregoriocolor}{M} & \textcolor{gregoriocolor}{N} & \textcolor{gregoriocolor}{P} \\
 20 & 21 & 22 & 23 & 24 & 25 & 25 & 26 & 27 & 28 & 29 & 30 \\
\end{tabularx}

\begin{paracol}{2}
\selectlanguage{latin}
\lettrine[lines=2]{H}{ierosólymis} Præsentátio beátæ Dei Genitrícis Vírginis Maríæ in Templo.
\switchcolumn
\selectlanguage{english}
\lettrine[lines=2]{I}{n} the temple at Jerusalem, the 
 Presentation of the Blessed Virgin Mary, Mother of God.
\switchcolumn*
\selectlanguage{latin}
Eódem die natális beáti 
 Rufi, de quo sanctus Paulus Apóstolus ad Romános scribit.
\switchcolumn
\selectlanguage{english}
Also, the birthday of blessed Rufus, 
 mentioned by the apostle St. Paul in his Epistle to the Romans.
\switchcolumn*
\selectlanguage{latin}
Romæ pássio sanctórum 
 Celsi et Cleméntis.
\switchcolumn
\selectlanguage{english}
At Rome, the martyrdom of the Saints 
 Celsus and Clement.
\switchcolumn*
\selectlanguage{latin}
Rhemis, in Gállia, 
 sancti Albérti, Epíscopi Leodiénsis et Mártyris; qui pro tuénda 
 ecclesiástica libertáte necátus est.
\switchcolumn
\selectlanguage{english}
At Rheims, St. Albert, bishop of 
 Liege and martyr, who was put to death for defending the liberty of the 
 Church.
\switchcolumn*
\selectlanguage{latin}
Apud Ostia Tiberína 
 natális sanctórum Mártyrum Demétrii et Honórii.
\switchcolumn
\selectlanguage{english}
At Ostia, the holy martyrs Demetrius 
 and Honorius.
\switchcolumn*
\selectlanguage{latin}
In Hispánia sanctórum 
 Mártyrum Honórii, Eutychii et Stéphani.
\switchcolumn
\selectlanguage{english}
In Spain, the holy martyrs Honorius, 
 Eutychius, and Stephen.
\switchcolumn*
\selectlanguage{latin}
In Pamphylia sancti 
 Heliodóri Mártyris, in persecutióne Aureliáni, sub Aétio Præside. Post 
 eum vero ipsi tortóres, convérsi ad fidem, in mare demérsi sunt.
\switchcolumn
\selectlanguage{english}
In Pamphylia, St. Heliodorus, 
 martyr, in the persecution of Aurelian under the governor Aetius. 
 After his death his executioners were converted to the faith and were cast 
 into the sea.
\switchcolumn*
\selectlanguage{latin}
Romæ sancti Gelásii 
 Papæ Primi, doctrína et sanctitáte conspícui.
\switchcolumn
\selectlanguage{english}
At Rome, Pope St. Gelasius, 
 distinguished for learning and sanctity.
\switchcolumn*
\selectlanguage{latin}
Verónæ sancti Mauri, 
 Epíscopi et Confessóris.
\switchcolumn
\selectlanguage{english}
At Verona, St. Maur, bishop and 
 confessor.
\switchcolumn*
\selectlanguage{latin}
In monastério Bobiénsi 
 deposítio sancti Columbáni Abbátis, qui, multórum cœnobiórum Fundátor, 
 plurimórum Monachórum éxstitit Pater, multísque virtútibus clarus, in 
 senectúte bona quiévit.
\switchcolumn
\selectlanguage{english}
In the monastery of Bobbio, the 
 death of St. Columban, abbot who founded many monasteries and governed a 
 large number of monks. He died at an advanced age, celebrated for many 
 virtues.
\switchcolumn*
\selectlanguage{latin}
\end{paracol}


% ---- martyrology/mart11/mart1122.htm
\needspace{10\baselineskip}
\begin{paracol}{2}
\selectlanguage{latin}
\begin{center}{\color{gregoriocolor} Décimo Kaléndas Decémbris. 
 Luna\dots\ }\end{center}
\switchcolumn
\selectlanguage{english}
\begin{center}{\color{gregoriocolor} The 
 22\textsuperscript{nd} Day of 
 November. The\dots\ Day of the Moon.}\end{center}
\end{paracol}

\noindent\begin{tabularx}{\linewidth}{*{19}{>{\centering\arraybackslash}X}}
 \textcolor{gregoriocolor}{a} & \textcolor{gregoriocolor}{b} & \textcolor{gregoriocolor}{c} & \textcolor{gregoriocolor}{d} & \textcolor{gregoriocolor}{e} & \textcolor{gregoriocolor}{f} & \textcolor{gregoriocolor}{g} & \textcolor{gregoriocolor}{h} & \textcolor{gregoriocolor}{i} & \textcolor{gregoriocolor}{k} & \textcolor{gregoriocolor}{l} & \textcolor{gregoriocolor}{m} & \textcolor{gregoriocolor}{n} & \textcolor{gregoriocolor}{p} & \textcolor{gregoriocolor}{q} & \textcolor{gregoriocolor}{r} & \textcolor{gregoriocolor}{s} & \textcolor{gregoriocolor}{t} & \textcolor{gregoriocolor}{u} \\
 2 & 3 & 4 & 5 & 6 & 7 & 8 & 9 & 10 & 11 & 12 & 13 & 14 & 15 & 16 & 17 & 18 & 19 & 20 \\
\end{tabularx}
\vspace{0.5\baselineskip}
\noindent\begin{tabularx}{\linewidth}{*{12}{>{\centering\arraybackslash}X}}
 \textcolor{gregoriocolor}{A} & \textcolor{gregoriocolor}{B} & \textcolor{gregoriocolor}{C} & \textcolor{gregoriocolor}{D} & \textcolor{gregoriocolor}{E} & F & \textcolor{gregoriocolor}{F} & \textcolor{gregoriocolor}{G} & \textcolor{gregoriocolor}{H} & \textcolor{gregoriocolor}{M} & \textcolor{gregoriocolor}{N} & \textcolor{gregoriocolor}{P} \\
 21 & 22 & 23 & 24 & 25 & 26 & 26 & 27 & 28 & 29 & 30 & 1 \\
\end{tabularx}

\begin{paracol}{2}
\selectlanguage{latin}
\lettrine[lines=2]{S}{anctæ} Cæcíliæ, 
 Vírginis et Mártyris, quæ ad cæléstem Sponsum, próprio sánguine purpuráta, 
 transívit sextodécimo Kaléndas Octóbris.
\switchcolumn
\selectlanguage{english}
\lettrine[lines=2]{S}{t.} Cecilia, virgin and martyr, who 
 on the 16th of September, purpled with her own blood, departed to her 
 heavenly Spouse.
\switchcolumn*
\selectlanguage{latin}
Colóssis, in Phrygia, 
 sanctórum Philémonis et Apphíæ, sancti Pauli discipulórum; qui, sub Neróne 
 Imperatóre, cum Gentíles, in die festo Diánæ, invasíssent Ecclésiam, ambo 
 céteris fugiéntibus, tenti, jussu Artoclis Præsidis verbéribus cæsi sunt, 
 et, usque ad renes in fóveam inclúsi, lapídibus opprimúntur.
\switchcolumn
\selectlanguage{english}
At Colossae in Phrygia, during the 
 reign of Nero, Saints Philemon and Apphias, disciples of St. Paul. 
 When the heathen rushed into the church on the feast of Diana, they were 
 arrested and the rest of the Christians fled. By command of the 
 governor Artocles they were scourged, enclosed up to their waists in a pit, 
 then overwhelmed with stones.
\switchcolumn*
\selectlanguage{latin}
Romæ sancti Mauri 
 Mártyris, qui, cum ex Africa venísset ad sepúlcra Apostolórum, sub 
 Imperatóre Numeriáno et Urbis Præfécto Celeríno agonizávit.
\switchcolumn
\selectlanguage{english}
At Rome, St. Maurus, martyr. He 
 came from Africa to visit the tombs of the apostles, and suffered martyrdom 
 there under Celerinus, prefect of the city in the reign of Emperor Numerian.
\switchcolumn*
\selectlanguage{latin}
Antiochíæ Pisídiæ 
 pássio sanctórum Marci et Stéphani, sub Diocletiáno Imperatóre.
\switchcolumn
\selectlanguage{english}
At Antioch in Pisidia, the martyrdom 
 of the Saints Mark and Stephen, under Emperor Diocletian.
\switchcolumn*
\selectlanguage{latin}
Augustodúni sancti 
 Pragmátii, Epíscopi et Confessóris.
\switchcolumn
\selectlanguage{english}
At Autun, St. Pragmatius, bishop and 
 confessor.
\switchcolumn*
\selectlanguage{latin}
\end{paracol}


% ---- martyrology/mart11/mart1123.htm
\needspace{10\baselineskip}
\begin{paracol}{2}
\selectlanguage{latin}
\begin{center}{\color{gregoriocolor} Nono Kaléndas Decémbris. 
 Luna\dots\ }\end{center}
\switchcolumn
\selectlanguage{english}
\begin{center}{\color{gregoriocolor} The 
 23\textsuperscript{rd} Day of 
 November. The\dots\ Day of the Moon.}\end{center}
\end{paracol}

\noindent\begin{tabularx}{\linewidth}{*{19}{>{\centering\arraybackslash}X}}
 \textcolor{gregoriocolor}{a} & \textcolor{gregoriocolor}{b} & \textcolor{gregoriocolor}{c} & \textcolor{gregoriocolor}{d} & \textcolor{gregoriocolor}{e} & \textcolor{gregoriocolor}{f} & \textcolor{gregoriocolor}{g} & \textcolor{gregoriocolor}{h} & \textcolor{gregoriocolor}{i} & \textcolor{gregoriocolor}{k} & \textcolor{gregoriocolor}{l} & \textcolor{gregoriocolor}{m} & \textcolor{gregoriocolor}{n} & \textcolor{gregoriocolor}{p} & \textcolor{gregoriocolor}{q} & \textcolor{gregoriocolor}{r} & \textcolor{gregoriocolor}{s} & \textcolor{gregoriocolor}{t} & \textcolor{gregoriocolor}{u} \\
 3 & 4 & 5 & 6 & 7 & 8 & 9 & 10 & 11 & 12 & 13 & 14 & 15 & 16 & 17 & 18 & 19 & 20 & 21 \\
\end{tabularx}
\vspace{0.5\baselineskip}
\noindent\begin{tabularx}{\linewidth}{*{12}{>{\centering\arraybackslash}X}}
 \textcolor{gregoriocolor}{A} & \textcolor{gregoriocolor}{B} & \textcolor{gregoriocolor}{C} & \textcolor{gregoriocolor}{D} & \textcolor{gregoriocolor}{E} & F & \textcolor{gregoriocolor}{F} & \textcolor{gregoriocolor}{G} & \textcolor{gregoriocolor}{H} & \textcolor{gregoriocolor}{M} & \textcolor{gregoriocolor}{N} & \textcolor{gregoriocolor}{P} \\
 22 & 23 & 24 & 25 & 26 & 27 & 27 & 28 & 29 & 30 & 1 & 2 \\
\end{tabularx}

\begin{paracol}{2}
\selectlanguage{latin}
\lettrine[lines=2]{N}{atális} sancti 
 Cleméntis Primi, Papæ et Mártyris, qui, tértius post beátum Petrum Apóstolum, 
 Pontificátum ténuit, et, in persecutióne Trajáni, apud Chersonésum relegátus, 
 ibi, alligáta ad ejus collum ánchora, præcipitátus in mare, martyrio 
 coronátur. Ipsíus autem corpus, Hadriáno Secúndo Summo Pontífice, a 
 sanctis Cyríllo et Methódio frátribus Romam translátum, in Ecclésia quæ ejus 
 nómine ántea fúerat exstrúcta, honorífice recónditum est.
\switchcolumn
\selectlanguage{english}
\lettrine[lines=2]{T}{he} birthday of Pope St. Clement, 
 who held the sovereign pontificate the third after the blessed apostle 
 Peter. In the persecution of Trajan, he was banished to Chersonesus, 
 where, being thrown into the sea with an anchor tied to his neck, he was 
 crowned with martyrdom. During the pontificate of Pope Adrian II, his 
 body was translated to Rome by the brothers Saints Cyril and Methodius, and 
 buried with honour in the church that had already been built and named 
 for him.
\switchcolumn*
\selectlanguage{latin}
Romæ sanctæ Felicitátis 
 Mártyris, septem filiórum Mártyrum matris; quæ, post eos, jubénte Marco 
 Antoníno Imperatóre, decolláta est pro Christo.
\switchcolumn
\selectlanguage{english}
At Rome, St. Felicitas, mother of 
 seven martyred sons. After them she was beheaded for Christ by order 
 of Emperor Marcus Antoninus.
\switchcolumn*
\selectlanguage{latin}
Cyzici, in Hellespónto, 
 sancti Sisínii Mártyris, qui in persecutióne Diocletiáni Imperatóris, post 
 multa torménta, gládio cæsus est.
\switchcolumn
\selectlanguage{english}
At Cyzicum, in the Hellespont, St. 
 Sisinius, martyr, who after many torments was put to the sword in the 
 persecution of Diocletian.
\switchcolumn*
\selectlanguage{latin}
Eméritæ, in Hispánia, 
 sanctæ Lucrétiæ, Vírginis et Mártyris; quæ in eádem persecutióne, sub 
 Daciáno Præside, martyrium consummávit.
\switchcolumn
\selectlanguage{english}
At Merida in Spain, St. Lucretia, 
 virgin and martyr, whose martyrdom was fulfilled in the same persecution, 
 under the governor Dacian.
\switchcolumn*
\selectlanguage{latin}
Icónii, in Lycaónia, 
 sancti Amphilóchii Epíscopi, qui, sanctórum Basilíi et Gregórii Nazianzéni 
 in erémo sócius et in Episcopátu colléga, tandem, post multa quæ suscépit 
 pro cathólica fide certámina, sanctitáte et doctrína clarus, quiévit in 
 pace.
\switchcolumn
\selectlanguage{english}
At Iconium in Lycaonia, the holy 
 bishop Amphilochius, who was the companion of St. Basil and St. Gregory 
 Nazianzen in the desert, and their colleague in the episcopate. After 
 enduring many trials for the Catholic faith, he rested in peace, renowned 
 for holiness and learning.
\switchcolumn*
\selectlanguage{latin}
Agrigénti deposítio 
 sancti Gregórii Epíscopi.
\switchcolumn
\selectlanguage{english}
At Girgenti, the death of St. 
 Gregory, bishop.
\switchcolumn*
\selectlanguage{latin}
In óppido Hasbániæ, in 
 Bélgio, sancti Trudónis, Presbyteri et Confessóris, cujus nómine póstmodum 
 insignítum fuit tum monastérium illic ab eódem Sancto in suis prædiis 
 eréctum, tum ipsum óppidum in eo loco paulátim exstrúctum.
\switchcolumn
\selectlanguage{english}
In the town of Hasbein in Belgium, 
 St. Trudo, priest and confessor. Both the monastery which he had 
 erected on his land, and the town which soon afterwards arose, were later 
 named for him.
\switchcolumn*
\selectlanguage{latin}
\end{paracol}


% ---- martyrology/mart11/mart1124.htm
\needspace{10\baselineskip}
\begin{paracol}{2}
\selectlanguage{latin}
\begin{center}{\color{gregoriocolor} Octávo Kaléndas Decémbris. 
 Luna\dots\ }\end{center}
\switchcolumn
\selectlanguage{english}
\begin{center}{\color{gregoriocolor} The 
 24\textsuperscript{th} Day of 
 November. The\dots\ Day of the Moon.}\end{center}
\end{paracol}

\noindent\begin{tabularx}{\linewidth}{*{19}{>{\centering\arraybackslash}X}}
 \textcolor{gregoriocolor}{a} & \textcolor{gregoriocolor}{b} & \textcolor{gregoriocolor}{c} & \textcolor{gregoriocolor}{d} & \textcolor{gregoriocolor}{e} & \textcolor{gregoriocolor}{f} & \textcolor{gregoriocolor}{g} & \textcolor{gregoriocolor}{h} & \textcolor{gregoriocolor}{i} & \textcolor{gregoriocolor}{k} & \textcolor{gregoriocolor}{l} & \textcolor{gregoriocolor}{m} & \textcolor{gregoriocolor}{n} & \textcolor{gregoriocolor}{p} & \textcolor{gregoriocolor}{q} & \textcolor{gregoriocolor}{r} & \textcolor{gregoriocolor}{s} & \textcolor{gregoriocolor}{t} & \textcolor{gregoriocolor}{u} \\
 4 & 5 & 6 & 7 & 8 & 9 & 10 & 11 & 12 & 13 & 14 & 15 & 16 & 17 & 18 & 19 & 20 & 21 & 22 \\
\end{tabularx}
\vspace{0.5\baselineskip}
\noindent\begin{tabularx}{\linewidth}{*{12}{>{\centering\arraybackslash}X}}
 \textcolor{gregoriocolor}{A} & \textcolor{gregoriocolor}{B} & \textcolor{gregoriocolor}{C} & \textcolor{gregoriocolor}{D} & \textcolor{gregoriocolor}{E} & F & \textcolor{gregoriocolor}{F} & \textcolor{gregoriocolor}{G} & \textcolor{gregoriocolor}{H} & \textcolor{gregoriocolor}{M} & \textcolor{gregoriocolor}{N} & \textcolor{gregoriocolor}{P} \\
 23 & 24 & 25 & 26 & 27 & 28 & 28 & 29 & 30 & 1 & 2 & 3 \\
\end{tabularx}

\begin{paracol}{2}
\selectlanguage{latin}
\lettrine[lines=2]{S}{ancti} Joánnis a Cruce 
 Presbyteri, Confessóris et Eccl\'esiæ Doctóris, sanctæ Terésiæ in 
 Carmelitárum reformatióne sócii, cujus dies natális décimo nono Kaléndas 
 Januárii recensétur.
\switchcolumn
\selectlanguage{english}
\lettrine[lines=2]{S}{t.} John of the Cross, priest and 
 confessor, and doctor of the Church, companion of St. Teresa in the reform 
 of Carmel, and whose birthday is the 14th of December.
\switchcolumn*
\selectlanguage{latin}
Eódem die natális 
 sancti Chrysógoni Mártyris, qui, post longa víncula et cárceres pro 
 constantíssima Christi confessióne tolerátos, Aquiléjam, jubénte Diocletiáno, 
 perdúctus, tandem, cæsus cápite et in mare projéctus, martyrium consummávit.
\switchcolumn
\selectlanguage{english}
Also, the birthday of St. 
 Chrysogonus, martyr. After a long imprisonment in chains for the 
 constant confession of Christ, he was ordered by Diocletian to be taken to 
 Aquileia, where he completed his martyrdom by being beheaded and thrown into 
 the sea.
\switchcolumn*
\selectlanguage{latin}
Romæ sancti 
 Crescentiáni Mártyris, qui in passióne beáti Marcélli Papæ memorátur.
\switchcolumn
\selectlanguage{english}
At Rome, St. Crescentian, martyr, 
 whose name is mentioned in the Acts of blessed Pope Marcellus.
\switchcolumn*
\selectlanguage{latin}
Apud Corínthum sancti 
 Alexándri Mártyris, qui sub Juliáno Apóstata et Sallústio Præside, pro 
 Christi fide certávit usque ad mortem.
\switchcolumn
\selectlanguage{english}
At Corinth, St. Alexander, martyr, 
 who fought unto death for the faith of Christ, under Julian the Apostate and 
 the governor Sallust.
\switchcolumn*
\selectlanguage{latin}
Perúsiæ sancti 
 Felicíssimi Mártyris.
\switchcolumn
\selectlanguage{english}
At Perugia, St. Felicissimus, 
 martyr.
\switchcolumn*
\selectlanguage{latin}
Amériæ, in Umbria, 
 sanctæ Firmínæ, Vírginis et Mártyris; quæ, in persecutióne Diocletiáni 
 Imperatóris, várie cruciáta est, ac demum, suspénsa et lampádibus ardéntibus 
 adústa, immaculátum spíritum Deo réddidit.
\switchcolumn
\selectlanguage{english}
At Amelia in Umbria, during the 
 persecution of Diocletian, St. Firmina, virgin and martyr. After being 
 subjected to various torments, to hanging, and to burning with flaming 
 torches, she yielded up her spirit.
\switchcolumn*
\selectlanguage{latin}
Córdubæ, in Hispánia, 
 sanctárum Vírginum et Mártyrum Floræ et Maríæ; quæ, post diutúrnos cárceres, 
 in persecutióne Arábica, gládio interémptæ sunt.
\switchcolumn
\selectlanguage{english}
At Cordova in Spain, the holy 
 virgins and martyrs Flora and Mary, who after a long imprisonment were slain 
 with the sword in the Arab persecution.
\switchcolumn*
\selectlanguage{latin}
Medioláni sancti 
 Protásii Epíscopi, qui apud Constántem Imperatórem in Concílio Sardicénsi 
 causam Athanásii deféndit, ac demum, pro Ecclésia sibi commíssa et pro 
 religióne multis perfúnctus labóribus, migrávit ad Dóminum.
\switchcolumn
\selectlanguage{english}
At Milan, St. Protase, bishop, who 
 defended the cause of Athanasius before Emperor Constans in the Council of 
 Sardica. Having sustained many labours for the church entrusted to him 
 and for religion, he departed this life to go to the Lord.
\switchcolumn*
\selectlanguage{latin}
In território 
 Arvernénsi sancti Portiáni Abbátis, qui, sub Theodoríco Rege, miráculis 
 cláruit; cujus étiam nomen índitum mansit tam monastério cui Sanctus ipse 
 præfuit, quam óppido quod in eódem loco póstea constrúctum fuit.
\switchcolumn
\selectlanguage{english}
In the territory of Auvergne, St. 
 Portian, an abbot who was renowned for miracles in the time of King 
 Theodoric. His name was given to the monastery that he had governed 
 and also the town which was later built there.
\switchcolumn*
\selectlanguage{latin}
In castro Blávio, in 
 Gállia, sancti Románi Presbyteri, cujus sanctitátis præcónium glória 
 miraculórum declárat.
\switchcolumn
\selectlanguage{english}
In the town of Blaye in France, St. 
 Romanus, priest, whose holiness is proclaimed by glorious miracles.
\switchcolumn*
\selectlanguage{latin}
\end{paracol}


% ---- martyrology/mart11/mart1125.htm
\needspace{10\baselineskip}
\begin{paracol}{2}
\selectlanguage{latin}
\begin{center}{\color{gregoriocolor} Séptimo Kaléndas Decémbris. 
 Luna\dots\ }\end{center}
\switchcolumn
\selectlanguage{english}
\begin{center}{\color{gregoriocolor} The 
 25\textsuperscript{th} Day of 
 November. The\dots\ Day of the Moon.}\end{center}
\end{paracol}

\noindent\begin{tabularx}{\linewidth}{*{19}{>{\centering\arraybackslash}X}}
 \textcolor{gregoriocolor}{a} & \textcolor{gregoriocolor}{b} & \textcolor{gregoriocolor}{c} & \textcolor{gregoriocolor}{d} & \textcolor{gregoriocolor}{e} & \textcolor{gregoriocolor}{f} & \textcolor{gregoriocolor}{g} & \textcolor{gregoriocolor}{h} & \textcolor{gregoriocolor}{i} & \textcolor{gregoriocolor}{k} & \textcolor{gregoriocolor}{l} & \textcolor{gregoriocolor}{m} & \textcolor{gregoriocolor}{n} & \textcolor{gregoriocolor}{p} & \textcolor{gregoriocolor}{q} & \textcolor{gregoriocolor}{r} & \textcolor{gregoriocolor}{s} & \textcolor{gregoriocolor}{t} & \textcolor{gregoriocolor}{u} \\
 5 & 6 & 7 & 8 & 9 & 10 & 11 & 12 & 13 & 14 & 15 & 16 & 17 & 18 & 19 & 20 & 21 & 22 & 23 \\
\end{tabularx}
\vspace{0.5\baselineskip}
\noindent\begin{tabularx}{\linewidth}{*{12}{>{\centering\arraybackslash}X}}
 \textcolor{gregoriocolor}{A} & \textcolor{gregoriocolor}{B} & \textcolor{gregoriocolor}{C} & \textcolor{gregoriocolor}{D} & \textcolor{gregoriocolor}{E} & F & \textcolor{gregoriocolor}{F} & \textcolor{gregoriocolor}{G} & \textcolor{gregoriocolor}{H} & \textcolor{gregoriocolor}{M} & \textcolor{gregoriocolor}{N} & \textcolor{gregoriocolor}{P} \\
 24 & 25 & 26 & 27 & 28 & 29 & 29 & 30 & 1 & 2 & 3 & 4 \\
\end{tabularx}

\begin{paracol}{2}
\selectlanguage{latin}
\lettrine[lines=2]{A}{lexandríæ} sanctæ 
 Catharínæ, Vírginis et Mártyris, quæ, ob fídei Christiánæ confessiónem, sub 
 Maximíno Imperatóre, in cárcerem trusa, et póstmodum scorpiónibus diutíssime 
 cæsa, tandem cápitis obtruncatióne martyrium complévit. Ipsíus corpus, 
 in montem Sínai mirabíliter ab Angelis delátum, ibídem, frequénti 
 Christianórum concúrsu, pia veneratióne cólitur.
\switchcolumn
\selectlanguage{english}
\lettrine[lines=2]{A}{t} Alexandria, St. Catherine, virgin 
 and martyr, in the time of Emperor Maximinus. For the confession of 
 the Christian faith she was cast into prison, endured a long scourging with 
 whips set with metal, and finally ended her martyrdom by having her head cut 
 off. Her body was miraculously carried by angels to Mount Sinai, where 
 pious veneration is paid to it by great gatherings of Christians.
\switchcolumn*
\selectlanguage{latin}
Romæ sancti Móysis, 
 Presbyteris et Mártyris; quem, cum áliis deténtum in cárcere, sanctus 
 Cypriánus per lítteras sæpe est consolátus. Ipse autem Móyses, cum non 
 tantum advérsus Gentíles, sed étiam advérsus schismáticos et hæréticos 
 Novatiános infrácto ánimo stetísset, demum (ut sanctus Cornélius Papa 
 testátur), in persecutióne Décii, exímio et admirábili martyrio decorátus 
 est.
\switchcolumn
\selectlanguage{english}
At Rome, St. Moses, priest and 
 martyr, who, along with others detained in prison, was often consoled by the 
 letters of St. Cyprian. He withstood with unbending courage not only 
 the heathen, but also the Novatian schismatics and heretics, and according 
 to the words of Pope St. Cornelius, he was finally crowned with a martyrdom 
 which fills the mind with admiration in the persecution of Decius.
\switchcolumn*
\selectlanguage{latin}
Antiochíæ sancti Erásmi 
 Mártyris.
\switchcolumn
\selectlanguage{english}
At Antioch, St. Erasmus, martyr.
\switchcolumn*
\selectlanguage{latin}
Cæsaréæ, in Cappadócia, pássio sancti Mercúrii mílitis, qui custodiéntis se Angeli patrocínio et 
 bárbaros vicit, et Décii sævítiam superávit; multísque auctus tormentórum 
 trophæis, martyrio coronátus migrávit in cælum.
\switchcolumn
\selectlanguage{english}
At Caesarea in Cappadocia, St. 
 Mercury, a soldier, who vanquished the barbarians and triumphed over the 
 cruelty of Decius through the protection of his guardian angel. 
 Finally, having acquired great glory from his sufferings, he was crowned 
 with martyrdom and went to reign forever in heaven.
\switchcolumn*
\selectlanguage{latin}
In Æmília, Itáliæ 
 província, sanctæ Jucúndæ Vírginis.
\switchcolumn
\selectlanguage{english}
In Emilia, a province of Italy, St. 
 Jucunda, virgin.
\switchcolumn*
\selectlanguage{latin}
\end{paracol}


% ---- martyrology/mart11/mart1126.htm
\needspace{10\baselineskip}
\begin{paracol}{2}
\selectlanguage{latin}
\begin{center}{\color{gregoriocolor} Sexto Kaléndas Decémbris. 
 Luna\dots\ }\end{center}
\switchcolumn
\selectlanguage{english}
\begin{center}{\color{gregoriocolor} The 
 26\textsuperscript{th} Day of 
 November. The\dots\ Day of the Moon.}\end{center}
\end{paracol}

\noindent\begin{tabularx}{\linewidth}{*{19}{>{\centering\arraybackslash}X}}
 \textcolor{gregoriocolor}{a} & \textcolor{gregoriocolor}{b} & \textcolor{gregoriocolor}{c} & \textcolor{gregoriocolor}{d} & \textcolor{gregoriocolor}{e} & \textcolor{gregoriocolor}{f} & \textcolor{gregoriocolor}{g} & \textcolor{gregoriocolor}{h} & \textcolor{gregoriocolor}{i} & \textcolor{gregoriocolor}{k} & \textcolor{gregoriocolor}{l} & \textcolor{gregoriocolor}{m} & \textcolor{gregoriocolor}{n} & \textcolor{gregoriocolor}{p} & \textcolor{gregoriocolor}{q} & \textcolor{gregoriocolor}{r} & \textcolor{gregoriocolor}{s} & \textcolor{gregoriocolor}{t} & \textcolor{gregoriocolor}{u} \\
 6 & 7 & 8 & 9 & 10 & 11 & 12 & 13 & 14 & 15 & 16 & 17 & 18 & 19 & 20 & 21 & 22 & 23 & 24 \\
\end{tabularx}
\vspace{0.5\baselineskip}
\noindent\begin{tabularx}{\linewidth}{*{12}{>{\centering\arraybackslash}X}}
 \textcolor{gregoriocolor}{A} & \textcolor{gregoriocolor}{B} & \textcolor{gregoriocolor}{C} & \textcolor{gregoriocolor}{D} & \textcolor{gregoriocolor}{E} & F & \textcolor{gregoriocolor}{F} & \textcolor{gregoriocolor}{G} & \textcolor{gregoriocolor}{H} & \textcolor{gregoriocolor}{M} & \textcolor{gregoriocolor}{N} & \textcolor{gregoriocolor}{P} \\
 25 & 26 & 27 & 28 & 29 & 1 & 30 & 1 & 2 & 3 & 4 & 5 \\
\end{tabularx}

\begin{paracol}{2}
\selectlanguage{latin}
\lettrine[lines=2]{A}{pud} Fabriánum, in 
 Picéno, beáti Silvéstri Abbátis, Institutóris Congregatiónis Monachórum 
 Silvestrinórum.
\switchcolumn
\selectlanguage{english}
\lettrine[lines=2]{A}{t} Fabriano in Piceno, St. 
 Sylvester, abbot, founder of the Congregation of Sylvestrine monks.
\switchcolumn*
\selectlanguage{latin}
Alexandríæ natális 
 sancti Petri, ejúsdem urbis Epíscopi et Mártyris; qui, cum esset ómnibus 
 virtútibus exornátus, ibídem, Galérii Maximiáni præcépto, cápite obtruncátus 
 est.
\switchcolumn
\selectlanguage{english}
At Alexandria, the birthday of St. 
 Peter, bishop of that city, graced with every virtue, who was beheaded by 
 command of Galerius Maximian.
\switchcolumn*
\selectlanguage{latin}
Passi sunt étiam 
 Alexandríæ, in eádem persecutióne, sancti Mártyres Faustus Présbyter, Dídius 
 et Ammónius, itémque Epíscopi quátuor Ægyptii, idest Philéas, Hesychius, 
 Pachómius et Theodórus, cum áliis sexcéntis sexagínta, quos persecutiónis 
 gládius evéxit ad cælos.
\switchcolumn
\selectlanguage{english}
There suffered also at Alexandria in 
 the same persecution the holy martyrs Faustus, a priest, Didius, and 
 Ammonius; likewise four bishops of Egypt, Phileas, Hesychius, Pachomius, and 
 Theodore, with others numbering six hundred and sixty, whom the sword of 
 persecution sent to heaven.
\switchcolumn*
\selectlanguage{latin}
Apud villam cui nomen 
 Fracta, in território Rhodigiénsi, sancti Bellíni, Epíscopi Patavíni et 
 Mártyris; qui a sicáriis, cum esset Ecclésiæ júrium defénsor exímius, 
 crudéliter impetítus ac multis illátis vulnéribus occísus est.
\switchcolumn
\selectlanguage{english}
In the village of Fracta, St. 
 Bellinus, bishop of Padua and martyr. The noble defender of the rights 
 of the Church was cruelly attacked by assassins, inflicting many wounds upon 
 him, and then slaying him.
\switchcolumn*
\selectlanguage{latin}
Nicomedíæ sancti 
 Marcélli Presbyteri, qui, Constántii témpore, ab Ariánis e rupe præcipitátus, 
 Martyr occúbuit.
\switchcolumn
\selectlanguage{english}
At Nicomedia, in the time of 
 Constantius, St. Marcellus, a priest, who died a martyr by being hurled from 
 a rock by the Arians.
\switchcolumn*
\selectlanguage{latin}
Romæ sancti Sirícii, 
 Papæ et Confessóris, doctrína, pietáte et religiónis zelo præclári, qui 
 vários damnávit hæréticos, et disciplínam ecclesiásticam salubérrimis 
 decrétis instaurávit.
\switchcolumn
\selectlanguage{english}
At Rome, St. Siricius, pope and 
 confessor, celebrated for his learning, piety, and zeal for religion, who 
 condemned various heretics and published salutary laws concerning 
 ecclesiastical discipline.
\switchcolumn*
\selectlanguage{latin}
Augustodúni sancti 
 Amatóris Epíscopi.
\switchcolumn
\selectlanguage{english}
At Autun, St. Amator, bishop.
\switchcolumn*
\selectlanguage{latin}
Constántiæ, in Germánia, sancti Conrádi Epíscopi.
\switchcolumn
\selectlanguage{english}
At Constance in Germany, St. Conrad, 
 bishop.
\switchcolumn*
\selectlanguage{latin}
Romæ sancti Leonárdi, a 
 Portu Maurítio, Sacerdótis ex Ordine Minórum et Confessóris, zelo animárum 
 et sacris per Itáliam expeditiónibus conspícui; quem Pius Nonus, Póntifex 
 Máximus, in Sanctórum cánonem rétulit, ac Pius Papa Undécimus cæléstem 
 Patrónum Sacerdótum qui ad sacras populáres Missiónes in regiónibus 
 cathólicis ubíque terrárum incúmbunt, elégit et constítuit.
\switchcolumn
\selectlanguage{english}
At Rome, St. Leonard of Port 
 Maurice, priest and confessor of the Order of Friars Minor. He was 
 remarkable for his zeal for souls and his holy expeditions throughout Italy. 
 He was canonized by Pope Pius IX, and Pope Pius XI chose and appointed him 
 the heavenly patron of priests to the preaching of missions to the people.
\switchcolumn*
\selectlanguage{latin}
In território Rheménsi 
 natális sancti Básoli Confessóris.
\switchcolumn
\selectlanguage{english}
In the district of Rheims, the 
 birthday of St. Basolus, confessor.
\switchcolumn*
\selectlanguage{latin}
Hadrianópoli, in 
 Paphlagónia, sancti Styliáni Anachorétæ, miráculis clari.
\switchcolumn
\selectlanguage{english}
At Adrianople in Paphlagonia, St. 
 Stylian, anchoret, renowned for miracles.
\switchcolumn*
\selectlanguage{latin}
In Arménia sancti 
 Nicónis Mónachi.
\switchcolumn
\selectlanguage{english}
In Armenia, St. Nicon, monk.
\switchcolumn*
\selectlanguage{latin}
\end{paracol}


% ---- martyrology/mart11/mart1127.htm
\needspace{10\baselineskip}
\begin{paracol}{2}
\selectlanguage{latin}
\begin{center}{\color{gregoriocolor} Quinto Kaléndas Decémbris. 
 Luna\dots\ }\end{center}
\switchcolumn
\selectlanguage{english}
\begin{center}{\color{gregoriocolor} The 
 27\textsuperscript{th} Day of 
 November. The\dots\ Day of the Moon.}\end{center}
\end{paracol}

\noindent\begin{tabularx}{\linewidth}{*{19}{>{\centering\arraybackslash}X}}
 \textcolor{gregoriocolor}{a} & \textcolor{gregoriocolor}{b} & \textcolor{gregoriocolor}{c} & \textcolor{gregoriocolor}{d} & \textcolor{gregoriocolor}{e} & \textcolor{gregoriocolor}{f} & \textcolor{gregoriocolor}{g} & \textcolor{gregoriocolor}{h} & \textcolor{gregoriocolor}{i} & \textcolor{gregoriocolor}{k} & \textcolor{gregoriocolor}{l} & \textcolor{gregoriocolor}{m} & \textcolor{gregoriocolor}{n} & \textcolor{gregoriocolor}{p} & \textcolor{gregoriocolor}{q} & \textcolor{gregoriocolor}{r} & \textcolor{gregoriocolor}{s} & \textcolor{gregoriocolor}{t} & \textcolor{gregoriocolor}{u} \\
 7 & 8 & 9 & 10 & 11 & 12 & 13 & 14 & 15 & 16 & 17 & 18 & 19 & 20 & 21 & 22 & 23 & 24 & 25 \\
\end{tabularx}
\vspace{0.5\baselineskip}
\noindent\begin{tabularx}{\linewidth}{*{12}{>{\centering\arraybackslash}X}}
 \textcolor{gregoriocolor}{A} & \textcolor{gregoriocolor}{B} & \textcolor{gregoriocolor}{C} & \textcolor{gregoriocolor}{D} & \textcolor{gregoriocolor}{E} & F & \textcolor{gregoriocolor}{F} & \textcolor{gregoriocolor}{G} & \textcolor{gregoriocolor}{H} & \textcolor{gregoriocolor}{M} & \textcolor{gregoriocolor}{N} & \textcolor{gregoriocolor}{P} \\
 26 & 27 & 28 & 29 & 1 & 2 & 1 & 2 & 3 & 4 & 5 & 6 \\
\end{tabularx}

\begin{paracol}{2}
\selectlanguage{latin}
\lettrine[lines=2]{A}{ntiochíæ} sanctórum 
 Mártyrum Basiléi Epíscopi, Auxílii et Saturníni.
\switchcolumn
\selectlanguage{english}
\lettrine[lines=2]{A}{t} Antioch, the holy martyrs 
 Basileus, bishop, Auxilius, and Saturninus.
\switchcolumn*
\selectlanguage{latin}
Sebáste, in Arménia, sanctórum Mártyrum Hirenárchi, Acácii Presbyteri, ac septem mulíerum. Harum 
 porro constántia Hirenárchus commótus, ad Christum convérsus, sub 
 Diocletiáno Imperatóre et Máximo Præside, una cum Acácio, secúri percútitur.
\switchcolumn
\selectlanguage{english}
At Sebaste in Armenia, in the reign 
 of Emperor Diocletian and under the governor Maximus, the holy martyrs 
 Hirenarchus, the priest Acacius, and seven women. Struck with the 
 constancy of these women, Hirenarchus was converted to Christ, and with 
 Acacius died under the axe.
\switchcolumn*
\selectlanguage{latin}
Apud Cæam flúvium, in 
 Gallæcia, sanctórum Facúndi et Primitívi, qui sub Attico Præside passi sunt.
\switchcolumn
\selectlanguage{english}
In Galicia, on the River Cea, the 
 Saints Facundus and Primitivus, who suffered under the governor Atticus.
\switchcolumn*
\selectlanguage{latin}
In Pérside sancti 
 Jacóbi intercísi, Mártyris conspícui, qui, témpore Theodósii junióris, cum 
 in Isdegérdis Regis grátiam Christum negásset, et proptérea mater ejus et 
 uxor ab ipsíus se consuetúdine subtraxíssent, hinc, in se revérsus, 
 intrépide coram Vararáne, Isdegérdis fílio ac successóre, se Christiánum 
 esse conféssus est; ideóque ab iráto Rege, lata in eum mortis senténtia, 
 membrátim jussus est concídi et cápite obtruncári. Quo étiam témpore 
 innúmeri álii Mártyres ibídem passi sunt.
\switchcolumn
\selectlanguage{english}
In Persia, St. James Intercisus, a 
 distinguished martyr. In the time of Theodosius the Younger he denied 
 Christ in order to please King Isdegerd, but his mother and his wife for 
 this reason withdrew from his company. Coming to himself, he returned 
 to the king, now Vararanus, the son and successor of Isdegerd, to declare his faith in our Lord, whereupon the angry monarch 
 condemned him to be cut in pieces and beheaded. Countless other 
 martyrs suffered at this time in the same country.
\switchcolumn*
\selectlanguage{latin}
Aquiléjæ sancti 
 Valeriáni Epíscopi.
\switchcolumn
\selectlanguage{english}
At Aquileia, St. Valerian, bishop.
\switchcolumn*
\selectlanguage{latin}
Apud Régium, in Gállia, 
 sancti Máximi, Epíscopi et Confessóris; qui, usque a primævæ ætátis annis 
 omni virtútum grátia præditus, primum Lirinénsis cœnóbii Pater, deínde 
 Regiénsis Ecclésiæ Epíscopus, signis et prodígiis ínclytus éxstitit.
\switchcolumn
\selectlanguage{english}
At Riez in France, St. Maximus, 
 bishop and confessor, who, from his tender years, was endowed with every 
 grace and virtue. Being first superior of the monastery of Lerins, and 
 afterwards bishop of the Church of Riez, he was celebrated for the working 
 of miracles and prodigies.
\switchcolumn*
\selectlanguage{latin}
Salisbúrgi, in Nórico, 
 sancti Virgílii, Epíscopi et Carinthiórum Apóstoli, qui a Gregório Nono, 
 Pontífice Máximo, in Sanctórum númerum adscríptus est.
\switchcolumn
\selectlanguage{english}
At Salzburg in Austria, St. Virgil, 
 bishop and apostle of Carinthia, who was placed among the number of saints 
 by Pope Gregory IX.
\switchcolumn*
\selectlanguage{latin}
Apud Indos, Persis 
 finítimos, sanctórum Bárlaam et Jósaphat, quorum actus mirándos sanctus 
 Joánnes Damascénus conscrípsit.
\switchcolumn
\selectlanguage{english}
In India, near the Persian boundary, 
 the Saints Barlaam and Josaphat, whose wonderful deeds were written by St. 
 John of Damascus.
\switchcolumn*
\selectlanguage{latin}
Lutétiæ Parisiórum 
 deposítio sancti Severíni, Mónachi et Solitárii.
\switchcolumn
\selectlanguage{english}
At Paris, the death of St. Severin, 
 monk and solitary.
\switchcolumn*
\selectlanguage{latin}
\end{paracol}


% ---- martyrology/mart11/mart1128.htm
\needspace{10\baselineskip}
\begin{paracol}{2}
\selectlanguage{latin}
\begin{center}{\color{gregoriocolor} Quarto Kaléndas Decémbris. 
 Luna\dots\ }\end{center}
\switchcolumn
\selectlanguage{english}
\begin{center}{\color{gregoriocolor} The 
 28\textsuperscript{th} Day of 
 November. The\dots\ Day of the Moon.}\end{center}
\end{paracol}

\noindent\begin{tabularx}{\linewidth}{*{19}{>{\centering\arraybackslash}X}}
 \textcolor{gregoriocolor}{a} & \textcolor{gregoriocolor}{b} & \textcolor{gregoriocolor}{c} & \textcolor{gregoriocolor}{d} & \textcolor{gregoriocolor}{e} & \textcolor{gregoriocolor}{f} & \textcolor{gregoriocolor}{g} & \textcolor{gregoriocolor}{h} & \textcolor{gregoriocolor}{i} & \textcolor{gregoriocolor}{k} & \textcolor{gregoriocolor}{l} & \textcolor{gregoriocolor}{m} & \textcolor{gregoriocolor}{n} & \textcolor{gregoriocolor}{p} & \textcolor{gregoriocolor}{q} & \textcolor{gregoriocolor}{r} & \textcolor{gregoriocolor}{s} & \textcolor{gregoriocolor}{t} & \textcolor{gregoriocolor}{u} \\
 8 & 9 & 10 & 11 & 12 & 13 & 14 & 15 & 16 & 17 & 18 & 19 & 20 & 21 & 22 & 23 & 24 & 25 & 26 \\
\end{tabularx}
\vspace{0.5\baselineskip}
\noindent\begin{tabularx}{\linewidth}{*{12}{>{\centering\arraybackslash}X}}
 \textcolor{gregoriocolor}{A} & \textcolor{gregoriocolor}{B} & \textcolor{gregoriocolor}{C} & \textcolor{gregoriocolor}{D} & \textcolor{gregoriocolor}{E} & F & \textcolor{gregoriocolor}{F} & \textcolor{gregoriocolor}{G} & \textcolor{gregoriocolor}{H} & \textcolor{gregoriocolor}{M} & \textcolor{gregoriocolor}{N} & \textcolor{gregoriocolor}{P} \\
 27 & 28 & 29 & 1 & 2 & 3 & 2 & 3 & 4 & 5 & 6 & 7 \\
\end{tabularx}

\begin{paracol}{2}
\selectlanguage{latin}
\lettrine[lines=2]{A}{pud} Corínthum natális 
 sancti Sósthenis, ex beáti Pauli Apóstoli discípulis; cujus mentiónem facit 
 idem Apóstolus Corínthiis scribens. Ipse autem Sósthenes, ex príncipe 
 Synagógæ convérsus ad Christum, fídei suæ primórdia, ante Galliónem 
 Procónsulem ácriter verberátus, præcláro inítio consecrávit.
\switchcolumn
\selectlanguage{english}
\lettrine[lines=2]{A}{t} Corinth, the birthday of St. 
 Sosthenes, disciple of the blessed apostle Paul, who is mentioned in his 
 Epistle to the Corinthians. He was chief of the synagogue when 
 converted to Christ, and as a glorious beginning, consecrated the first 
 fruits of his faith by being scourged before the proconsul Gallio.
\switchcolumn*
\selectlanguage{latin}
Romæ sancti Rufi, quem, 
 cum omni família sua, Christi Mártyrem Diocletiánus fecit.
\switchcolumn
\selectlanguage{english}
At Rome, St. Rufus, who was martyred 
 with all his family by Diocletian.
\switchcolumn*
\selectlanguage{latin}
In Africa sanctórum 
 Mártyrum Papiniáni et Mansuéti Episcopórum, qui, témpore Wandálicæ 
 persecutiónis, sub Rege Ariáno Genseríco, pro fídei cathólicæ defensióne, 
 candéntibus ferri láminis toto córpore adústi, gloriósum agónem consummárunt. 
 Quo étiam témpore álii novem sancti Epíscopi, scílicet Valeriánus, Urbánus, 
 Crescens, Eustáchius, Crescónius, Crescentiánus, Felix, Hortulánus et 
 Florentiánus, damnáti exsílio, cursum vitæ suæ implevérunt.
\switchcolumn
\selectlanguage{english}
In Africa, under the Arian king 
 Genseric, in the persecution of the Vandals, the holy martyrs Papinian and 
 Mansuetus, bishops, who, for the Catholic faith, were burned in every part 
 of their bodies with hot plates of iron, which ended their glorious trial. 
 At this time also, other holy bishops, Valerian, Urban, Crescens, Eustachius, 
 Cresconius, Crescentian, Felix, Hortulanus, and Florentian ended the course 
 of their lives in exile.
\switchcolumn*
\selectlanguage{latin}
Constantinópoli 
 sanctórum Mártyrum Stéphani junióris, Basilíi, Petri, Andréæ, et Sociórum 
 trecentórum et trigínta novem Monachórum; qui, sub Constantíno Coprónymo, 
 pro sanctárum Imáginum cultu váriis excruciáti supplíciis, veritátem 
 cathólicam effúso sánguine confirmárunt.
\switchcolumn
\selectlanguage{english}
At Constantinople, in the time of 
 Constantine Copronymus, the holy martyrs Stephen the Younger, Basil, Peter, 
 Andrew, and their companions, numbering three hundred and thirty-nine monks, 
 who were subjected to diverse torments for the veneration of holy images, 
 and confirmed the Catholic truth with the shedding of their blood.
\switchcolumn*
\selectlanguage{latin}
Neápoli, in Campánia, deposítio sancti Jacóbi Picéni, Sacerdótis ex Ordine Minórum et Confessóris, 
 vitæ asperitáte, apostólica prædicatióne ac plúribus pro re Christiána óbitis legatiónibus præclári; quem Benedíctus Décimus tértius, Póntifex 
 Máximus, Sanctórum fastis adjúnxit.
\switchcolumn
\selectlanguage{english}
At Naples in Campania, the death of 
 St. James della Marca, priest and confessor of the Order of Friars Minor, 
 celebrated for the austerity of his life, his apostolic preaching, and his 
 many diplomatic missions undertaken for the success of the affairs of 
 Christianity. His name was added to the calendar of the saints by the 
 Sovereign Pontiff, Benedict XIII.
\switchcolumn*
\selectlanguage{latin}
\end{paracol}


% ---- martyrology/mart11/mart1129.htm
\needspace{10\baselineskip}
\begin{paracol}{2}
\selectlanguage{latin}
\begin{center}{\color{gregoriocolor} Tértio Kaléndas Decémbris. 
 Luna\dots\ }\end{center}
\switchcolumn
\selectlanguage{english}
\begin{center}{\color{gregoriocolor} The 
 29\textsuperscript{th} Day of 
 November. The\dots\ Day of the Moon.}\end{center}
\end{paracol}

\noindent\begin{tabularx}{\linewidth}{*{19}{>{\centering\arraybackslash}X}}
 \textcolor{gregoriocolor}{a} & \textcolor{gregoriocolor}{b} & \textcolor{gregoriocolor}{c} & \textcolor{gregoriocolor}{d} & \textcolor{gregoriocolor}{e} & \textcolor{gregoriocolor}{f} & \textcolor{gregoriocolor}{g} & \textcolor{gregoriocolor}{h} & \textcolor{gregoriocolor}{i} & \textcolor{gregoriocolor}{k} & \textcolor{gregoriocolor}{l} & \textcolor{gregoriocolor}{m} & \textcolor{gregoriocolor}{n} & \textcolor{gregoriocolor}{p} & \textcolor{gregoriocolor}{q} & \textcolor{gregoriocolor}{r} & \textcolor{gregoriocolor}{s} & \textcolor{gregoriocolor}{t} & \textcolor{gregoriocolor}{u} \\
 9 & 10 & 11 & 12 & 13 & 14 & 15 & 16 & 17 & 18 & 19 & 20 & 21 & 22 & 23 & 24 & 25 & 26 & 27 \\
\end{tabularx}
\vspace{0.5\baselineskip}
\noindent\begin{tabularx}{\linewidth}{*{12}{>{\centering\arraybackslash}X}}
 \textcolor{gregoriocolor}{A} & \textcolor{gregoriocolor}{B} & \textcolor{gregoriocolor}{C} & \textcolor{gregoriocolor}{D} & \textcolor{gregoriocolor}{E} & F & \textcolor{gregoriocolor}{F} & \textcolor{gregoriocolor}{G} & \textcolor{gregoriocolor}{H} & \textcolor{gregoriocolor}{M} & \textcolor{gregoriocolor}{N} & \textcolor{gregoriocolor}{P} \\
 28 & 29 & 1 & 2 & 3 & 4 & 3 & 4 & 5 & 6 & 7 & 8 \\
\end{tabularx}

\begin{paracol}{2}
\selectlanguage{latin}
\lettrine[lines=1]{V}{igília} sancti Andréæ 
 Apóstoli.
\switchcolumn
\selectlanguage{english}
\lettrine[lines=1]{T}{he} Vigil of St. Andrew, apostle.
\switchcolumn*
\selectlanguage{latin}
Romæ, via Salária, 
 natális sanctórum Mártyrum Saturníni senis, et Sisínii Diáconi, sub 
 Maximiáno Príncipe; quos, diu in cárcere macerátos, jussit Urbis Præféctus 
 in equúleum levári et áttrahi nervis, fústibus ac scorpiónibus cædi, deínde 
 eis flammas appóni, et, depósitos de equúleo, cápite truncári.
\switchcolumn
\selectlanguage{english}
At Rome, on the Salarian Way, the 
 birthday of the holy martyr, Saturninus, an aged man, and the deacon 
 Sisinius, in the time of Emperor Maximian. After a long imprisonment, 
 by order of the prefect of the city they were placed on the rack, stretched 
 with ropes, scourged with rods and whips garnished with metal, then exposed 
 to the flames, taken down from the rack and beheaded.
\switchcolumn*
\selectlanguage{latin}
Tolósæ sancti Saturníni 
 Epíscopi, qui, tempóribus Décii, in Capitólio ejúsdem urbis a Pagánis tentus 
 atque de summa Capitólii arce per omnes gradus præcipitátus est, atque ita, 
 cápite collíso, excussóque cérebro, et toto córpore dilaniáto, dignam 
 Christo ánimam réddidit.
\switchcolumn
\selectlanguage{english}
At Toulouse, in the time of Decius, 
 the holy bishop Saturninus, who was taken to the capitol of that city by the 
 heathen and thrown down the steps from the highest part of the building. 
 The fall having crushed his head, dashed out his brain and mangled his whole 
 body, he rendered his worthy soul to our Lord.
\switchcolumn*
\selectlanguage{latin}
Item pássio sanctórum 
 Parámonis et Sociórum trecentórum septuagínta quinque, sub Décio Imperatóre 
 et Aquilíno Præside.
\switchcolumn
\selectlanguage{english}
Also, the martyrdom of the Saints 
 Paramon and his companions, to the number of three hundred and seventy-five 
 under Emperor Decius and the governor Aquilinus.
\switchcolumn*
\selectlanguage{latin}
Ancyræ, in Galátia, sancti Philómeni Mártyris, qui, in persecutióne Aureliáni Imperatóris, sub 
 Felíce Præside, igne probátus, mánibus pedibúsque ac demum cápite clavis 
 confíxo, martyrium consummávit.
\switchcolumn
\selectlanguage{english}
At Ancyra in Galatia, St. Philomenus, 
 martyr. During the persecution of Emperor Aurelian, under the governor 
 Felix, he was first exposed to the flames, then having his hands, feet, and 
 head pierced with nails, he fulfilled his martyrdom.
\switchcolumn*
\selectlanguage{latin}
Vérulis, in Hérnicis, 
 sanctórum Mártyrum Blásii et Demétrii.
\switchcolumn
\selectlanguage{english}
At Veroli, the holy martyrs Blaise 
 and Demetrius.
\switchcolumn*
\selectlanguage{latin}
Tudérti, in Umbria, 
 sanctæ Illuminátæ Vírginis.
\switchcolumn
\selectlanguage{english}
At Todi in Umbria, St. Illuminata, 
 virgin.
\switchcolumn*
\selectlanguage{latin}
\end{paracol}


% ---- martyrology/mart11/mart1130.htm
\needspace{10\baselineskip}
\begin{paracol}{2}
\selectlanguage{latin}
\begin{center}{\color{gregoriocolor} Prídie Kaléndas Decémbris. 
 Luna\dots\ }\end{center}
\switchcolumn
\selectlanguage{english}
\begin{center}{\color{gregoriocolor} The 
 30\textsuperscript{th} Day of 
 November. The\dots\ Day of the Moon.}\end{center}
\end{paracol}

\noindent\begin{tabularx}{\linewidth}{*{19}{>{\centering\arraybackslash}X}}
 \textcolor{gregoriocolor}{a} & \textcolor{gregoriocolor}{b} & \textcolor{gregoriocolor}{c} & \textcolor{gregoriocolor}{d} & \textcolor{gregoriocolor}{e} & \textcolor{gregoriocolor}{f} & \textcolor{gregoriocolor}{g} & \textcolor{gregoriocolor}{h} & \textcolor{gregoriocolor}{i} & \textcolor{gregoriocolor}{k} & \textcolor{gregoriocolor}{l} & \textcolor{gregoriocolor}{m} & \textcolor{gregoriocolor}{n} & \textcolor{gregoriocolor}{p} & \textcolor{gregoriocolor}{q} & \textcolor{gregoriocolor}{r} & \textcolor{gregoriocolor}{s} & \textcolor{gregoriocolor}{t} & \textcolor{gregoriocolor}{u} \\
 10 & 11 & 12 & 13 & 14 & 15 & 16 & 17 & 18 & 19 & 20 & 21 & 22 & 23 & 24 & 25 & 26 & 27 & 28 \\
\end{tabularx}
\vspace{0.5\baselineskip}
\noindent\begin{tabularx}{\linewidth}{*{12}{>{\centering\arraybackslash}X}}
 \textcolor{gregoriocolor}{A} & \textcolor{gregoriocolor}{B} & \textcolor{gregoriocolor}{C} & \textcolor{gregoriocolor}{D} & \textcolor{gregoriocolor}{E} & F & \textcolor{gregoriocolor}{F} & \textcolor{gregoriocolor}{G} & \textcolor{gregoriocolor}{H} & \textcolor{gregoriocolor}{M} & \textcolor{gregoriocolor}{N} & \textcolor{gregoriocolor}{P} \\
 29 & 1 & 2 & 3 & 4 & 5 & 4 & 5 & 6 & 7 & 8 & 9 \\
\end{tabularx}

\begin{paracol}{2}
\selectlanguage{latin}
\lettrine[lines=2]{A}{pud} Patras, in Achája, 
 natális sancti Andréæ Apóstoli, qui in Thrácia et Scythia sacrum Christi 
 Evangélium prædicávit. Is, ab Ægéa Proncónsule comprehénsus, primum in cárcere clausus est, deínde gravíssime cæsus, ad 
 últimum suspénsus in cruce, 
 in ea pópulum docens bíduo supervíxit; et, rogáto Dómino ne eum síneret de 
 cruce depóni, circúmdatus est magno splendóre de cælo, et, abscedénte 
 póstmodum lúmine, emísit spíritum.
\switchcolumn
\selectlanguage{english}
\lettrine[lines=2]{A}{t} Patras in Achaia, the birthday of 
 the apostle St. Andrew, who preached the gospel of Christ in Thrace and 
 Sythia. He was apprehended by the proconsul Aegeas, imprisoned, and 
 severely scourged, and finally, being hung on a cross, he lived two days on 
 it, teaching the people. Having besought our Lord not to permit him to 
 be taken down from the cross, he was surrounded with a great brightness from 
 heaven, and when the light disappeared he breathed his last.
\switchcolumn*
\selectlanguage{latin}
Romæ pássio sanctórum 
 Cástuli et Euprépitis.
\switchcolumn
\selectlanguage{english}
At Rome, the martyrdom of the Saints 
 Castulus and Euprepis.
\switchcolumn*
\selectlanguage{latin}
Constantinópoli sanctæ 
 Mauræ, Vírginis et Mártyris.
\switchcolumn
\selectlanguage{english}
At Constantinople, St. Maura, virgin 
 and martyr.
\switchcolumn*
\selectlanguage{latin}
Item sanctæ Justínæ, 
 Vírginis et Mártyris.
\switchcolumn
\selectlanguage{english}
Also, St. Justina, virgin and 
 martyr.
\switchcolumn*
\selectlanguage{latin}
Romæ sancti Constántii 
 Confessóris, qui, fórtiter Pelagiánis resístens, ab eórum factióne pértulit 
 multa, quæ illum sanctis Confessóribus sociárunt.
\switchcolumn
\selectlanguage{english}
At Rome, St. Constantius, confessor, 
 who bravely opposed the Pelagians, and by enduring many injuries from them, 
 gained a place among the holy confessors.
\switchcolumn*
\selectlanguage{latin}
Apud Sántonas, in 
 Gállia, sancti Trojáni Epíscopi, magnæ sanctitátis viri, qui, sepúltus in 
 terris, se in cælis vívere multis virtútibus maniféstat.
\switchcolumn
\selectlanguage{english}
At Saintes in France, St. Trojan, 
 bishop and confessor, a man of great sanctity, who shews by many miracles 
 that he lives in heaven, though his body is buried on earth.
\switchcolumn*
\selectlanguage{latin}
In Palæstína beáti 
 Zósimi Confessóris, qui, sub Justíno Imperatóre, sanctitáte et miráculis 
 fuit insígnis.
\switchcolumn
\selectlanguage{english}
In Palestine, blessed Zosimus, 
 confessor, who was distinguished for his sanctity and miracles in the time 
 of Emperor Justin.
\switchcolumn*
\selectlanguage{latin}
\end{paracol}
